% English translation, made from our own transcription (original.tex), not from any existing
% translation. Every math expression, symbol, \tag{}, \origpage{N} marker and \ednote is reproduced
% unchanged; only the prose is translated, including the short prose set inside math via a text
% insert (HOUSESTYLE R21). The \ednote notes flagging this printing's misprints are our own notes,
% already in English, and are carried over verbatim.
%
% Period terminology, kept consistent with abel-1828-remarques:
%   fonction entiere = integral function;  developpement = development;
%   premier/second membre = first/second member;  quantite = quantity;  geometres = geometers
\documentclass{article}
\usepackage{readmasters}
\begin{document}

\origpage{101}

\section*{Researches on elliptic functions.}

(By Mr. \emph{N. H. Abel}.)

For a long time the logarithmic functions, and the exponential and circular functions, were the only transcendental functions that attracted the attention of the geometers. It is only in recent times that one has begun to consider some others. Among these one must distinguish the functions called elliptic, as much for their beautiful analytical properties as for their application in the various branches of mathematics. The first idea of these functions was given by the immortal \emph{Euler}, in demonstrating that the separated equation
\[ \frac{\partial x}{\sqrt{(\alpha+\beta x+\gamma x^{2}+\delta x^{3}+\varepsilon x^{4})}}+\frac{\partial y}{\sqrt{(\alpha+\beta y+\gamma y^{2}+\delta y^{3}+\varepsilon y^{4})}}=0 \tag{1.} \]
is algebraically integrable. After \emph{Euler}, \emph{Lagrange} added something to it, in giving his elegant theory of the transformation of the integral $\displaystyle\int\frac{R.\partial x}{\sqrt{[(1-p^{2}x^{2})(1-q^{2}x^{2})]}}$, where $R$ is a rational function of $x$. But the first and, if I am not mistaken, the only one who has gone deeply into the nature of these functions is Mr. \emph{Legendre}, who, first in a memoir on elliptic functions, and then in his excellent exercises of mathematics, has developed a number of elegant properties\ednote{Printed poropriétés instead of propriétés; an apparent misprint.} of these functions, and has shown their application. At the time of the publication of that work, nothing had been added to the theory of Mr. \emph{Legendre}. I believe that further researches on these functions will not be seen here without pleasure.

In general one understands under the denomination of elliptic functions every function contained in the integral
\[ \int\frac{R\,\partial x}{\sqrt{(\alpha+\beta x+\gamma x^{2}+\delta x^{3}+\varepsilon x^{4})}}, \]
where $R$ is a rational function and $\alpha$, $\beta$, $\gamma$, $\delta$, $\varepsilon$ are constant and \emph{real} quantities. Mr. \emph{Legendre} has demonstrated that by suitable substitutions one can always reduce this integral to the form

\origpage{102}

\[ \int\frac{P\,\partial y}{\sqrt{(a+by^{2}+cy^{4})}}, \]
where $P$ is a rational function of $y^{2}$. By suitable reductions, this integral can then be brought to the form
\[ \int\frac{A+By^{2}}{C+Dy^{2}}\cdot\frac{\partial y}{\sqrt{(a+by^{2}+cy^{4})}}, \]
and this to:
\[ \int\frac{A+B\sin^{2}\theta}{C+D\sin^{2}\theta}\cdot\frac{\partial\theta}{\sqrt{(1-c^{2}\sin^{2}\theta)}}, \]
where $c$ is real, and less than unity.

From this it follows that every elliptic function can be reduced to one of the three forms:
\[ \int\frac{\partial\theta}{\sqrt{(1-c^{2}\sin^{2}\theta)}};\ \int\partial\theta\sqrt{(1-c^{2}\sin^{2}\theta)};\ \int\frac{\partial\theta}{(1+n\sin^{2}\theta)\sqrt{(1-c^{2}\sin^{2}\theta)}}, \]
to which Mr. \emph{Legendre} gives the names of elliptic functions of the first, second and third species. It is these three functions that Mr. \emph{Legendre} has considered; above all the first, which has the most remarkable and the simplest properties.

I propose, in this memoir, to consider the inverse function, that is, the function $\Phi\alpha$, determined by the equations
\[ \begin{gathered} \alpha=\int\frac{\partial\theta}{\sqrt{(1-c^{2}\sin^{2}\theta)}}\ \text{and}\\ \sin\theta=\varphi(\alpha)=x. \end{gathered} \]
The last equation gives
\[ \partial\theta.\sqrt{(1-\sin^{2}\theta)}=\partial.\Phi\alpha=\partial x, \]
hence
\[ \alpha=\int_{0}\frac{\partial x}{\sqrt{[(1-x^{2})(1-c^{2}x^{2})]}}. \]
Mr. \emph{Legendre} supposes $c^{2}$ positive, but I have remarked that the formulas become simpler on supposing $c^{2}$ negative, $=-e^{2}$. Likewise, for the sake of greater symmetry, I write $1-c^{2}x^{2}$ in place of $1-x^{2}$. So that the function $\Phi\alpha=x$ will be given by the equation
\[ \alpha=\int_{0}\frac{\partial x}{\sqrt{[(1-c^{2}x^{2})(1+e^{2}x^{2})]}}, \]
or else
\[ \Phi'\alpha=\partial\alpha.\sqrt{[(1-c^{2}\Phi^{2}\alpha)(1+e^{2}\Phi^{2}\alpha)]}.\]
\ednote{Printed $\partial\alpha$ before the radical; the factor is spurious, since the derivative is the square root alone.}
For brevity, I introduce two other functions of $\alpha$, namely:
\[ f\alpha=\sqrt{(1-c^{2}\Phi^{2}\alpha)};\ F\alpha=\sqrt{(1+e^{2}\Phi^{2}\alpha)}. \]

\origpage{103}

Several properties of these functions are deduced immediately from the known properties of the elliptic function of the first species, but others are more hidden. For example, one demonstrates that the equations $\Phi\alpha=0$, $f\alpha=0$, $F\alpha=0$ have an infinite number of roots, which one can find all. One of the most remarkable properties is that one can express $\Phi(m\alpha)$, $f(m\alpha)$, $F(m\alpha)$ ($m$ being an integer) rationally in terms of $\Phi\alpha$, $f\alpha$, $F\alpha$. Also nothing is easier than to find $\Phi(m\alpha)$, $f(m\alpha)$, $F(m\alpha)$, when one knows $\Phi\alpha$, $f\alpha$, $F\alpha$; but the inverse problem, namely to determine $\Phi\alpha$, $f\alpha$, $F\alpha$ in terms of $\Phi(m\alpha)$, $f(m\alpha)$, $F(m\alpha)$, is more difficult, because it depends on an equation of a higher degree (namely of the degree $m^{2}$).

The solution of this equation is the principal object of this memoir. First it will be shown how one\ednote{Printed ont instead of on; an apparent misprint.} can find all the roots, by means of the functions $\Phi,f,F$. Next the algebraic solution of the equation in question will be treated, and one will arrive at this remarkable result, that $\varphi\left(\frac{\alpha}{m}\right)$; $f\left(\frac{\alpha}{m}\right)$; $F\left(\frac{\alpha}{m}\right)$ can be expressed in terms of $\Phi\alpha$, $f\alpha$, $F\alpha$, by means of a function which, with respect to $\alpha$, contains no other irrationalities than radicals. This produces a very general class of equations which are algebraically soluble. It is to be remarked that the expressions of the roots contain constant quantities which, in general, are not expressible by algebraic quantities. These constant quantities depend on an equation of the degree $\mathrm{m}^{2}-1$. It will be shown how, by means of algebraic functions, the solution of it can be reduced to that of an equation of the degree $\mathrm{m}+1$. Several expressions will be given of the functions $\Phi(2n+1)\alpha$; $f(2n+1)\alpha$; $F(2n+1)\alpha$ in terms of $\Phi\alpha$, $f\alpha$, $F\alpha$. From these will then be deduced the values of $\Phi\alpha$, $f\alpha$, $F\alpha$ in terms of $\alpha$. It will be demonstrated that these functions\ednote{Printed fontions instead of fonctions; an apparent misprint.} can be decomposed into an infinite number of factors, and even into an infinity of partial fractions.

\section*{§. I.}

\emph{Fundamental properties of the functions} $\Phi\alpha$; $f\alpha$; $F\alpha$.

\subsection*{1.}

Supposing that $\Phi\alpha=x$ (1.), we shall have by virtue of what precedes:
\[ \alpha=\int_{0}^{x}\frac{\partial x}{\sqrt{[(1-c^{2}x^{2})(1+e^{2}x^{2})]}}. \tag{2.} \]

\origpage{104}
Hence one sees that $\alpha$, considered as a function of $x$, is positive from $x=0$ to $x=\frac{1}{c}$. Making therefore
\[ \frac{\omega}{2}=\int_{0}^{\frac{1}{c}}\frac{\partial x}{\sqrt{[(1-c^{2}x^{2})(1+e^{2}x^{2})]}}, \tag{3.} \]
it is evident that $\Phi\alpha$ and\ednote{Printed et instead of est; an apparent misprint.} positive and goes on increasing from $\alpha=0$ to $\alpha=\frac{\omega}{2}$, and that we shall have
\[ \varphi(0)=0,\ \varphi\left(\frac{\omega}{2}\right)=\frac{1}{c}. \tag{4.} \]
Because $\alpha$ changes sign when one writes $-x$ in place of $x$, the same is true of the function $\Phi\alpha$ with respect to $\alpha$, and consequently we shall have the equation
\[ \varphi(-\alpha)=-\varphi(\alpha). \tag{5.} \]
Putting in (1.) $xi$ in place of $x$ (where $i$, for brevity, represents the imaginary quantity $\sqrt{-1}$) and denoting the value of $\alpha$ by $\beta i$, there will result
\[ xi=\Phi(\beta i)\ \text{and}\ \beta=\int_{0}^{x}\frac{\partial x}{\sqrt{[(1+c^{2}x^{2})(1-e^{2}x^{2})]}}. \tag{6.} \]
$\beta$ is real and positive from $x=0$ to $x=\frac{1}{e}$, hence making
\[ \frac{\varpi}{2}=\int_{0}^{\frac{1}{e}}\frac{\partial x}{\sqrt{[(1-e^{2}x^{2})(1+c^{2}x^{2})]}}, \tag{7.} \]
$x$ will be positive from $\beta=$ zero to $\beta=\frac{\varpi}{2}$, that is, the function $\frac{1}{i}\Phi(\beta i)$ will be positive between the same limits. Making $\beta=\alpha$ and $y=\frac{\varphi(\alpha i)}{i}$, we have
\[ \alpha=\int_{0}^{y}\frac{\partial y}{\sqrt{[(1-e^{2}y^{2})(1+c^{2}y^{2})]}}, \]
hence one sees that, on supposing $c$ in place of $e$ and $e$ in place of \uncertain{$\varepsilon$},\ednote{Printed as a curled e (apparently $\varepsilon$) at the last place; the sense requires $c$ (interchange of $c$ and $e$).}
\[ \frac{\varphi(\alpha i)}{i}\ \text{will be changed into}\ \Phi\alpha. \]

And because
\[ \begin{gathered} f\alpha=\sqrt{(1-c^{2}\Phi^{2}\alpha)},\\ F\alpha=\sqrt{(1+e^{2}\Phi^{2}\alpha)}, \end{gathered} \]
one sees that by the change of $c$ into $e$ and $e$ into $c$, $f(\alpha i)$ and $F(\alpha i)$ will return respectively into $F\alpha$ and $f\alpha$. Finally the equations (2.) and (7.) show that by the same transformation, $\omega$ and $\varpi$ will return respectively into $\varpi$ and $\omega$.

\origpage{105}

According to (7.) we shall have $x=\frac{1}{e}$ for $\beta=\frac{\varpi}{2}$, hence by virtue of the equation $xi=\varphi(\beta i)$, there will result
\[ \varphi\left(\frac{\varpi i}{2}\right)=i.\frac{1}{e}. \tag{8.} \]

\subsection*{2.}

By virtue of what precedes, we shall have the values of $\Phi\alpha$ for every real value of $\alpha$ lying between $-\frac{\omega}{2}$ and $+\frac{\omega}{2}$, and for every imaginary value of this quantity of the form $\beta i$, if $\beta$ is a quantity continuous between the limits $-\frac{\varpi}{2}$ and $+\frac{\varpi}{2}$. It is now a question of finding the value of this function for any value whatever, real or imaginary, of the variable. To attain this, we shall first establish the fundamental properties of the functions $\Phi$, $f$ and $F$.

Because we have
\[ \begin{gathered} f^{2}\alpha=1-c^{2}\Phi^{2}\alpha,\\ F^{2}\alpha=1+e^{2}\Phi^{2}\alpha, \end{gathered} \]
we shall have, on differentiating:
\[ \begin{gathered} f\alpha.f'\alpha=-c^{2}\Phi\alpha.\varphi'\alpha,\\ F\alpha.F'\alpha=e^{2}\Phi\alpha\,\varphi'\alpha, \end{gathered} \]
Now by (2.) we have
\[ \varphi'\alpha=\sqrt{[(1-c^{2}\Phi^{2}\alpha)(1+e^{2}\Phi^{2}\alpha)]}=f\alpha.F\alpha, \]
hence, on substituting this value of $\Phi'\alpha$ in the two preceding equations, we shall find that the functions $\Phi\alpha$, $f\alpha$, $F\alpha$ are connected with one another by the equations
\[ \left\{\begin{aligned} \varphi'\alpha&=f\alpha.F\alpha,\\ f'\alpha&=-c^{2}\Phi\alpha.F\alpha,\\ F'\alpha&=e^{2}\Phi\alpha.f\alpha. \end{aligned}\right. \tag{9.} \]
This being premised, I say that, denoting by $\alpha$ and $\beta$ two indeterminates, we shall have
\[ \left\{\begin{aligned} \Phi(\alpha+\beta)&=\frac{\varphi\alpha.f\beta.F\beta+\varphi\beta.f\alpha.F\alpha}{1+e^{2}c^{2}.\varphi^{2}\alpha.\varphi^{2}\beta},\\ f(\alpha+\beta)&=\frac{f\alpha.f\beta-c^{2}.\varphi\alpha.\varphi\beta.F\alpha.F\beta}{1+e^{2}c^{2}\varphi^{2}\alpha.\varphi^{2}\beta},\\ F(\alpha+\beta)&=\frac{F\alpha.F\beta+e^{2}.\varphi\alpha.\varphi\beta.f\alpha.f\beta}{1+e^{2}c^{2}\varphi^{2}\alpha.\varphi^{2}\beta}. \end{aligned}\right. \tag{10.} \]

These formulas can be deduced at once from the known properties of the elliptic functions${}^{*)}$; but one can also verify them easily in the following manner.

\textbf{*)} \emph{Legendre} Exercices de calcul intégral.

\origpage{106}

Denoting by $r$ the second member of the first of the equations (10.), we shall have, on differentiating with respect to $\alpha$:
\[ \begin{aligned} \left(\frac{\partial r}{\partial\alpha}\right)&=\frac{\{\varphi'\alpha.f\beta.F\beta+\varphi\beta.F\alpha.f'\alpha+\varphi\beta.f\alpha.F'\alpha\}}{1+e^{2}c^{2}.\varphi^{2}\alpha.\varphi^{2}\beta}\\ &\quad-\frac{(\varphi\alpha.f\beta F\beta+\varphi\beta.f\alpha\ F\alpha).2e^{2}c^{2}\varphi\alpha.\varphi^{2}\beta.\varphi'\alpha}{(1+e^{2}c^{2}\varphi^{2}\alpha.\varphi^{2}\beta)^{2}}. \end{aligned} \]
On substituting for $\Phi'\alpha$, $f'\alpha$, $F'\alpha$ their values given by the equations (9.), there will result
\[ \begin{aligned} \left(\frac{\partial r}{\partial\alpha}\right)&=\frac{f\alpha.F\alpha.f\beta.F\beta}{1+e^{2}c^{2}\varphi^{2}\alpha.\varphi^{2}\beta}-\frac{2e^{2}c^{2}\varphi^{2}\alpha.\varphi^{2}\beta.f\alpha.f\beta.F\alpha.F\beta}{(1+e^{2}c^{2}\varphi^{2}\alpha.\varphi^{2}\beta)^{2}}\\ &\quad+\frac{\varphi\alpha.\varphi\beta.\{1+e^{2}c^{2}\varphi^{2}\alpha\varphi^{2}\beta\}\{-c^{2}F^{2}\alpha+e^{2}f^{2}\alpha\}-2e^{2}c^{2}\varphi\alpha\,\varphi\beta.\varphi^{2}\beta.f^{2}\alpha.F^{2}\alpha}{(1+e^{2}c^{2}\varphi^{2}\alpha\varphi^{2}\beta)^{2}}, \end{aligned} \]
whence, on substituting for $f^{2}\alpha$ and $F^{2}\alpha$ their values: $1-c^{2}\Phi^{2}\alpha$, $1+e^{2}\Phi^{2}\alpha$, and on reducing, we derive
\[ \left(\frac{\partial r}{\partial\alpha}\right)=\frac{(1-e^{2}c^{2}\varphi^{2}\alpha.\varphi^{2}\beta)\{(e^{2}-c^{2})+\varphi\alpha.\varphi\beta+f\alpha.f\beta.F\alpha.F\beta\}-2e^{2}c^{2}\varphi\alpha\,\varphi\beta(\varphi^{2}\alpha+\varphi^{2}\beta)}{(1+e^{2}c^{2}\varphi^{2}\alpha.\varphi^{2}\beta)^{2}}. \]
\ednote{Printed with a plus sign (and a stray hairline before it) between $(e^{2}-c^{2})$ and $\varphi\alpha.\varphi\beta$; a product $(e^{2}-c^{2})\varphi\alpha.\varphi\beta$ is expected.}
Now $\alpha$ and $\beta$ enter symmetrically into the expression of $r$; hence we shall have the value of $\left(\frac{\partial r}{\partial\beta}\right)$ by permuting $\alpha$ and $\beta$ in the value of $\left(\frac{\partial r}{\partial\alpha}\right)$. Now by this the expression of $\left(\frac{\partial r}{\partial\alpha}\right)$ does not change its value, hence we shall have
\[ \left(\frac{\partial r}{\partial\alpha}\right)=\left(\frac{\partial r}{\partial\beta}\right). \]

This equation in partial differentials shows that $r$ is a function of $\alpha+\beta$; hence we shall have
\[ r=\psi(\alpha+\beta). \]
The form of the function $\psi$ will be found by giving to $\beta$ a particular value. Supposing for example $\beta=0$, and remarking that $\Phi(0)=0$, $f(0)=1$, $F(0)=1$, the two values of $r$ will become
\[ r=\Phi(\alpha)\ \text{and}\ r=\psi(\alpha), \]
hence
\[ \psi(\alpha)=\Phi(\alpha), \]
and from this
\[ r=\psi(\alpha+\beta)=\Phi(\alpha+\beta). \]
The first of the formulas (10.) therefore holds good indeed.

In the same manner one will verify the two other formulas.

\origpage{107}

\subsection*{3.}

From the formulas (10.) one can deduce a multitude of others. I shall set down some of the most remarkable: For brevity I make
\[ 1+e^{2}c^{2}\Phi^{2}\alpha.\Phi^{2}\beta=R. \tag{11.} \]

On changing first the sign of $\beta$, we shall obtain
\[ \left\{\begin{aligned} \Phi(\alpha+\beta)+\Phi(\alpha-\beta)&=\frac{2\varphi\alpha.f\beta.F\beta}{R},\\ \Phi(\alpha+\beta)-\Phi(\alpha-\beta)&=\frac{2\varphi\beta.f\alpha.F\alpha}{R},\\ f(\alpha+\beta)+f(\alpha-\beta)&=\frac{2f\alpha.f\beta}{R},\\ f(\alpha+\beta)-f(\alpha-\beta)&=\frac{-2c^{2}.\varphi\alpha.\varphi\beta.F\alpha.F\beta}{R},\\ F(\alpha+\beta)+F(\alpha-\beta)&=\frac{2F\alpha.F\beta}{R},\\ F(\alpha+\beta)-F(\alpha-\beta)&=\frac{2e^{2}.\varphi\alpha.\varphi\beta.f\alpha.f\beta}{R}. \end{aligned}\right. \tag{12.} \]

On forming the product of $\Phi(\alpha+\beta)$ and $\Phi(\alpha-\beta)$, we shall find
\[ \begin{aligned} \Phi(\alpha+\beta).\Phi(\alpha-\beta)&=\frac{\varphi\alpha.f\beta.F\beta+\varphi\beta.f\alpha.F\alpha}{R}.\frac{\varphi\alpha.f\beta.F\beta-\varphi\beta.f\alpha.F\alpha}{R}\\ &=\frac{\varphi^{2}\alpha.f^{2}\beta.F^{2}\beta-\varphi^{2}\beta.f^{2}\alpha.F^{2}\alpha}{R^{2}}, \end{aligned} \]
or, on substituting the values of $f^{2}\beta$, $F^{2}\beta$, $f^{2}\alpha$, $F^{2}\alpha$ in $\Phi\beta$ and $\Phi\alpha$:
\[ \begin{aligned} \Phi(\alpha+\beta).\Phi(\alpha-\beta)&=\frac{\varphi^{2}\alpha-\varphi^{2}\beta-e^{2}c^{2}\varphi^{2}\alpha.\varphi^{4}\beta+e^{2}c^{2}\varphi^{2}\beta.\varphi^{4}\alpha}{R^{2}}\\ &=\frac{(\varphi^{2}\alpha-\varphi^{2}\beta)(1+e^{2}c^{2}\varphi^{2}\alpha.\varphi^{2}\beta)}{R^{2}}; \end{aligned} \]
now $R=1+e^{2}c^{2}\Phi^{2}\alpha.\Phi^{2}\beta$, hence
\[ \Phi(\alpha+\beta).\Phi(\alpha-\beta)=\frac{\varphi^{2}\alpha-\varphi^{2}\beta}{R}. \tag{13.} \]
We shall find in the same way
\[ \left\{\begin{aligned} f(\alpha+\beta).f(\alpha-\beta)&=\frac{f^{2}\alpha-c^{2}\varphi^{2}\beta.F^{2}\alpha}{R}=\frac{f^{2}\beta-c^{2}\varphi^{2}\alpha.F^{2}\beta}{R}\\ &=\frac{1-c^{2}\varphi^{2}\alpha-c^{2}\varphi^{2}\beta-c^{2}e^{2}\varphi^{2}\alpha.\varphi^{2}\beta}{R}=\frac{f^{2}\alpha.f^{2}\beta-c^{2}(c^{2}+e^{2})\varphi^{2}\alpha.\varphi^{2}\beta}{R},\\ F(\alpha+\beta).F(\alpha-\beta)&=\frac{F^{2}\alpha+e^{2}\varphi^{2}\beta.f^{2}\alpha}{R}=\frac{F^{2}\beta+e^{2}\varphi^{2}\alpha.f^{2}\beta}{R}\\ &=\frac{1+e^{2}\varphi^{2}\alpha+e^{2}\varphi^{2}\beta-e^{2}c^{2}\varphi^{2}\alpha.\varphi^{2}\beta}{R}=\frac{F^{2}\alpha.F^{2}\beta-e^{2}(c^{2}+e^{2})\varphi^{2}\alpha.\varphi^{2}\beta}{R}. \end{aligned}\right. \tag{14.} \]

\origpage{108}

\subsection*{4.}

Making in the formulas (10.) $\beta=\pm\frac{\omega}{2}$, $\beta=\pm\frac{\varpi}{2}i$, and remarking that $f\left(\pm\frac{\omega}{2}\right)=0$, $F\left(\pm\frac{\varpi}{2}i\right)=0$, we shall have
\[ \left\{\begin{aligned} &\varphi\left(\alpha\pm\frac{\omega}{2}\right)=\pm\varphi\frac{\omega}{2}.\frac{f\alpha}{F\alpha};\ f\left(\alpha\pm\frac{\omega}{2}\right)=\mp\frac{F\frac{\omega}{2}}{\varphi\frac{\omega}{2}}.\frac{\varphi\alpha}{F\alpha};\\ &F\left(\alpha\pm\frac{\omega}{2}\right)=\frac{F\frac{\omega}{2}}{F\alpha};\\ &\varphi\left(\alpha\pm\frac{\varpi}{2}i\right)=\pm\varphi\frac{\varpi}{2}i.\frac{F\alpha}{f\alpha};\ F\left(\alpha\pm\frac{\varpi}{2}i\right)=\mp\frac{f\frac{\varpi}{2}i}{\varphi\frac{\varpi}{2}i}.\frac{\varphi\alpha}{f\alpha};\\ &f\left(\alpha\pm\frac{\varpi}{2}i\right)=\frac{f\frac{\varpi}{2}i}{f\alpha}; \end{aligned}\right. \tag{15.} \]
or else:
\[ \left\{\begin{aligned} &\varphi\left(\alpha\pm\frac{\omega}{2}\right)=\pm\frac{1}{c}.\frac{f\alpha}{F\alpha};\ f\left(\alpha\pm\frac{\omega}{2}\right)=\mp\sqrt{(e^{2}+c^{2})}.\frac{\varphi\alpha}{F\alpha};\\ &F\left(\alpha\pm\frac{\omega}{2}\right)=\frac{\sqrt{(e^{2}+c^{2})}}{c}.\frac{1}{F\alpha};\\ &\varphi\left(\alpha\pm\frac{\varpi}{2}i\right)=\pm\frac{i}{e}.\frac{F\alpha}{f\alpha};\ F\left(\alpha\pm\frac{\varpi}{2}i\right)=\pm i\sqrt{(e^{2}+c^{2})}.\frac{\varphi\alpha}{f\alpha};\\ &f\left(\alpha\pm\frac{\varpi}{2}i\right)=\frac{\sqrt{(e^{2}+c^{2})}}{e}.\frac{1}{f\alpha}. \end{aligned}\right. \tag{16.} \]

From this we derive at once:
\[ \left\{\begin{aligned} &\varphi\left(\frac{\omega}{2}+\alpha\right)=\Phi\left(\frac{\omega}{2}-\alpha\right);\ f\left(\frac{\omega}{2}+\alpha\right)=-f\left(\frac{\omega}{2}-\alpha\right);\\ &F\left(\frac{\omega}{2}+\alpha\right)=F\left(\frac{\omega}{2}-\alpha\right);\\ &\varphi\left(\frac{\varpi}{2}i+\alpha\right)=\Phi\left(\frac{\varpi}{2}i-\alpha\right);\ F\left(\frac{\varpi}{2}i+\alpha\right)=-F\left(\frac{\varpi}{2}i-\alpha\right);\\ &f\left(\frac{\varpi}{2}i+\alpha\right)=f\left(\frac{\varpi}{2}i-\alpha\right); \end{aligned}\right. \tag{17.} \]
\[ \varphi\left(\alpha\pm\frac{\omega}{2}\right).\varphi\left(\alpha+\frac{\varpi}{2}i\right)=\pm\frac{i}{ce};\ F\left(\alpha\pm\frac{\omega}{2}\right).F\alpha=\frac{b}{c};\ f\left(\alpha\pm\frac{\varpi}{2}i\right).f\alpha=\frac{b}{e}. \tag{18.} \]
Making $\alpha=\frac{\omega}{2}$ and $\frac{\varpi}{2}i$, we find from this
\[ \Phi\left(\frac{\omega}{2}+\frac{\varpi}{2}i\right)=\frac{1}{0},\ f\left(\frac{\omega}{2}+\frac{\varpi}{2}i\right)=\frac{1}{0},\ F\left(\frac{\omega}{2}+\frac{\varpi}{2}i\right)=\frac{1}{0}. \]

\origpage{109}
Next the equations (17.), on putting in the first three $\alpha+\frac{\omega}{2}$ in place of $\alpha$, and in the last three $\alpha+\frac{\varpi}{2}i$ in place of $\alpha$, give the following
\[ \left\{\begin{aligned} &\varphi(\alpha+\omega)=-\Phi\alpha;\ f(\alpha+\omega)=-f\alpha;\ F(\alpha+\omega)=F\alpha;\\ &\varphi(\alpha+\varpi i)=-\Phi\alpha;\ f(\alpha+\varpi i)=f\alpha;\ F(\alpha+\varpi i)=-F\alpha; \end{aligned}\right. \tag{19.} \]
and on putting $\alpha+\omega$ and $\alpha+\varpi i$ in place of $\alpha$:
\[ \left\{\begin{aligned} &\Phi(2\omega+\alpha)=\Phi\alpha;\ \Phi(2\varpi i+\alpha)=\Phi\alpha;\ \Phi(\omega+\varpi i+\alpha)=\Phi\alpha;\\ &f(2\omega+\alpha)=f\alpha;\ f(\varpi i+\alpha)=f\alpha;\\ &F(\omega+\alpha)=F\alpha;\ F(2\varpi i+\alpha)=F\alpha. \end{aligned}\right. \tag{20.} \]

These equations show that the functions $\Phi\alpha$, $f\alpha$, $F\alpha$ are \emph{periodic} functions. One will deduce from them without difficulty the following, where $m$ and $n$ are two integers positive or negative:
\[ \left\{\begin{aligned} &\varphi((m+n)\omega+(m-n)\varpi i+\alpha)=\varphi\alpha;\ \varphi((m+n)\omega+(m-n+1)\varpi i+\alpha)=-\varphi\alpha;\\ &f(2m\omega+n\varpi i+\alpha)=f\alpha;\ f((2m+1)\omega+n\varpi i+\alpha)=-f\alpha;\\ &F(m\omega+2n\varpi i+\alpha)=+F\alpha;\ F(m\omega+(2n+1)\varpi i+\alpha)=-F\alpha. \end{aligned}\right. \tag{21.} \]
These formulas can also be written as follows:
\[ \left\{\begin{aligned} &\Phi(m\omega+n\varpi i\pm\alpha)=\pm(-1)^{m+n}.\Phi\alpha,\\ &f(m\omega+n\varpi i\pm\alpha)=(-1)^{m}.f\alpha,\\ &F(m\omega+n\varpi i\pm\alpha)=(-)^{n}.F\alpha. \end{aligned}\right. \tag{22.} \]
\ednote{In the last line of this system the base of the exponent is printed as a bare minus sign, without the digit 1; an apparent misprint for $(-1)^{n}$.}
One can remark as particular cases:
\[ \left\{\begin{aligned} &\Phi(m\omega\pm\alpha)=\pm(-1)^{m}.\Phi\alpha;\ \Phi(n\varpi i\pm\alpha)=\pm(-1)^{n}.\Phi\alpha;\\ &f(m\omega\pm\alpha)=(-1)^{m}.f\alpha;\ f(n\varpi i\pm\alpha)=f\alpha;\\ &F(m\omega\pm\alpha)=F\alpha;\ F(n\varpi i\pm\alpha)=(-1)^{n}.F\alpha. \end{aligned}\right. \tag{22$'$.} \]

\subsection*{5.}

The formulas which we have just established show that we shall have the values of the functions $\Phi\alpha$; $f\alpha$; $F\alpha$ for all the real or imaginary values of the variable, when we know them for the real values of this quantity lying between $\frac{\omega}{2}$ and $-\frac{\omega}{2}$ and for the imaginary values of the form $\beta i$, where $\beta$ lies between $\frac{\varpi}{2}$ and $-\frac{\varpi}{2}$.

In effect, suppose that the value of the functions $\Phi(\alpha+\beta i)$, $f(\alpha+\beta i)$, $F(\alpha+\beta i)$ is asked, where $\alpha$ and $\beta$ are any real quantities whatever. On putting in the formulas (10.) $\beta i$ in place of $\beta$, it is clear that we shall have the three functions in question expressed by the functions: $\Phi\alpha$; $f\alpha$; $F\alpha$; $\Phi(\beta i)$; $f(\beta i)$; $F(\beta i)$. It remains therefore only to determine these latter.

\origpage{110}
Now, whatever the values of $\alpha$ and $\beta$ may be, one can always find two integral numbers $m$ and $n$, such that $\alpha=m.\omega\pm\alpha'$, $\beta=n\varpi\pm\beta'$, where $\alpha'$ is a quantity comprised between $0$ and $+\frac{\omega}{2}$, and $\beta'$ between $0$ and $+\frac{\varpi}{2}$. Therefore, by virtue of equations (22.), on substituting the preceding values of $\alpha$ and $\beta$, there will result:
\[ \begin{aligned} \Phi(\alpha)&=\varphi(m\omega\pm\alpha')=\pm(-1)^{m}.\Phi\alpha',\\ f(\alpha)&=f(m\omega\pm\alpha')=(-1)^{m}.f\alpha',\\ F(\alpha)&=F(m\omega\pm\alpha')=F\alpha',\\ \Phi(\beta i)&=\varphi(+n\varpi i\pm\beta'i)=\pm(-1)^{n}.\varphi(\beta'i),\\ f(\beta i)&=f(n\varpi i\pm\beta'i)=f(\beta'i),\\ F(\beta i)&=F(n\varpi i\pm\beta'i)=(-1)^{n}.F(\beta'i). \end{aligned} \]
Therefore the functions $\Phi\alpha$, $f\alpha$, $F\alpha$, $\Phi(\beta i)$, $f(\beta i)$, $F(\beta i)$ will be expressed as has just been said, and consequently also the functions $\Phi(\alpha+\beta i)$; $f(\alpha+\beta i)$; $F(\alpha+\beta i)$.

We have seen previously that $\Phi\alpha$ is real from $\alpha=-\frac{\omega}{2}$ to $\alpha=+\frac{\omega}{2}$, and that $\frac{\varphi(\alpha i)}{i}$ is real from $\alpha=-\frac{\varpi}{2}$ to $\alpha=+\frac{\varpi}{2}$.

Therefore, by virtue of equations (22.), it is clear:

1) that $\Phi(\alpha)$ and $\frac{\varphi(\alpha i)}{i}$ are real for every real value of $\alpha$; $\Phi\alpha$ is comprised between $-\frac{1}{c}$ and $+\frac{1}{c}$, and $\frac{\varphi(\alpha i)}{i}$ between $-\frac{1}{e}$ and $+\frac{1}{e}$;

2) that $\Phi(\alpha)$ vanishes for $\alpha=m\omega$, and $\frac{\varphi(\alpha i)}{i}$ for $\alpha=m\varpi$, $m$ being a positive or negative integral number; but $\Phi(\alpha)$ is not zero for any other real value of $\alpha$.

On remarking that $f\alpha=\sqrt{(1-c^{2}\Phi^{2}\alpha)}$, $F\alpha=\sqrt{(1+e^{2}\Phi^{2}\alpha)}$, it follows from what we have just said:

1) that the functions $f(\alpha)$; $F(\alpha)$; $f(\alpha i)$; $F(\alpha i)$ are real for every value of $\alpha$;

2) that $f(\alpha)$ is comprised between the limits $-1$ and $+1$ and $F(\alpha)$ between the limits $+1$ and $+\sqrt{\left(1+\frac{e^{2}}{c^{2}}\right)}$, so that $F\alpha$ is positive for every real value of $\alpha$;

3) that $f(\alpha i)$ is positive and comprised between the limits $+1$ and $\sqrt{\left(1+\frac{c^{2}}{e^{2}}\right)}$ and $F(\alpha i)$ between the limits $-1$ and $+1$ for every real value of $\alpha$;

\origpage{111}

4) that $f(\alpha)$ vanishes for $\alpha=(m+\frac{1}{2})\omega$ and $F(i\alpha)$ for $\alpha=(m+\frac{1}{2})\varpi$; but for no other value of $\alpha$.

The following will be remarked as corollaries drawn from the formulas (22.):

1) Let $\alpha=0$. In this case, on remarking that $\Phi(0)=0$, $f(0)=1$, $F(0)=1$, we shall have
\[ \left\{\begin{aligned} \varphi(m\omega+n\varpi i)&=0,\\ f(m\omega+n\varpi i)&=(-1)^{m},\\ F(m\omega+n\varpi i)&=-1. \end{aligned}\right. \tag{23.} \]
\ednote{The last line is printed as $-1$ over a blotted impression; by formula (22.) one would expect $(-1)^{n}$.}

2) Let $\alpha=\frac{\omega}{2}$. By virtue of the equations:
\[ \Phi\left(\frac{\omega}{2}\right)=\frac{1}{c};\ f\left(\frac{\omega}{2}\right)=0;\ F\left(\frac{\omega}{2}\right)=\frac{b}{c}; \]
we shall have
\[ \left\{\begin{aligned} \Phi((m+\tfrac{1}{2})\omega+n\varpi i)&=(-1)^{m+n}.\frac{1}{c},\\ f((m+\tfrac{1}{2})\omega+n\varpi i)&=0,\\ F((m+\tfrac{1}{2})\omega+n\varpi i)&=(-1)^{n}.\frac{b}{c}. \end{aligned}\right. \tag{24.} \]

3) Let $\alpha=\frac{\varpi}{2}i$. By virtue of the equations
\[ \Phi\left(\frac{\varpi}{2}i\right)=\frac{i}{e};\ f\left(\frac{\varpi}{2}i\right)=\frac{b}{e};\ F\left(\frac{\varpi}{2}i\right)=0; \]
we shall have:
\[ \left\{\begin{aligned} \Phi(m\omega+(n+\tfrac{1}{2})\varpi i)&=(-1)^{m+n}\frac{i}{e},\\ f(m\omega+(n+\tfrac{1}{2})\varpi i)&=(-1)^{m}.\frac{b}{e},\\ F(m\omega+(n+\tfrac{1}{2})\varpi i)&=0. \end{aligned}\right. \tag{25.} \]

4) Let $\alpha=\frac{\omega}{2}+\frac{\varpi}{2}i$. By virtue of the equations above we shall have
\[ \left\{\begin{aligned} \Phi((m+\tfrac{1}{2})\omega+(n+\tfrac{1}{2})\varpi i)&=\frac{1}{0},\\ f((m+\tfrac{1}{2})\omega+(n+\tfrac{1}{2})\varpi i)&=\frac{1}{0},\\ F((m+\tfrac{1}{2})\omega+(n+\tfrac{1}{2})\varpi i)&=\frac{1}{0}. \end{aligned}\right. \tag{26.} \]

\subsection*{6.}

The equations (23.), (24.), (25.) show that the function $\Phi(\alpha)$ vanishes whenever $\alpha$ is of the form $\alpha=m\omega+n\varpi i$; that $f\alpha$ vanishes whenever $\alpha$ is of the form $\alpha=(m+\frac{1}{2})\omega+n\varpi i$, and that $F\alpha$ vanishes whenever $\alpha$ is of the form $\alpha=m\omega+(n+\frac{1}{2})\varpi i$.

\origpage{112}
Now I say that for every other value of $\alpha$, the functions $\Phi\alpha$, $f\alpha$, $F\alpha$ will necessarily have a value different from zero.

Suppose in fact that we have
\[ \Phi(\alpha+\beta i)=0, \]
$\alpha$ and $\beta$ being real quantities. By virtue of the first of the formulas (10.), this equation may be written as follows:
\[ \frac{\varphi\alpha.f(\beta i).F(\beta i)+\varphi(\beta i).f\alpha.F\alpha}{1+e^{2}c^{2}\varphi^{2}\alpha.\varphi^{2}(\beta i)}=0. \]

Now the quantities $\Phi\alpha$, $f(\beta i)$, $F(\beta i)$ are real and $\Phi(\beta i)$ is of the form $i.A$, where $A$ is real; therefore this equation cannot subsist unless we have separately:
\[ \Phi(\alpha).f(\beta i).F(\beta i)=0;\ \Phi(\beta i).f\alpha.F\alpha=0. \]
These equations can be satisfied only in two ways, namely by making
\[ \Phi(\alpha)=0,\ \Phi(\beta i)=0, \]
or
\[ f(\beta i).F(\beta i)=0,\ f\alpha.F\alpha=0. \]

The first two equations give $\alpha=m\omega$; $\beta=n\varpi$.

The last two, on remarking that $F\alpha$ and $f(\beta i)$ can never vanish, give
\[ f\alpha=0,\ F(\beta i)=0, \]
whence
\[ \alpha=(m+\tfrac{1}{2})\omega,\ \beta=(n+\tfrac{1}{2})\varpi. \]
But for these values of $\alpha$ and $\beta$ the value of $\Phi(\alpha+\beta i)$ will become infinite; therefore the only values of $\alpha$ and $\beta$ are $\alpha=m\omega$ and $\beta=n\varpi$, and consequently all the roots of the equation
\[ \Phi(x)=0, \]
may be represented by
\[ x=m\omega+n\varpi i. \tag{27.} \]
In the same manner we shall find that all the roots of the equation
\[ f(x)=0, \]
may be represented by
\[ x=(m+\tfrac{1}{2})\omega+n\varpi i, \tag{28.} \]
and those of the equation
\[ F(x)=0, \]
by
\[ x=m\omega+(n+\tfrac{1}{2})\varpi i. \tag{29.} \]

\origpage{113}

\subsection*{7.}

The formulas (26.) show that the three equations
\[ \Phi(x)=\frac{1}{0},\ f(x)=\frac{1}{0},\ F(x)=\frac{1}{0}, \]
are satisfied by giving to $x$ one of the values of the form
\[ x=(m+\tfrac{1}{2})\omega+(n+\tfrac{1}{2})\varpi i. \tag{30.} \]

Now it can be demonstrated that the equations in question have no other roots.

In fact, having
\[ \begin{gathered} F(x)=\frac{b}{c}.\frac{1}{F\left(x-\frac{\omega}{2}\right)};\ \Phi(x)=\frac{-i}{ec}.\frac{1}{\varphi\left(x-\frac{\omega}{2}-\frac{\varpi}{2}i\right)};\\ f(x)=\frac{b}{e}.\frac{1}{f\left(x-\frac{\varpi}{2}i\right)}; \end{gathered} \]
the equations in question will entail these:
\[ \varphi\left(x-\frac{\omega}{2}-\frac{\varpi}{2}i\right)=0;\ f\left(x-\frac{\varpi}{2}i\right)=0;\ F\left(x-\frac{\omega}{2}\right)=0, \]
but by virtue of what we have just seen in the preceding number, these equations give respectively:
\[ \begin{gathered} x-\frac{\omega}{2}-\frac{\varpi}{2}i=m\omega+n\varpi i;\ x-\frac{\varpi}{2}i=(m+\tfrac{1}{2})\omega+n\varpi i;\\ x-\frac{\omega}{2}=m\omega+(n+\tfrac{1}{2})\varpi i; \end{gathered} \]
that is: we shall have for the three equations:
\[ x=(m+\tfrac{1}{2})\omega+(n+\tfrac{1}{2})\varpi i, \]
Q.\ E.\ D.

\subsection*{8.}

Having found as above all the roots of the equations
\[ \begin{gathered} \Phi(x)=0;\ f(x)=0;\ F(x)=0;\\ \Phi(x)=\frac{1}{0};\ f(x)=\frac{1}{0};\ F(x)=\frac{1}{0}; \end{gathered} \]
I shall now seek the roots of the more general equations.
\[ \Phi(x)=\Phi a,\ f(x)=fa,\ F(x)=Fa, \]
where $a$ is any quantity whatever, real or imaginary.

Let us consider first the equation
\[ \Phi(x)-\Phi a=0. \]
On making in the second of the formulas (12.)
\[ \alpha=\frac{x+a}{2},\ \beta=\frac{x-a}{2}, \]
we shall find
\[ \Phi x-\Phi a=\frac{2\varphi\left(\frac{x-a}{2}\right).f\left(\frac{x+a}{2}\right).F\left(\frac{x+a}{2}\right)}{1+e^{2}c^{2}.\varphi^{2}\left(\frac{x+a}{2}\right).\varphi^{2}\left(\frac{x-a}{2}\right)}=0. \]

\origpage{114}
This equation can subsist only in one of the five following cases:

1. if $\Phi\left(\frac{x-a}{2}\right)=0$, whence $x=a+2m\omega+2n\varpi i$,

2. if $f\left(\frac{x+a}{2}\right)=0$, whence $x=-a+(2m+1)\omega+2n\varpi i$,

3. if $F\left(\frac{x+a}{2}\right)=0$, whence $x=-a+2m\omega+(2n+1)\varpi i$,

4. if $\Phi\left(\frac{x-a}{2}\right)=\frac{1}{0}$, whence $x=a+(2m+1)\omega+(2n+1)\varpi i$,

5. if $\Phi\left(\frac{x+a}{2}\right)=\frac{1}{0}$, whence $x=-a+(2m+1)\omega+(2n+1)\varpi i$.

The solution of these five equations is contained in the formulas (27.), (28.), (29.).

Of the values found for $x$, those must be rejected which the formula
\[ x=-a+(2m+1)\omega+(2n+1)\varpi i, \]
gives, for such a value of $x$ gives by virtue of (22.):
\[ \Phi x=-\Phi a, \]
whereas we must have $\Phi x=\Phi a$; but the other values of $x$, expressed by the first four formulas, may be admitted. They are, as one sees, contained in the single formula:
\[ x=(-1)^{m+n}.a+m\omega+n\varpi i. \tag{31.} \]
Such then is the general expression of all the roots of the equation
\[ \Phi x=\Phi a. \]
In the same manner we shall find that all the roots of the equation
\[ fx=fa \]
are represented by the formula
\[ x=\pm a+2m\omega+n\varpi i, \tag{32.} \]
and all those of the equation
\[ Fx=Fa, \]
by the formula
\[ x=\pm a+m\omega+2n\varpi i. \tag{33.} \]

\section*{§. II.}

\emph{Formulas which give the values of} $\Phi(n\alpha)$, $f(n\alpha)$, $F(n\alpha)$ \emph{expressed in rational functions of} $\Phi\alpha$, $f\alpha$, $F\alpha$.

\subsection*{9.}

Let us resume the formulas (12.). On making in the $1{}^{\mathrm{e}}$, $3{}^{\mathrm{e}}$ and $5{}^{\mathrm{e}}$ $\alpha=n\beta$, there will result:

\origpage{115}
\[ \left\{\begin{aligned} \Phi(n+1)\beta&=-\Phi(n-1)\beta+\frac{2\varphi(n\beta).f\beta.F\beta}{R},\\ f(n+1)\beta&=-f(n-1)\beta+\frac{2f(n\beta).f\beta}{R},\\ F(n+1)\beta&=-F(n-1)\beta+\frac{2F(n\beta).F\beta}{R}, \end{aligned}\right. \tag{34.} \]
where $R=1+c^{2}e^{2}.\Phi^{2}(n\beta).\Phi^{2}\beta$.

These formulas give the value of $\Phi(n+1)\beta$ in terms of $\Phi(n-1)\beta$ and $\Phi(n\beta)$; that of $f(n+1)\beta$ in terms of $f(n-1)\beta$ and $f(n\beta)$, and that of $F(n+1)\beta$ in terms of $F(n-1)\beta$ and $F(n\beta)$. Therefore, on making successively $n=1$, 2, 3\ldots, we shall find successively the values of the functions:
\[ \begin{gathered} \Phi(2\beta);\ \Phi(3\beta);\ \Phi(4\beta)\ .....\ \Phi(n\beta),\\ f(2\beta);\ f(3\beta);\ f(4\beta)\ .....\ f(n\beta),\\ F(2\beta);\ F(3\beta);\ F(4\beta)\ .....\ F(n\beta) \end{gathered} \]
expressed in rational functions of the three quantities
\[ \Phi\beta;\ f\beta;\ F\beta. \]
On making e.g.\ $n=1$, we shall have:
\[ \left\{\begin{aligned} \Phi(2\beta)&=\frac{2.\varphi\beta.f\beta.F\beta}{1+e^{2}c^{2}\varphi^{4}\beta},\\ f(2\beta)&=-1+\frac{2f^{2}\beta}{1+e^{2}c^{2}\varphi^{4}\beta},\\ F(2\beta)&=-1+\frac{2F^{2}\beta}{1+e^{2}c^{2}\varphi^{4}\beta}. \end{aligned}\right. \tag{35.} \]

The functions $\Phi(n\beta)$, $f(n\beta)$, $F(n\beta)$ being rational functions of $\Phi\beta$, $f\beta$, $F\beta$, they can always be reduced to the form $\frac{P}{Q}$, where $P$ and $Q$ are integral functions of $\Phi\beta$, $f\beta$, $F\beta$. Likewise it is clear that the denominator $Q$ will have the same value for the three functions under consideration. Let then
\[ \Phi(n\beta)=\frac{P_n}{Q_n},\ f(n\beta)=\frac{P'_n}{Q_n},\ F(n\beta)=\frac{P''_n}{Q_n}, \tag{35$'$.} \]
we shall likewise have
\[ \begin{gathered} \Phi(n+1)\beta=\frac{P_{n+1}}{Q_{n+1}},\ f(n+1)\beta=\frac{P'_{n+1}}{Q_{n+1}},\ F(n+1)\beta=\frac{P''_{n+1}}{Q_{n+1}},\\ \Phi(n-1)\beta=\frac{P_{n-1}}{Q_{n-1}},\ f(n-1)\beta=\frac{P'_{n-1}}{Q_{n-1}},\ F(n-1)\beta=\frac{P''_{n-1}}{Q_{n-1}}. \end{gathered} \]
On substituting these values, the first of the formulas (34.) will become:
\[ \frac{P_{n+1}}{Q_{n+1}}=-\frac{P_{n-1}}{Q_{n-1}}+\frac{2f\beta.F\beta.\dfrac{P_n}{Q_n}}{1+c^{2}e^{2}\varphi^{2}\beta.\dfrac{P_n^{2}}{Q_n^{2}}}, \]

\origpage{116}
or else
\[ \frac{P_{n+1}}{Q_{n+1}}=\frac{-P_{n-1}.(Q_n^{2}+c^{2}e^{2}\varphi^{2}\beta.P_n^{2})+2P_n.Q_n.Q_{n-1}.f\beta.F\beta}{Q_{n-1}.(Q_n^{2}+e^{2}c^{2}\varphi^{2}\beta.P_n^{2})}. \]

On equating the numerators and the denominators of these two fractions, we shall have
\[ P_{n+1}=-P_{n-1}(Q_n^{2}+c^{2}e^{2}\Phi^{2}\beta.P_n^{2})+2f\beta.F\beta.P_n.Q_n.Q_{n-1}, \tag{36.} \]
\[ Q_{n+1}=Q_{n-1}(Q_n^{2}+e^{2}c^{2}\Phi^{2}\beta.P_n^{2}). \tag{37.} \]
The second and the third of the equations (34.) will give in the same manner:
\[ P'_{n+1}=-P'_{n-1}.(Q_n^{2}+c^{2}e^{2}\Phi^{2}\beta.P_n^{2})+2f\beta\,P'_n.Q_n.Q_{n-1}, \tag{38.} \]
\[ P''_{n+1}=-P''_{n-1}.(Q_n^{2}+c^{2}e^{2}\Phi^{2}\beta.P_n^{2})+2F\beta.P''_n.Q_n.Q_{n-1}. \tag{39.} \]
On making in these four formulas $n=1$, 2, 3\ ....., and remarking that we shall have:
\[ \begin{gathered} Q_0=1;\ Q_1=1;\ P_0=0;\ P_1=\Phi\beta;\\ P'_0=1;\ P'_1=f\beta;\ P''_0=1;\ P''_1=F\beta; \end{gathered} \]
we shall find successively the integral functions $Q_n$, $P_n$, $P'_n$, $P''_n$, for all values of $n$.

Let, for brevity:
\[ \Phi\beta=x,\ f\beta=y,\ F\beta=z\ \text{and} \tag{40.} \]
\[ R_n=Q_n^{2}+e^{2}c^{2}x^{2}P_n^{2}, \tag{41.} \]
the preceding formulas will give:
\[ \left\{\begin{aligned} Q_{n+1}&=Q_{n-1}.R_n,\\ P_{n+1}&=-P_{n-1}.R_n+2yz.P_n.Q_n.Q_{n-1},\\ P'_{n+1}&=-P'_{n-1}.R_n+2y.P'_n.Q_n.Q_{n-1},\\ P''_{n+1}&=-P''_{n-1}.R_n+2z.P''_n.Q_n.Q_{n-1}. \end{aligned}\right. \tag{42.} \]
On putting $n=1$, 2, we shall have:
\[ \left\{\begin{aligned} R_1&=Q_1^{2}+e^{2}c^{2}x^{2}P_1^{2}=1+e^{2}c^{2}x^{4},\\ Q_2&=Q_0R_1=1+e^{2}c^{2}x^{4},\\ P_2&=-P_0R_1+2yzP_1.Q_1.Q_0=2xyz,\\ P'_2&=-P'_0R_1+2yP'_1.Q_1.Q_0=-1-e^{2}c^{2}x^{4}+2y^{2},\\ P''_2&=-P''_0R_1+2zP''_1.Q_1.Q_0=-1-e^{2}c^{2}x^{4}+2z^{2}. \end{aligned}\right. \tag{43.} \]
\[ \left\{\begin{aligned} R_2&=Q_2^{2}+e^{2}c^{2}x^{2}.P_2^{2}=(1+e^{2}c^{2}x^{4})^{2}+e^{2}c^{2}x^{2}.4x^{2}y^{2}z^{2},\\ Q_3&=Q_1R_2=R_2,\\ P_3&=-P_1R_2+2yzP_2Q_2Q_1=-xR_2+2y^{2}z^{2}x.Q_2\\ &=x\{2y^{2}z^{2}.Q_2-R_2\},\\ P'_3&=-P'_1.R_2+2y.P'_2.Q_2Q_1=-y.R_2+2yP'_2Q_2\\ &=y\{2Q_2P'_2-R_2\},\\ P''_3&=z\{2Q_2P''_2-R_2\}. \end{aligned}\right. \tag{44.} \]

\origpage{117}
On continuing in this manner, and remarking that $y^{2}=1-c^{2}x^{2}$, $z^{2}=1+e^{2}x^{2}$, one will readily see that the quantities:
\[ Q_n,\ \frac{P_{2n}}{xyz},\ \frac{P_{2n+1}}{x},\ P'_{2n},\ \frac{P'_{2n+1}}{y},\ P''_{2n},\ \frac{P''_{2n+1}}{z} \]
are integral functions of the three quantities $x^{2}$, $y^{2}$, $z^{2}$ and consequently also of any one whatever of these quantities, for any integral value whatever of $n$.

This shows that the expressions of $\Phi(n\beta)$, $f(n\beta)$, $F(n\beta)$ will be of the following form:
\[ \left\{\begin{aligned} \Phi(2n\beta)&=\Phi\beta.f\beta.F\beta.T, & \Phi(2n+1)\beta&=\Phi\beta.T',\\ f(2n\beta)&=T_1, & f(2n+1)\beta&=f\beta.T'',\\ F(2n\beta)&=T_2, & F(2n+1)\beta&=F\beta.T''', \end{aligned}\right. \tag{45.} \]
where $T$ etc.\ represent rational functions of the quantities $(\Phi\beta)^{2}$, $(f\beta)^{2}$, $(F\beta)^{2}$.

\section*{§. III.}

\emph{Solution of the equations}
\[ \Phi(n\beta)=\frac{P_n}{Q_n},\ f(n\beta)=\frac{P'_n}{Q_n},\ F(n\beta)=\frac{P''_n}{Q_n}. \]

\subsection*{10.}

According to what has been seen, the functions $\Phi(n\beta)$, $f(n\beta)$, $F(n\beta)$ are expressed rationally in $x$, $y$, $z$. The converse does not hold, for the equations (35.) are in general of a very high degree. They have for this reason a certain number of roots. We shall see how one can readily express all these roots by means of the functions $\Phi$, $f$, $F$.

A. Let us consider first the equation $\Phi(n\beta)=\frac{P_n}{Q_n}$, or $Q_n.\Phi(n\beta)=P_n$, and let us seek all the values of $x$.

Two cases must be distinguished, according as $n$ is even or odd:

1) \emph{If $n$ is an even number.}

From what has been seen in the preceding paragraph (45.), we shall have in this case
\[ \Phi(2n\beta)=\psi(x^{2}).x.y.z, \]
that is, by virtue of the formulas
\[ y=\sqrt{(1-c^{2}x^{2})},\ z=\sqrt{(1+e^{2}x^{2})}, \]
\[ \Phi(2n\beta)=x\psi(x^{2}).\sqrt{[(1-c^{2}x^{2})(1+e^{2}x^{2})]}. \]
Therefore the equation in $x$ will become,
\[ \Phi^{2}(2n\beta)=x^{2}(\psi(x^{2}))^{2}.(1-c^{2}x^{2})(1+e^{2}x^{2}). \]

\origpage{118}
On denoting the second member by $\theta(x^{2})$, we shall have
\[ \Phi^{2}(2n\beta)=\theta(x^{2}). \]
$\Phi\beta$ being one of the values of $x$, we shall have
\[ \Phi^{2}(2n\beta)=\theta(\Phi^{2}\beta), \tag{46.} \]
an equation which holds, whatever the value of $\beta$ may be. To find the other values of $x$, let $x=\Phi\alpha$ be any root whatever; we must have
\[ \Phi^{2}(2n\beta)=\theta(\Phi^{2}\alpha). \]
Now, on putting in (46.) $\alpha$ in place of $\beta$, there will result
\[ \Phi^{2}(2n\alpha)=\theta(\Phi^{2}\alpha),\ \text{therefore:} \]
\[ \Phi^{2}(2n\beta)=\Phi^{2}(2n\alpha); \tag{47.} \]
an equation which comes to these two:
\[ \Phi(2n\alpha)=\Phi(2n\beta)\ \text{and}\ \Phi(2n\alpha)=-\Phi(2n\beta). \]
The first gives by virtue of (31.)
\[ 2n\alpha=2n\beta.(-1)^{m+\mu}+m\omega+\mu\varpi i, \]
where $m$ and $\mu$ are any two integral numbers, positive or negative, zero included.

The second gives the same values of $2m\alpha$\ednote{Printed $2m\alpha$ instead of $2n\alpha$; an apparent misprint.}, but of contrary sign, as it is easy to see, on writing it as follows:
\[ \Phi(-2n\alpha)=\Phi(2n\beta). \]
Every value of $2n\alpha$ which satisfies the equation (47.) can therefore be represented by
\[ 2n\alpha=\pm(2n\beta.(-1)^{m+\mu}+m\omega+\mu\varpi i). \]
From this we draw the value of $\alpha$, on dividing by $2n$, namely:
\[ \alpha=\pm\left((-1)^{m+\mu}.\beta+\frac{m}{2n}\omega+\frac{\mu}{2n}\varpi i\right). \]
Having the value of $\alpha$, we shall have
\[ \Phi\alpha=\pm\Phi\left((-1)^{m+\mu}.\beta+\frac{m}{2n}\omega+\frac{\mu}{2n}\varpi i\right)=x. \tag{48.} \]
Therefore all the values of $x$ are contained in this expression, and they will be found on giving to the numbers $m$ and $\mu$ all the integral values from $-\infty$ to $+\infty$. Now to have all those which are different from one another, it suffices to give to $m$ and $\mu$ integral values less than $2n$. In fact, whatever these numbers may be, one can always suppose them reduced to the form:
\[ m=2n.k+m',\ \mu=2n.k'+\mu', \]
where $k$, $k'$ are integral numbers, and $m'$, $\mu'$, integral numbers less than $2n$. On substituting these values in the expression of $x$, it will become:

\origpage{119}
\[ x=\pm\Phi\left\{(-1)^{m'+\mu'}\beta+\frac{m'}{2n}\omega+\frac{\mu'}{2n}\varpi i+k\omega+k'\varpi i\right\}, \]
now by virtue of (22.) this expression reduces to
\[ x=\pm\Phi\left\{(-1)^{m'+\mu'}\beta+\frac{m'}{2n}\omega+\frac{\mu'}{2n}\varpi i\right\}. \tag{49.} \]
This value of $x$ is of the same form as the preceding (48.), $m$ and $\mu$ only being replaced by $m'$ and $\mu'$, which, both of them, are positive and less than $2n$; therefore we shall obtain all the different values of $x$, on giving only to $m$ and $\mu$ all the integral values from zero up to $2n$ excl. All these values are necessarily different from one another. In fact, suppose for example that we have
\[ \begin{gathered} \pm\Phi\left\{(-1)^{m'+\mu'}\beta+\frac{m'}{2n}\omega+\frac{\mu'}{2n}\varpi i\right\}\\ =\pm\Phi\left\{(-1)^{m+\mu}\beta+\frac{m}{2n}\omega+\frac{\mu}{2n}\varpi i\right\}, \end{gathered} \]
it would follow, according to (31.):
\[ (-1)^{m'+\mu'}\beta+\frac{m'}{2n}\omega+\frac{\mu'}{2n}\varpi i=\pm\left\{(-1)^{m+\mu}.\beta+\frac{m}{2n}\omega+\frac{\mu}{2n}\varpi i\right\}+k\omega+k'\varpi i, \]
$k$ and $k'$ being integers.

This equation gives:
\[ \mu'=k'.2n\pm\mu,\ m'=k.2n\pm m,\ (-1)^{m'+\mu'}=\pm(-1)^{m+\mu}. \]
The first two equations cannot subsist unless $k'=1$, $k=1$, $\mu'=2n-\mu$, $m'=2n-m$, and then the last will become:
\[ (-1)^{4n-m-\mu}=-(-1)^{m+\mu}, \]
whence we draw:
\[ (-1)^{2m+2\mu}=-1, \]
an absurd result.

Therefore all the values of $x$, contained in the formula (48.), are different from one another, if $m$ and $\mu$ are positive and less than $m$\ednote{Printed $m$ instead of $2n$; an apparent misprint.}.

The total number of the values of $x$ is, as it is easy to see, equal to $2(2n)^{2}=8n^{2}$, now the equation $\Phi^{2}(2n\beta)=\theta(x^{2})$ cannot have equal roots, for in that case we should have $\frac{\partial.\theta(x^{2})}{\partial x}=0$, which would give for $x$ a value independent of $\beta$. Therefore the degree of the equation $\Phi^{2}(2n\beta)=\theta(x^{2})$ is equal to the number of the roots, that is to $8n^{2}$. If for example $n=1$, we shall have the equation
\[ \Phi^{2}(2\beta)=\theta(x^{2})=\frac{4x^{2}(1-c^{2}x^{2})(1+e^{2}x^{2})}{(1+e^{2}c^{2}x^{4})^{2}}, \]
or else
\[ (1+e^{2}c^{2}x^{4})^{2}.\Phi^{2}(2\beta)=4x^{2}(1-c^{2}x^{2})(1+e^{2}x^{2}), \]

\origpage{120}
and, according to formula (48.), the roots of this equation, eight in number, will be:
\[ \begin{gathered} x=\pm\Phi\beta,\ x=\pm\Phi\left(-\beta+\frac{\omega}{2}\right),\\ x=\pm\Phi\left(-\beta+\frac{\varpi}{2}i\right),\ x=\pm\Phi\left(\beta+\frac{\omega}{2}+\frac{\varpi}{2}i\right). \end{gathered} \]

2) \emph{If $n$ is an odd number} $=2n+1$.

In this case $\frac{P_{2n+1}}{Q_{2n+1}}$ is, as we have seen, a rational function of $x$, and consequently the equation in $x$ will be:
\[ \Phi(2n+1)\beta=\frac{P_{2n+1}}{Q_{2n+1}}. \tag{50.} \]
Precisely as in the preceding case, one will find that all the roots of this equation may be represented by
\[ x=\Phi\left((-1)^{m+\mu}.\beta+\frac{m}{2n+1}\omega+\frac{\mu}{2n+1}\varpi i\right), \tag{51.} \]
where one must give to $m$ and $\mu$ all integral values from $-n$ to $+n$ inclusive. Hence the number of different roots is $(2n+1)^{2}$. This is also the degree of the equation in question. One may also express the roots by
\[ x=(-1)^{m+\mu}.\Phi\left(\beta+\frac{m}{2n+1}\omega+\frac{\mu}{2n+1}\varpi i\right). \]
If for example $n=1$, one will have an equation of the degree $3^{2}=9$.

The formula (51.) gives for $x$ the following 9 values:
\[ \begin{gathered} \Phi(\beta);\\ \Phi\left(-\beta-\frac{\omega}{3}\right);\\ \Phi\left(-\beta+\frac{\omega}{3}\right);\\ \Phi\left(-\beta-\frac{\varpi}{3}i\right);\\ \Phi\left(-\beta+\frac{\varpi}{3}i\right);\\ \Phi\left(\beta-\frac{\omega}{3}-\frac{\varpi}{3}i\right);\\ \Phi\left(\beta-\frac{\omega}{3}+\frac{\varpi}{3}i\right);\\ \Phi\left(\beta+\frac{\omega}{3}-\frac{\varpi}{3}i\right);\\ \Phi\left(\beta+\frac{\omega}{3}+\frac{\varpi}{3}i\right); \end{gathered} \]

\origpage{121}

B. Let us now consider the equation
\[ f(n\beta)=\frac{P'_n}{Q_n} \tag{52.} \]
and seek the values of $y$ which satisfy this equation. The function $\frac{P'_n}{Q_n}$ being, as has been seen above, rational in $y$: the equation in $y$, on putting $\frac{P'_n}{Q_n}=\psi(y)$, will be
\[ f(n\beta)=\psi(y). \]
One of the roots of this equation is $y=f\beta$, therefore, whatever be the value\ednote{Printed ``leur'' instead of ``valeur''; the first syllable of the word is missing at the line break.} of $\beta$:
\[ f(n\beta)=\psi(f\beta). \tag{53.} \]
To find the other values of $y$, let $\alpha$ be a new unknown, such that $y=f\alpha$; there will be
\[ f(n\beta)=\psi(f\alpha); \]
now, by virtue of (53.), the second member is equal to $f(n\alpha)$, therefore to determine $\alpha$ one will have the equation:
\[ f(n\alpha)=f(n\beta). \]
By virtue of (32.) this equation gives for the general expression of $n\alpha$:
\[ n\alpha=\pm n\beta+2m\omega+\mu\varpi i, \]
$m$ and $\mu$ being two integral numbers, positive or negative, zero included.

From this we derive
\[ \alpha=\pm\beta+\frac{2m}{n}\omega+\frac{\mu}{n}\varpi i, \]
and consequently:
\[ f\alpha=f\left(\pm\beta+\frac{2m}{n}\omega+\frac{\mu}{n}\varpi i\right)=y. \]
This is the general value of $y$. Now, to have the different values of $y$, I say that it suffices to take $\beta$ with the sign $+$ and to give to $m$ and $\mu$ all integral values less than $n$. Indeed, since one has $f(+\alpha)=f(-\alpha)$, one will have first:
\[ f\left(-\beta+\frac{2m}{n}\omega+\frac{\mu}{n}\varpi i\right)=f\left(\beta-\frac{2m}{n}\omega-\frac{\mu}{n}\varpi i\right). \]
Therefore one may always take $\beta$ with the sign $+$ in the expression of $y$. Thus all the values of $y$ are contained in the expression
\[ y=f\left(\beta+\frac{2m}{n}\omega+\frac{\mu}{n}\varpi i\right). \tag{54.} \]

Now, whatever the numbers $m$ and $\mu$ may be, one may always suppose,
\[ m=k.n+m',\ \mu'=k'.n+\mu', \]
\ednote{Printed $\mu'$ on the left of the second equation instead of $\mu$; an apparent misprint.}

\origpage{122}
where $k$, $k'$, $m'$, $\mu'$ are integral numbers, the last two being at the same time positive and less than $n$.

On substituting, there will result
\[ y=f\left(\beta+\frac{2m'}{n}\omega+\frac{\mu'}{n}\varpi i+2k\omega+k'\varpi i\right). \]
Now, by virtue of (22.), the second member of this equation is equal to
\[ f\left(\beta+\frac{2m'}{n}\omega+\frac{\mu'}{n}\varpi i\right)=y, \tag{55.} \]
a quantity of the same form as the second member of (54.); only $m'$ and $\mu'$ are positive and less than $n$. Therefore etc.

On giving to $m$ and $\mu$ all possible values less than $n$, one will find a number $n^{2}$ of values of $y$. Now, in general all these quantities are different from one another. Indeed, suppose for example
\[ f\left(\beta+\frac{2m}{n}\omega+\frac{\mu}{n}\varpi i\right)=f\left(\beta+\frac{2\mu'}{n}+\frac{\mu'}{n}\varpi i\right), \]
\ednote{In the right member the print has $2\mu'/n$ where $2m'\omega/n$ is expected; the $\omega$ is absent. An apparent misprint.}
one will have by virtue of (32.), denoting by $k$, $k'$ two integral numbers:
\[ \beta+\frac{2m}{n}\omega+\frac{\mu}{n}\varpi i=\pm\left(\beta+\frac{2m'}{n}\omega+\frac{\mu'}{n}\varpi i\right)+2k\omega+k'\varpi i. \]
Since $\beta$ may have any irrational value whatever, it is clear that this equation cannot subsist unless one takes in the second member the upper sign. Then there will result
\[ \frac{2m}{n}\omega+\frac{\mu}{n}\varpi i=\frac{2m'}{n}\omega+\frac{\mu'}{n}\varpi i+2k\omega+k'\varpi i, \]
whence one derives, on equating the real parts and the imaginary parts:
\[ m=m'+kn;\ \mu=\mu'+k'n, \]
absurd equations, on observing that the numbers $m$, $m'$, $\mu$ and $\mu'$ are all positive and less than $n$. Therefore in general the equation
\[ f(n\beta)=\psi(y) \]
has a number $n^{2}$ of roots different from one another, and not a greater number. Now in general all the roots of this equation are different from one another. Indeed, if two of them were equal, one would have at once:
\[ f(n\beta)=\psi(y)\ \text{and}\ 0=\psi'(y), \]
and this is impossible, if one observes that the coefficients of $y$ in $\psi(y)$ do not contain $\beta$. Therefore in general the equation (52.) is necessarily of the degree $n^{2}$.

\origpage{123}

C. The equation
\[ F(n\beta)=\frac{P''_n}{Q_n}, \tag{56.} \]
being treated in absolutely the same manner with respect to $z$ as the equation $f(n\beta)=\frac{P'_n}{Q_n}$ has been with respect to $y$, gives for the general expression of the values of $z$:
\[ z=F\left(\beta+\frac{m}{n}\omega+\frac{2\mu}{n}\varpi i\right), \tag{57.} \]
where $m$ and $\mu$ are integral, positive and less than $n$. The number of values of $z$ is $n^{2}$, and they are in general all different from one another.

Therefore in general the equation (56.) is of the degree $n^{2}$.

\subsection*{11.}

We have found above all the roots of the equations
\[ \varphi(n\beta)=\frac{P_n}{Q_n},\ f(n\beta)=\frac{P'_n}{Q_n},\ F(n\beta)=\frac{P''_n}{Q_n}, \]
roots which are expressed by the formulas (48.), (51.), (54.), (57.). All these roots are different from one another, except in the cases of particular values of $\beta$; but for these values the different roots are contained in the same formulas. --- In this latter case a certain number of the values of the quantities $x$, $y$, $z$ will be equal; but it is clear that all the values, equal or unequal, will nevertheless be the roots of the equations in question. This is made evident by making $\beta$ converge towards a particular value which gives for $x$, or $y$, or $z$ equal values.

On putting in the formula (48.) $\beta=\frac{\alpha}{2n}$, one will have the equation
\[ \varphi^{2}\alpha=\frac{P_{2n}^{2}}{Q_{2n}^{2}},\ \text{whose roots are}\ x=\pm\varphi\left((-1)^{m+\mu}\frac{\alpha}{2n}+\frac{m}{2n}\omega+\frac{\mu}{2n}\varpi i\right); \tag{58.} \]
and where $m$ and $\mu$ have all integral and positive values less than $2n$.

On putting likewise in the formula (50.) $\beta=\frac{\alpha}{2n+1}$, one will have
\[ \Phi\alpha=\frac{P_{2n+1}}{Q_{2n+1}},\ \text{whose roots are} \]
\[ x=(-1)^{m+\mu}.\Phi\left(\frac{\alpha}{2n+1}+\frac{m\omega+\mu\varpi i}{2n+1}\right), \tag{59.} \]
$m$ and $\mu$ having for values all integral numbers from $-n$ to $+n$.

Finally, on putting in (52.), (56.) $\beta=\frac{\alpha}{n}$, one will have the equation
\[ f\alpha=\frac{P'_n}{Q_n},\ \text{whose roots are} \]
\[ y=f\left(\frac{\alpha}{n}+\frac{2m}{n}\omega+\frac{\mu}{n}\varpi i\right), \tag{60.} \]

\origpage{124}

and the equation
\[ F\alpha=\frac{P''_n}{Q_n},\ \text{whose roots are} \]
\[ z=F\left(\frac{\alpha}{n}+\frac{m}{n}\omega+\frac{2\mu}{n}\varpi i\right), \tag{61.} \]
where $m$ and $\mu$ are enclosed between the limits 0 and $n-1$ inclusive. If $n$ is odd $=2n+1$, one may also suppose
\[ \begin{gathered} y=(-1)^{m}.f\left(\frac{\alpha}{2n+1}+\frac{m}{2n+1}\omega+\frac{\mu}{2n+1}\varpi i\right),\\ z=(-1)^{\mu}.F\left(\frac{\alpha}{2n+1}+\frac{m}{2n+1}\omega+\frac{\mu}{2n+1}\varpi i\right), \end{gathered} \]
$m$ and $\mu$ having all integral values from $-n$ to $+n$.

In all these equations the quantity $\alpha$ may have any value whatever.

As particular cases one must notice the following:

1) On putting in (58.) and (59.) $\alpha=0$, one will have the equations
\[ \left\{\begin{aligned} &P_{2n}^{2}=0,\ \text{whose roots are}\ x=\pm\Phi\left(\frac{m}{2n}\omega+\frac{\mu}{2n}\varpi i\right)\\ &\text{(the limits of $m$ and $\mu$ being 0 and $2n-1$),}\\ &P_{2n+1}=0,\ \text{whose roots are}\ x=\Phi\left(\frac{m}{2n+1}\omega+\frac{\mu}{2n+1}\varpi i\right)\\ &\text{(the limits of $m$ and $\mu$ being $-n$ and $+n$).} \end{aligned}\right. \tag{62.} \]

2) On putting in (60.) $\alpha=\frac{\omega}{2}$ and in (61.) $\alpha=\frac{\varpi}{2}i$, and observing that $f\left(\frac{\omega}{2}\right)=0$, $F\left(\frac{\varpi}{2}i\right)=0$, one will obtain the two equations:
\[ P'_n=0,\ \text{whose roots are}\ y=f\left((2m+\tfrac{1}{2})\frac{\omega}{n}+\frac{\mu}{n}\varpi i\right) \tag{63.} \]
\[ P''_n=0,\ \text{whose roots are}\ z=F\left(\frac{m}{n}\omega+(2\mu+\tfrac{1}{2})\frac{\varpi i}{n}\right) \tag{64.} \]
(the limits of $m$ and $\mu$ being 0 and $n-1$).

3) On putting in (58.) $\alpha=\frac{\omega}{2}+\frac{\varpi}{2}i$, and observing that $\varphi\left(\frac{\omega}{2}+\frac{\varpi}{2}i\right)=\frac{1}{0}$, one will have the equation
\[ Q_{2n}^{2}=0, \]
whose roots will be:
\[ x=\pm\Phi\left((m+\tfrac{1}{2}(-1)^{m+\mu})\frac{\omega}{2n}+(\mu+\tfrac{1}{2}(-1)^{m+\mu}).\frac{\varpi i}{2n}\right). \]

The values of $x$ must be equal in pairs, and one will readily see that the unequal values may be represented by
\[ x=\Phi\left((m+\tfrac{1}{2})\frac{\omega}{2n}+(\mu+\tfrac{1}{2})\frac{\varpi i}{2n}\right), \tag{65.} \]

\origpage{125}
on giving to $m$ and $\mu$ all integral values from 0 to $2n-1$. Therefore these are the roots of the equation
\[ Q_{2n}=0,\ \text{with respect to } x. \]

On putting likewise in (59.) $\alpha=\frac{\omega}{2}+\frac{\varpi}{2}i$, one will have the equation
\[ Q_{2n+1}=0, \]
whose roots will be:
\[ \left\{\begin{aligned}
x&=(-1)^{m+\mu}.\varphi\left((m+\tfrac{1}{2})\frac{\omega}{2n+1}+(\mu+\tfrac{1}{2})\frac{\varpi i}{2n+1}\right),\\
y&=(-1)^{m}.f\left((m+\tfrac{1}{2})\frac{\omega}{2n+1}+(\mu+\tfrac{1}{2})\frac{\varpi i}{2n+1}\right),\\
z&=(-1)^{\mu}.F\left((m+\tfrac{1}{2})\frac{\omega}{2n+1}+(\mu+\tfrac{1}{2})\frac{\varpi i}{2n+1}\right),
\end{aligned}\right. \tag{66.} \]
$m$ and $\mu$ having for values all integral numbers from $-n$ to $+n$.

Among the values of $x$, $y$, $z$, one must notice that which answers to $m=n$, $\mu=n$. Then one has
\[ \begin{aligned}
x&=(-1)^{n}.\varphi\left(\frac{\omega}{2}+\frac{\varpi}{2}i\right)=\frac{1}{0},\\
y&=(-1)^{n}.f\left(\frac{\omega}{2}+\frac{\varpi}{2}i\right)=\frac{1}{0},\\
z&=(-1)^{n}.F\left(\frac{\omega}{2}+\frac{\varpi}{2}i\right)=\frac{1}{0}.
\end{aligned} \]

These infinite values show that the equation $Q_{2n+1}=0$ is of a degree lower by one unit than that of the equations from which it comes. Setting aside these values, the remaining ones, $(2n+1)^{2}-1$ in number, will be the roots of the equation $Q_{2n+1}=0$.

\section*{§. IV.}

\emph{Algebraic resolution of the equations}
\[ \Phi\alpha=\frac{P_{2n+1}}{Q_{2n+1}},\quad f\alpha=\frac{P'_{2n+1}}{Q_{2n+1}},\quad F\alpha=\frac{P''_{2n+1}}{Q_{2n+1}}. \]

\subsection*{12.}

We have seen in the preceding §. how one may readily express the roots of the equations in question by means of the functions $\Phi$, $f$, $F$. We are now going to deduce from this the resolution of these same equations, or the determination of the functions $\varphi\left(\frac{\alpha}{n}\right)$, $f\left(\frac{\alpha}{n}\right)$, $F\left(\frac{\alpha}{n}\right)$, in terms of the functions $\varphi\alpha$, $f\alpha$, $F\alpha$.

\origpage{126}

As one has
\[ \varphi\left(\frac{\alpha}{m\mu}\right)=\varphi\left(\frac{1}{m}\left(\frac{\alpha}{\mu}\right)\right), \]
one may suppose that $n$ is a prime number. First we shall consider the case where $n=2$, and afterwards that where $n$ is an odd number.

A. \emph{Expressions of the functions} $\varphi\left(\frac{\alpha}{2}\right)$, $f\left(\frac{\alpha}{2}\right)$, $F\left(\frac{\alpha}{2}\right)$.

\subsection*{13.}

The values of $\varphi\left(\frac{\alpha}{2}\right)$, $f\left(\frac{\alpha}{2}\right)$, $F\left(\frac{\alpha}{2}\right)$ may be found very easily in the following manner. Supposing in the formulas (35.) $\beta=\frac{\alpha}{2}$, and putting
\[ x=\varphi\left(\frac{\alpha}{2}\right),\ y=f\left(\frac{\alpha}{2}\right),\ z=F\left(\frac{\alpha}{2}\right), \]
there will result:
\[ f(\alpha)=\frac{y^{2}-c^{2}x^{2}z^{2}}{1+e^{2}c^{2}x^{4}},\ F(\alpha)=\frac{z^{2}+e^{2}y^{2}x^{2}}{1+e^{2}c^{2}x^{4}}, \]
or else, on substituting the values of $y^{2}$ and $z^{2}$ in $x^{2}$:
\[ f(\alpha)=\frac{1-2c^{2}x^{2}-c^{2}e^{2}x^{4}}{1+e^{2}c^{2}x^{4}},\ F(\alpha)=\frac{1+2e^{2}x^{2}-e^{2}c^{2}x^{4}}{1+e^{2}c^{2}x^{4}}. \]
These equations give
\[ \begin{aligned}
1+f\alpha&=\frac{2(1-c^{2}x^{2})}{1+e^{2}c^{2}x^{4}},\ 1-f\alpha=\frac{2c^{2}x^{2}(1+e^{2}x^{2})}{1+e^{2}c^{2}x^{4}},\\
F\alpha-1&=\frac{2e^{2}x^{2}(1-c^{2}x^{2})}{1+e^{2}c^{2}x^{4}},\ F\alpha+1=\frac{2(1+e^{2}x^{2})}{1+e^{2}c^{2}x^{4}},
\end{aligned} \]
\ednote{A blot resembling a 2 is printed after the minus sign in the numerator of the first fraction; the sense requires $2(1-c^{2}x^{2})$, as in the later fractions.}
whence
\[ \frac{F\alpha-1}{1+f\alpha}=e^{2}.x^{2},\ \frac{1-f\alpha}{F\alpha+1}=c^{2}x^{2}, \]
and from this, observing that $y^{2}=1-c^{2}x^{2}$, $z^{2}=1+e^{2}x^{2}$:
\[ z^{2}=\frac{F\alpha+f\alpha}{1+f\alpha},\ y^{2}=\frac{F\alpha+f\alpha}{1+F\alpha}. \]

From these equations one derives, on extracting the square root, and replacing $x$, $y$, $z$ by their values $\varphi\left(\frac{\alpha}{2}\right)$, $f\left(\frac{\alpha}{2}\right)$, $F\left(\frac{\alpha}{2}\right)$:
\[ \left\{\begin{aligned}
\varphi\left(\frac{\alpha}{2}\right)&=\frac{1}{c}\sqrt{\left(\frac{1-f\alpha}{1+F\alpha}\right)}=\frac{1}{e}\sqrt{\left(\frac{F\alpha-1}{f\alpha+1}\right)},\\
f\left(\frac{\alpha}{2}\right)&=\sqrt{\left(\frac{F\alpha+f\alpha}{1+F\alpha}\right)},\ F\left(\frac{\alpha}{2}\right)=\sqrt{\left(\frac{F\alpha+f\alpha}{1+f\alpha}\right)}.
\end{aligned}\right. \tag{67.} \]
Such are the simplest forms that can be given to the values of the functions $\varphi\left(\frac{\alpha}{2}\right)$, $f\left(\frac{\alpha}{2}\right)$, $F\left(\frac{\alpha}{2}\right)$. In this manner one can express

\origpage{127}
algebraically $\varphi\left(\frac{\alpha}{2}\right)$, $f\left(\frac{\alpha}{2}\right)$, $F\left(\frac{\alpha}{2}\right)$ in $f\alpha$, $F\alpha$. In the same manner $\varphi\left(\frac{\alpha}{4}\right)$, $f\left(\frac{\alpha}{4}\right)$, $F\left(\frac{\alpha}{4}\right)$ will be expressed in $f\left(\frac{\alpha}{2}\right)$, $F\left(\frac{\alpha}{2}\right)$, and so on. Therefore in general the functions $\varphi\left(\frac{\alpha}{2^{n}}\right)$, $f\left(\frac{\alpha}{2^{n}}\right)$, $F\left(\frac{\alpha}{2^{n}}\right)$ may be expressed by means of extractions of square roots in terms of the three quantities $\varphi\alpha$, $f\alpha$, $F\alpha$.

To apply the formulas found above for the \emph{bisection} to an example, let us suppose $\alpha=\frac{\omega}{2}$.

Then one will have $f\left(\frac{\omega}{2}\right)=0$, $F\left(\frac{\omega}{2}\right)=\frac{\sqrt{(e^{2}+c^{2})}}{c}$, therefore on substituting:
\[ \varphi\left(\frac{\omega}{4}\right)=\frac{1}{c}\sqrt{\left(\frac{1}{1+\frac{1}{c}\sqrt{(e^{2}+c^{2})}}\right)}=\frac{1}{e}.\sqrt{\left(\frac{1}{c}\sqrt{(e^{2}+c^{2})}-1\right)}, \]
\[ f\left(\frac{\omega}{4}\right)=\sqrt{\left(\frac{\frac{1}{c}\sqrt{(e^{2}+c^{2})}}{1+\frac{1}{c}\sqrt{(e^{2}+c^{2})}}\right)}, \]
\[ F\left(\frac{\omega}{4}\right)=\sqrt{\left(\frac{1}{c}\sqrt{(e^{2}+c^{2})}\right)}, \]
or else
\[ \begin{aligned}
\varphi\left(\frac{\omega}{4}\right)&=\frac{1}{\sqrt{(c^{2}+c\sqrt{(e^{2}+c^{2})})}}=\frac{\sqrt{(\sqrt{(e^{2}+c^{2})}-c)}}{e\sqrt{c}},\\
f\left(\frac{\omega}{4}\right)&=\frac{\sqrt[4]{(e^{2}+c^{2})}}{\sqrt{(c+\sqrt{(e^{2}+c^{2})})}}=\frac{1}{e}.\sqrt{(e^{2}+c^{2}-c\sqrt{(e^{2}+c^{2})})}\\
F\left(\frac{\omega}{4}\right)&=\sqrt[4]{\left(1+\frac{e^{2}}{c^{2}}\right)}=\sqrt{\left[F\left(\frac{\omega}{2}\right)\right]}.
\end{aligned} \]

B. \emph{Expressions of the functions} $\varphi\left(\frac{\alpha}{2n+1}\right)$, $f\left(\frac{\alpha}{2n+1}\right)$, $F\left(\frac{\alpha}{2n+1}\right)$, \emph{in algebraic functions of the quantities} $\varphi\alpha$, $f\alpha$, $F\alpha$.

\subsection*{14.}

To find the values of $\varphi\left(\frac{\alpha}{2n+1}\right)$, $f\left(\frac{\alpha}{2n+1}\right)$, $F\left(\frac{\alpha}{2n+1}\right)$ in $\Phi\alpha$, $f\alpha$, $F\alpha$, one must resolve the equations
\[ \Phi\alpha=\frac{P_{2n+1}}{Q_{2n+1}},\ f\alpha=\frac{P'_{2n+1}}{Q_{2n+1}},\ F\alpha=\frac{P''_{2n+1}}{Q_{2n+1}}, \]
which are all of the degree $(2n+1)^{2}$. We shall see that it is always possible to effect this resolution algebraically.

\origpage{128}

Let
\[ \varphi_{1}\beta=\overset{+n}{\underset{-n}{\Sigma}}_{m}\varphi\left(\beta+\frac{2m\omega}{2n+1}\right) \tag{68.} \]
and
\[ \psi\beta=\overset{+n}{\underset{-n}{\Sigma}}_{\mu}\theta^{\mu}.\varphi_{1}\left(\beta+\frac{2\mu\varpi i}{2n+1}\right),\ \psi_{1}\beta=\overset{+n}{\underset{-n}{\Sigma}}_{\mu}\theta^{\mu}\varphi_{1}\left(\beta-\frac{2\mu\varpi i}{2n+1}\right), \tag{69.} \]
where $\theta$ is any imaginary root whatever of the equation $\theta^{2n+1}-1=0$.
This being premised, I say that the two quantities
\[ \psi\beta.\psi_{1}\beta\ \text{and}\ (\psi\beta)^{2n+1}+(\psi_{1}\beta)^{2n+1}, \]
may be expressed rationally in $\varphi(2n+1)\beta$.

First, on writing $\varphi_{1}\beta$ as follows:
\[ \begin{aligned}
\varphi_{1}\beta&=\varphi\beta.\overset{n}{\underset{1}{\Sigma}}.\left\{\varphi\left(\beta+\frac{2m\omega}{2n+1}\right)+\varphi\left(\beta-\frac{2m\omega}{2n+1}\right)\right\}\\
&=\varphi\beta.\overset{n}{\underset{1}{\Sigma}}(-1)^{m}.\frac{2\varphi\beta.f\left(\frac{2m\omega}{2n+1}\right).F\left(\frac{2m\omega}{2n+1}\right)}{1+e^{2}c^{2}\varphi^{2}\left(\frac{2m\omega}{2n+1}\right).\varphi^{2}\beta},
\end{aligned} \]
\ednote{Printed with a point after $\varphi\beta$ before the sum in both lines, where the sense requires a plus sign; the first line shows a blot over the point and a further point after the sum sign.}
one sees that $\varphi_{1}\beta$ can be expressed rationally in $\varphi\beta$. Let then $\varphi_{1}\beta=\chi(\varphi\beta)$; one has likewise:
\[ \begin{aligned}
\varphi_{1}\left(\beta\pm\frac{2\mu\varpi i}{2n+1}\right)&=\chi\left(\varphi\left(\beta\pm\frac{2\mu\varpi i}{2n+1}\right)\right)\\
&=\chi\left(\frac{\varphi\beta.f\left(\frac{2\mu\varpi i}{2n+1}\right).F\left(\frac{2\mu\varpi i}{2n+1}\right)\pm\varphi\left(\frac{2\mu\varpi i}{2n+1}\right).f\beta.F\beta}{1+e^{2}c^{2}.\varphi^{2}\left(\frac{2\mu\varpi i}{2n+1}\right).\varphi^{2}\beta}\right),
\end{aligned} \]
or else, on putting $\Phi\beta=x$:
\[ f\left(\frac{2\mu\varpi i}{2n+1}\right).F\left(\frac{2\mu\varpi i}{2n+1}\right)=a,\ \varphi\left(\frac{2\mu\varpi i}{2n+1}\right)=b, \]
and on substituting for $f\beta$ and $F\beta$ their values $\sqrt{(1-c^{2}x^{2})}$ and $\sqrt{(1+e^{2}x^{2})}$:
\[ \varphi_{1}\left(\beta\pm\frac{2\mu\varpi i}{2n+1}\right)=\chi\left(\frac{ax\pm b.\sqrt{[(1-c^{2}x^{2})(1+e^{2}x^{2})]}}{1+e^{2}c^{2}b^{2}.x^{2}}\right); \]
now, $\chi$ denoting a rational function, the second member of this equation may be put under the form
\[ R_{\mu}\pm R'_{\mu}.\sqrt{[(1-c^{2}x^{2})(1+e^{2}x^{2})]}, \]
where $R_{\mu}$ and $R'_{\mu}$ are rational functions of $x$.

Therefore one has
\[ \varphi_{1}\left(\beta\pm\frac{2\mu\varpi i}{2n+1}\right)=R_{\mu}\pm R'_{\mu}.\sqrt{[(1-c^{2}x^{2})(1+e^{2}x^{2})]}. \]
On substituting in the expressions of $\psi\beta$ and $\psi_{1}\beta$, there will result:

\origpage{129}
\[ \left\{\begin{aligned}
\psi\beta&=\overset{+n}{\underset{-n}{\Sigma}}_{\mu}(+\theta)^{\mu}.R_{\mu}+\sqrt{[(1-c^{2}x^{2})(1+e^{2}x^{2})]}.\overset{+n}{\underset{-n}{\Sigma}}_{\mu}(+\theta)^{\mu}.R'_{\mu},\\
\psi_{1}\beta&=\overset{+n}{\underset{-n}{\Sigma}}_{\mu}(+\theta)^{\mu}.R_{\mu}-\sqrt{[(1-c^{2}x^{2})(1+e^{2}x^{2})]}.\overset{+n}{\underset{-n}{\Sigma}}_{\mu}(+\theta)^{\mu}.R'_{\mu}.
\end{aligned}\right. \tag{70.} \]
\ednote{The plus sign inside the parentheses before $\theta$ is printed overstruck with a short horizontal bar throughout (70.); kept here as a plain plus.}

Now $R_{\mu}$ and $R'_{\mu}$ being rational functions of $x$, the quantities $\overset{+n}{\underset{-n}{\Sigma}}_{\mu}(+\theta)^{\mu}.R_{\mu}$ and $\overset{+n}{\underset{-n}{\Sigma}}_{\mu}(+\theta)^{\mu}.R'_{\mu}$ are so likewise. On raising then $\psi\beta$ and $\psi_{1}\beta$ to the $(2n+1)^{\text{th}}$ power: the two quantities $(\psi\beta)^{2n+1}$ and $(\psi_{1}\beta)^{2n+1}$ may be put under the form:
\[ \begin{aligned}
(\psi\beta)^{2n+1}&=t+t'.\sqrt{[(1-c^{2}x^{2})(1+e^{2}x^{2})]},\\
(\psi_{1}\beta)^{2n+1}&=t-t'.\sqrt{[(1-c^{2}x^{2})(1+e^{2}x^{2})]},
\end{aligned} \]
$t$ and $t'$ being rational functions of $x$. On taking the sum of the values of $(\psi\beta)^{2n+1}$ and $(\psi_{1}\beta)^{2n+1}$, one will have
\[ (\psi\beta)^{2n+1}+(\psi_{1}\beta)^{2n+1}=2t. \]

Therefore the quantity $(\psi\beta)^{2n+1}+(\psi_{1}\beta)^{2n+1}$ can be expressed rationally in $x$. The same holds of the product $\psi\beta.\psi_{1}\beta$, as one sees by the equations (70.).

Therefore one may put
\[ \left\{\begin{aligned}
\psi\beta.\psi_{1}\beta&=\lambda(x),\\
(\psi\beta)^{2n+1}+(\psi_{1}\beta)^{2n+1}&=\lambda_{1}(x),
\end{aligned}\right. \tag{71.} \]
$\lambda(x)$ and $\lambda_{1}(x)$ denoting rational functions of $x$. Now these functions have the property of not changing in value when one puts in place of $x$ any other root whatever of the equation
\[ \Phi(2n+1)\beta=\frac{P_{2n+1}}{Q_{2n+1}}. \]

Let us consider first the function $\lambda(x)$. On restoring the value of $x=\Phi\beta$, one will have
\[ \psi\beta.\psi_{1}\beta=\lambda(\Phi\beta), \]
whence one derives, on putting $\beta+\frac{2k\omega}{2n+1}+\frac{2k'\varpi i}{2n+1}$ in place of $\beta$:
\[ \lambda\left(\varphi\left(\beta+\frac{2k\omega}{2n+1}+\frac{2k'\varpi i}{2n+1}\right)\right)=\psi\left(\beta+\frac{2k'\varpi i}{2n+1}+\frac{2k\omega}{2n+1}\right).\psi_{1}\left(\beta+\frac{2k'\varpi i}{2n+1}+\frac{2k\omega}{2n+1}\right). \]
This being premised, observing that:
\[ \overset{+n}{\underset{-n}{\Sigma}}\psi(m+k)=\overset{+n}{\underset{-n}{\Sigma}}_{m}\psi(m)+\overset{k}{\underset{1}{\Sigma}}_{m}.(\psi(m+n)-\psi(m-n-1)), \tag{72.} \]
one will have, on putting in the expression of $\psi\beta$, $\beta=\beta+\frac{2k\omega}{2n+1}$:

\origpage{130}
\[ \begin{aligned}
\varphi_{1}\left(\beta+\frac{2k\omega}{2n+1}\right)&=\overset{+n}{\underset{-n}{\Sigma}}_{m}\varphi\left(\beta+\frac{2(k+m)}{2n+1}\omega\right)\\
&=\varphi_{1}\beta+\overset{k}{\underset{1}{\Sigma}}_{m}\left\{\varphi\left(\beta+\frac{2(m+n)\omega}{2n+1}\right)-\varphi\left(\beta+\frac{2(m-n-1)\omega}{2n+1}\right)\right\},
\end{aligned} \]
but: $\varphi\left(\beta+\frac{2(m-n-1)}{2n+1}\omega\right)=\varphi\left(\beta+\frac{2(m+n)}{2n+1}\omega-2\omega\right)=\varphi\left(\beta+\frac{2(m+n)}{2n+1}\omega\right)$,
therefore:
\[ \varphi_{1}\left(\beta+\frac{2k\omega}{2n+1}\right)=\varphi_{1}\beta. \tag{73.} \]

On putting in the expression of $\psi\beta$, $\beta+\frac{2k'\varpi i}{2n+1}+\frac{2k\omega}{2n+1}$ in place of $\beta$, we shall find
\[ \psi\left(\beta+\frac{2k'\varpi i}{2n+1}+\frac{2k\omega}{2n+1}\right)=\overset{+n}{\underset{-n}{\Sigma}}_{\mu}\theta^{\mu}.\varphi_{1}\left(\beta+\frac{2(k'+\mu)\varpi i}{2n+1}+\frac{2k\omega}{2n+1}\right), \]
but by virtue of (73.) we have:
\[ \varphi_{1}\left(\beta+\frac{2(k'+\mu)\varpi i}{2n+1}+\frac{2k\omega}{2n+1}\right)=\varphi_{1}\left(\beta+\frac{2(k'+\mu)\varpi i}{2n+1}\right), \]
therefore
\[ \psi\left(\beta+\frac{2k'\varpi i}{2n+1}+\frac{2k\omega}{2n+1}\right)=\overset{+n}{\underset{-n}{\Sigma}}_{\mu}\theta^{\mu}.\varphi_{1}\left(\beta+\frac{2(k'+\mu)\varpi i}{2n+1}\right). \]
By virtue of (72.) we have
\[ \begin{aligned}
&\overset{+n}{\underset{-n}{\Sigma}}_{\mu}\theta^{\mu}.\varphi_{1}\left(\beta+\frac{2(k'+\mu)\varpi i}{2n+1}\right)\\
&=\theta^{-k'}.\overset{+n}{\underset{-n}{\Sigma}}_{\mu}\theta^{\mu}.\varphi_{1}\left(\beta+\frac{2\mu\varpi i}{2n+1}\right)+\overset{k'}{\underset{1}{\Sigma}}_{\mu}.\theta^{n+\mu-k'}.\varphi_{1}\left(\beta+\frac{2(\mu+n)\varpi i}{2n+1}\right)\\
&\quad-\overset{k'}{\underset{1}{\Sigma}}_{\mu}\theta^{\mu-n-1-k'}.\varphi_{1}\left(\beta+\frac{2(\mu-n-1)\varpi i}{2n+1}\right),
\end{aligned} \]
therefore, on remarking that $\theta^{n+\mu-k'}=\theta^{\mu-n-1-k'}$ and
\[ \varphi_{1}\left(\beta+\frac{2(\mu-n-1)\varpi i}{2n+1}\right)=\varphi_{1}\left(\beta+\frac{2(\mu+n)\varpi i}{2n+1}-2\varpi i\right)=\varphi_{1}\left(\beta+\frac{2(\mu+n)\varpi i}{2n+1}\right), \]
there will result
\[ \psi\left(\beta+\frac{2k'\varpi i}{2n+1}+\frac{2k\omega}{2n+1}\right)=\theta^{-k'}.\psi\beta. \tag{74.} \]
In the same manner we shall also find
\[ \psi_{1}\left(\beta+\frac{2k'\varpi i}{2n+1}+\frac{2k\omega}{2n+1}\right)=\theta^{k'}.\psi_{1}\beta. \]
These two equations will give
\[ \begin{aligned}
&\psi\left(\beta+\frac{2k\omega+2k'\varpi i}{2n+1}\right).\psi_{1}\left(\beta+\frac{2k\omega+2k'\varpi i}{2n+1}\right)=\psi\beta.\psi_{1}\beta,\\
&\left\{\psi\left(\beta+\frac{2k\omega+2k'\varpi i}{2n+1}\right)\right\}^{2n+1}+\left\{\psi_{1}\left(\beta+\frac{2k\omega+2k'\varpi i}{2n+1}\right)\right\}^{2n+1}=(\psi\beta)^{2n+1}+(\psi_{1}\beta)^{2n+1}.
\end{aligned} \]

\origpage{131}
By virtue of these equations we shall obtain, on putting in the values of $\lambda(\Phi\beta)$ and $\lambda_{1}(\Phi\beta)$, $\beta+\frac{2k\omega+2k'\varpi i}{2n+1}$ in place of $\beta$:
\[ \begin{aligned}
\lambda(\Phi\beta)&=\lambda\left[\Phi\left(\beta+\frac{2k\omega+2k'\varpi i}{2n+1}\right)\right],\\
\lambda_{1}(\Phi\beta)&=\lambda_{1}\left[\Phi\left(\beta+\frac{2k\omega+2k'\varpi i}{2n+1}\right)\right].
\end{aligned} \]
Now, $\Phi\left(\beta+\frac{2k\omega+2k'\varpi i}{2n+1}\right)$ expresses any root whatever of the equation
\[ \Phi(2n+1)\beta=\frac{P_{2n+1}}{Q_{2n+1}}. \]

Therefore, as we have said, the functions $\lambda(x)$ and $\lambda_{1}(x)$ will have the same values, whatever root be put in place of $x$.

Let then $x_{0}$, $x_{1}$, $x_{2}....x_{2n+1}$ be these roots; we shall have
\[ \begin{aligned}
\lambda(x)&=\frac{1}{2n+1}\{\lambda(x_{0})+\lambda(x_{1})+....+\lambda(x_{2n})\},\\
\lambda_{1}(x)&=\frac{1}{2n+1}\{\lambda_{1}(x_{0})+\lambda_{1}(x_{1})+....+\lambda_{1}(x_{2n})\}.
\end{aligned} \]
\ednote{The list of roots is printed as ending with $x_{2n+1}$, while the sums run only to $x_{2n}$.}
Now the second member of these equations is a \emph{rational and symmetric} function of the roots of the equation $\Phi(2n+1)\beta=\frac{P_{2n+1}}{Q_{2n+1}}$, therefore $\lambda(x)$ and $\lambda_{1}(x)$ can be expressed rationally in $\Phi(2n+1)\beta$. On making
\[ \lambda(x)=B,\ \lambda_{1}(x)=2A, \]
the equations (71.) will give
\[ (\psi\beta)^{2n+1}.(\psi_{1}\beta)^{2n+1}=B^{2n+1};\ (\psi\beta)^{2n+1}+(\psi_{1}\beta)^{2n+1}=2A, \]
whence one derives
\[ \psi\beta=\sqrt[2n+1]{[A+\sqrt{(A^{2}-B^{2n+1})}]}=\overset{+n}{\underset{-n}{\Sigma}}_{\mu}\theta^{\mu}.\Phi_{1}\left(\beta+\frac{2\mu\varpi i}{2n+1}\right). \tag{75.} \]

\subsection*{15.}

The value of $\psi\beta$ having been found, one will easily deduce from it those of $\Phi_{1}\beta$.

In fact, on taking for $\theta$ successively all the imaginary roots of the equation $\theta^{2n+1}-1=0$, and denoting the corresponding values of $A$ and $B$ by $A_{1}$, $B_{1}$, $A_{2}$, $B_{2}$ etc., we shall obtain:
\[ \begin{aligned}
\sqrt[2n+1]{[A_{1}+\sqrt{(A_{1}^{2}-B_{1}^{2n+1})}]}&=\overset{+n}{\underset{-n}{\Sigma}}_{\mu}\theta_{1}^{\mu}.\Phi_{1}\left(\beta+\frac{2\mu\varpi i}{2n+1}\right),\\
\sqrt[2n+1]{[A_{2}+\sqrt{(A_{2}^{2}-B_{2}^{2n+1})}]}&=\overset{+n}{\underset{-n}{\Sigma}}_{\mu}\theta_{2}^{\mu}.\Phi_{1}\left(\beta+\frac{2\mu\varpi i}{2n+1}\right),\\
\cdots\cdots\cdots&\cdots\cdots\cdots\\
\sqrt[2n+1]{[A_{2n}+\sqrt{(A_{2n}^{2}-B_{2n}^{2n+1})}]}&=\overset{+n}{\underset{-n}{\Sigma}}_{\mu}\theta_{2n}^{\mu}.\Phi_{1}\left(\beta+\frac{2\mu\varpi i}{2n+1}\right).
\end{aligned} \]

\origpage{132}
Likewise we know the sum of the roots:
\[ \overset{+n}{\underset{-n}{\Sigma}}_{m}\overset{+n}{\underset{-n}{\Sigma}}_{\mu}\varphi\left(\beta+\frac{2m\omega}{2n+1}+\frac{2\mu\varpi i}{2n+1}\right)=\overset{+n}{\underset{-n}{\Sigma}}_{\mu}\varphi_{1}\left(\beta+\frac{2\mu\varpi i}{2n+1}\right), \]
which is equal to $(2n+1)\varphi(2n+1)\beta$, as we shall see in what follows. On adding these equations member to member, after having multiplied the first by $\theta_{0}^{-k}$, the second by $\theta_{1}^{-k}$, the third by $\theta_{2}^{-k}....$ and the $(2n)^{\text{th}}$ by $\theta_{2n}^{-k}$, there will result:
\[ \begin{aligned}
&\overset{+n}{\underset{-n}{\Sigma}}_{\mu}.(\theta_{0}^{\mu-k}+\theta_{1}^{\mu-k}+\theta_{2}^{\mu-k}+....+\theta_{2n}^{\mu-k}+1).\varphi_{1}\left(\beta+\frac{2\mu\varpi i}{2n+1}\right)\\
&=(2n+1).\varphi(2n+1)\beta+\overset{2n}{\underset{1}{\Sigma}}_{\mu}\theta_{\mu}^{-k}.\sqrt[2n+1]{[A_{\mu}+\sqrt{(A_{\mu}^{2}-B_{\mu}^{2n+1})}]};
\end{aligned} \]
\ednote{In the first sum the term $\theta_{0}^{\mu-k}$ and the final $1$ are both printed, so 1 is counted twice; the next sum is printed without $\theta_{0}^{\mu-k}$.}
now the sum
\[ 1+\theta_{1}^{\mu-k}+\theta_{2}^{\mu-k}+....+\theta_{2n}^{\mu-k} \]
reduces to zero for all values of $k$, except for $k=\mu$. In this case it becomes equal to $2n+1$. Therefore the first member of the preceding equation becomes
\[ (2n+1).\varphi_{1}\left(\beta+\frac{2k\varpi i}{2n+1}\right), \]
therefore, on substituting and dividing by $(2n+1)$, we have:
\[ \begin{aligned}
\varphi_{1}\left(\beta+\frac{2k\varpi i}{2n+1}\right)=&\\
\varphi(2n+1)\beta+\frac{1}{2n+1}.&\left\{\theta_{1}^{-k}.\sqrt[2n+1]{[A_{1}+\sqrt{(A_{1}^{2}-B_{1}^{2n+1})}]}+\theta_{2}^{-k}\sqrt[2n+1]{[A_{2}+\sqrt{(A_{2}^{2}-B_{2}^{2n+1})}]}+....\right.\\
&\left.....+\theta_{2n}^{-k}\sqrt[2n+1]{[A_{2n}+\sqrt{(A_{2n}^{2}-B_{2n}^{2n+1})}]}\right\}.
\end{aligned} \tag{76.} \]
For $k=0$, we have:
\[ \begin{aligned}
\varphi_{1}\beta=\varphi(2n+1)\beta+\frac{1}{2n+1}&\left\{\sqrt[2n+1]{[A_{1}+\sqrt{(A_{1}^{2}-B_{1}^{2n+1})}]}+\sqrt[2n+1]{[A_{2}+\sqrt{(A_{2}^{2}-B_{2}^{2n+1})}]}+\ldots\right.\\
&\left.....+\sqrt[2n+1]{[A_{2n}+\sqrt{(A_{2n}^{2}-B_{2n}^{2n+1})}]}\right\}.
\end{aligned} \tag{77.} \]

\subsection*{16.}

The value of $\varphi_{1}\beta$ having thus been found, the question is to derive from it that of $\varphi\beta$. Now this can easily be done as follows:

Let
\[ \psi_{2}\beta=\overset{+n}{\underset{-n}{\Sigma}}_{m}\theta^{m}.\varphi\left(\beta+\frac{2m\omega}{2n+1}\right);\ \psi_{3}\beta=\overset{+n}{\underset{-n}{\Sigma}}_{m}\theta^{m}.\varphi\left(\beta-\frac{2m\omega}{2n+1}\right), \tag{78.} \]
we have
\[ \varphi\left(\beta\pm\frac{2m\omega}{2n+1}\right)=\frac{\varphi\beta.f\left(\frac{2m\omega}{2n+1}\right).F\left(\frac{2m\omega}{2n+1}\right)\pm f\beta.F\beta.\varphi\left(\frac{2m\omega}{2n+1}\right)}{1+e^{2}c^{2}.\varphi^{2}\left(\frac{2m\omega}{2n+1}\right).\varphi^{2}\beta}. \]

\origpage{133}
From this it follows that we may make
\[ \psi_{2}\beta=r+f\beta.F\beta.s;\ \psi_{3}\beta=r-f\beta.F\beta.s, \]
where $r$ and $s$ are rational functions of $\Phi\beta$.

From this we derive
\[ \left\{\begin{aligned}
\psi_{2}\beta.\psi_{3}\beta&=\chi(\Phi\beta),\\
(\psi_{2}\beta)^{2n+1}+(\psi_{3}\beta)^{2n+1}&=\chi_{1}(\Phi\beta),
\end{aligned}\right. \tag{79.} \]
$\chi(\Phi\beta)$ and $\chi_{1}(\Phi\beta)$ being two rational functions of $\Phi\beta$.

This being premised, I say that $\chi(\Phi\beta)$ and $\chi_{1}(\Phi\beta)$ can be expressed rationally in $\Phi_{1}\beta$.

We have seen that
\[ \varphi_{1}\beta=\varphi\beta+\overset{n}{\underset{1}{\Sigma}}_{m}.\left\{\frac{2\varphi\beta.f\left(\frac{2m\omega}{2n+1}\right).F\left(\frac{2m\omega}{2n+1}\right)}{1+e^{2}c^{2}.\varphi^{2}\left(\frac{2m\omega}{2n+1}\right)\varphi^{2}\beta}\right\}. \tag{80.} \]
On making $\Phi\beta=x$, we shall have an equation in $x$ of the degree $(2n+1)$. One root of this equation is $x=\Phi\beta$; now, on putting $\beta+\frac{2k\omega}{2n+1}$ in place of $\beta$, $\Phi_{1}\beta$ does not change its value, therefore $x=\Phi\left(\beta+\frac{2k\omega}{2n+1}\right)$ will be a root, whatever the integer $k$ may be. Now, on giving to $k$ all the integer values from $-n$ up to $+n$, $\Phi\left(\beta+\frac{2k\omega}{2n+1}\right)$ will take $2n+1$ different values, therefore these $2n+1$ quantities will be precisely the $2n+1$ roots of the equation in $x$.

This being premised, on putting $\beta+\frac{2k\omega}{2n+1}$ in place of $\beta$ in the expression of $\psi_{2}\beta$, there will result by virtue of (72.):
\[ \begin{aligned}
\psi_{2}\left(\beta+\frac{2k\omega}{2n+1}\right)&=\overset{+n}{\underset{-n}{\Sigma}}_{m}\theta^{m}\Phi\left(\beta+\frac{2(k+m)\omega}{2n+1}\right)\\
&=\theta^{-k}.\psi_{2}3+\overset{k}{\underset{1}{\Sigma}}_{m}\theta^{m+n-k}.\Phi\left(\beta+\frac{2(m+n)\omega}{2n+1}\right)\\
&\quad-\overset{k}{\underset{1}{\Sigma}}_{m}\theta^{m-n-1-k}.\Phi\left(\beta+\frac{2(m-n-1)\omega}{2n+1}\right),
\end{aligned} \]
\ednote{Printed $\psi_{2}3$ with a digit three instead of $\psi_{2}\beta$; an apparent misprint.}
therefore, on paying attention that $\theta^{m+n-k}=\theta^{m-n-1-k}$ and $\Phi\left(\beta+\frac{2(m-n-1)}{2n+1}\omega\right)=\Phi\left(\beta+\frac{2(m+n)\omega}{2n+1}\right)$, it will result
\[ \psi_{2}\left(\beta+\frac{2k\omega}{2n+1}\right)=\theta^{-k}.\psi_{2}\beta. \tag{81.} \]

\origpage{134}
Likewise we shall have
\[ \psi_{3}\left(\beta+\frac{2k\omega}{2n+1}\right)=\theta^{+k}.\psi_{3}\beta. \]
One sees by virtue of these relations that the equations which give the values of the functions $\chi(\Phi\beta)$ and $\chi_{1}(\Phi\beta)$\ednote{The second printed argument looks like a damaged or misprinted $\beta$.}, lead to these two equalities:
\[ \begin{aligned}
\chi_{1}\left[\Phi\left(\beta+\frac{2k\omega}{2n+1}\right)\right]&=\chi_{1}(\Phi\beta),\\
\chi\left[\Phi\left(\beta+\frac{2k\omega}{2n+1}\right)\right]&=\chi(\Phi\beta).
\end{aligned} \]
From this we derive
\[ \begin{aligned}
\chi(\Phi\beta)&=\frac{1}{2n+1}.\overset{+n}{\underset{-n}{\Sigma}}_{k}.\chi\left[\Phi\left(\beta+\frac{2k\omega}{2n+1}\right)\right],\\
\chi_{1}(\Phi\beta)&=\frac{1}{2n+1}.\overset{+n}{\underset{-n}{\Sigma}}_{k}.\chi_{1}\left[\Phi\left(\beta+\frac{2k\omega}{2n+1}\right)\right].
\end{aligned} \]

Now, these values of $\chi(\Phi\beta)$ and $\chi_{1}(\Phi\beta)$ are rational and symmetric functions of all the roots of the equation (80.). Therefore they can be expressed rationally by the coefficients of the same equation, that is rationally in $\Phi_{1}\beta$.

Let
\[ \chi_{1}(\Phi\beta)=2C,\ \chi(\Phi\beta)=D, \]
\ednote{The subscripts of both printed $\chi$ are damaged and hard to read; they are read here as $\chi_1$ for $2C$ and $\chi$ for $D$, as (79.) requires.}
the equations (79.) will give
\[ \psi_{2}\beta=\sqrt[2n+1]{[C+\sqrt{(C^{2}-D^{2n+1})}]}, \]
whence, on restoring the value of $\psi_{2}\beta$:
\[ \sqrt[2n+1]{[C+\sqrt{(C^{2}-D^{2n+1})}]}=\overset{+n}{\underset{-n}{\Sigma}}_{m}\theta^{m}.\Phi\left(\beta+\frac{2m\omega}{2n+1}\right). \tag{82.} \]
From this we derive, on putting $\theta_{\mu}$ in place of $\theta$, and denoting the corresponding values of $C$ and $D$ by $C_{\mu}$ and $D_{\mu}$:
\[ \theta_{\mu}^{-k}\sqrt[2n+1]{[C_{\mu}+\sqrt{(C_{\mu}^{2}-D_{\mu}^{2n+1})}]}=\overset{+n}{\underset{-n}{\Sigma}}_{m}\theta_{\mu}^{m-k}.\Phi\left(\beta+\frac{2m\omega}{2n+1}\right). \]
On joining to it the equation
\[ \Phi_{1}\beta=\overset{+n}{\underset{-n}{\Sigma}}_{m}\Phi\left(\beta+\frac{2m\omega}{2n+1}\right), \]
one will easily derive from it:
\[ (2n+1).\Phi\left(\beta+\frac{2k\omega}{2n+1}\right)=\Phi_{1}\beta+\overset{2n}{\underset{1}{\Sigma}}_{\mu}\theta_{\mu}^{-k}.\sqrt[2n+1]{[C_{\mu}+\sqrt{(C_{\mu}^{2}-D_{\mu}^{2n+1})}]}. \tag{83.} \]
Supposing $k=0$, there will result:
\[ \Phi\beta=\frac{1}{2n+1}\left\{\Phi_{1}\beta+\sqrt[2n+1]{[C_{1}+\sqrt{(C_{1}^{2}-D_{1}^{2n+1})}]}+\ldots+\sqrt[2n+1]{[C_{2n}+\sqrt{(C_{2n}^{2}-D_{2n}^{2n+1})}]}\right\}. \tag{84.} \]

\origpage{135}

This equation gives $\Phi\beta$ as an algebraic function of $\Phi_{1}\beta$, and previously we found $\Phi_{1}\beta$ as an algebraic function of $\Phi(2n+1)\beta$. Therefore, on putting $\frac{\alpha}{2n+1}$ in place of $\beta$, we shall have $\Phi\left(\frac{\alpha}{2n+1}\right)$ as an algebraic function of $\Phi\alpha$.

By an entirely similar analysis one will find $f\left(\frac{\alpha}{2n+1}\right)$ in $f\alpha$ and $F\left(\frac{\alpha}{2n+1}\right)$ in $F\alpha$.

\subsection*{17.}

The values which we have just found for the quantities $\Phi_{1}\beta$ and $\Phi\beta$, the first in $\Phi(2n+1)\beta$ and the second in $\Phi_{1}\beta$, each contain the sum of $2n$ different radicals of the $(2n+1)$ degree. There will result for $\Phi\beta$, $\Phi_{1}\beta....$ a number $(2n+1)^{2n}$ of values, whereas each of these quantities is the root of an equation of the $(2n+1)^{\text{th}}$ degree. Now one can give to the expressions of $\Phi\beta$ and $\Phi_{1}\beta$ such a form that the number of the values of these quantities shall be precisely equal to $2n+1$.

For this let
\[ \theta=\cos\frac{2\pi}{2n+1}+i\sin\frac{2\pi}{2n+1}, \]
we may make
\[ \theta_{1}=\theta,\ \theta_{2}=\theta^{2},\ \theta_{3}=\theta^{3}....\theta_{2n}=\theta^{2n}. \]

Let likewise
\[ \left\{\begin{aligned}
\psi^{k}(\beta)&=\overset{+n}{\underset{-n}{\Sigma}}_{\mu}\theta^{k\mu}.\Phi_{1}\left(\beta+\frac{2\mu\varpi i}{2n+1}\right),\\
\psi_{1}^{k}(\beta)&=\overset{+n}{\underset{-n}{\Sigma}}_{\mu}\theta^{k\mu}.\Phi_{1}\left(\beta-\frac{2\mu\varpi i}{2n+1}\right),
\end{aligned}\right. \tag{85.} \]
we shall have by virtue of (74.)
\[ \begin{aligned}
\psi^{k}\left(\beta+\frac{2\nu\varpi i}{2n+1}\right)&=\theta^{-k\nu}.\psi^{k}(\beta),\\
\psi_{1}^{k}\left(\beta+\frac{2\nu\varpi i}{2n+1}\right)&=\theta^{+k\nu}.\psi_{1}^{k}(\beta),\\
\psi^{1}\left(\beta+\frac{2\nu\varpi i}{2n+1}\right)&=\theta^{-\nu}.\psi^{1}(\beta),\\
\psi_{1}^{1}\left(\beta+\frac{2\nu\varpi i}{2n+1}\right)&=\theta^{+\nu}.\psi_{1}^{1}(\beta).
\end{aligned} \]

Let now
\[ \left\{\begin{aligned}
\frac{\psi^{k}(\beta)}{(\psi^{1}\beta)^{k}}+\frac{\psi_{1}^{k}(\beta)}{(\psi_{1}^{1}\beta)^{k}}&=P(\Phi\beta),\\
\frac{\psi^{k}(\beta)}{(\psi^{1}\beta)^{k-2n-1}}+\frac{\psi_{1}^{k}(\beta)}{\psi_{1}^{1}(\beta)^{k-2n-1}}&=Q(\Phi\beta),
\end{aligned}\right. \tag{86.} \]

\origpage{136}
$P(\Phi\beta)$ and $Q(\Phi\beta)$ will be rational functions of $\Phi\beta$; now, on putting $\beta+\frac{2m\omega+2\mu\varpi i}{2n+1}$ in place of $\beta$, it is clear by virtue of the formulas that $P$ and $Q$ do not change their values; therefore we shall have
\[ P(\Phi\beta)=\frac{1}{(2n+1)^{2}}.\overset{+n}{\underset{-n}{\Sigma}}_{m}\overset{+n}{\underset{-n}{\Sigma}}_{\mu}.P\left[\Phi\left(\beta+\frac{2m\omega+2\mu\varpi i}{2n+1}\right)\right]; \]
now, the second member being a symmetric and rational function of the roots of the equation $\Phi(2n+1)\beta=\frac{P_{2n+1}}{Q_{2n+1}}$, $P(\Phi\beta)$ can be expressed rationally in $\Phi(2n+1)\beta$. The same holds of $Q(\Phi\beta)$. These two quantities being known, the equations (86.) will give
\[ \frac{\psi^{k}(\beta)}{(\psi^{1}\beta)^{k}}\left\{1-\left(\frac{\psi_{1}^{1}\beta}{\psi_{1}^{1}\beta}\right)^{2n+1}\right\}=P(\Phi\beta)-\frac{Q(\varphi\beta)}{(\psi_{1}^{1}\beta)^{2n+1}}, \]
now
\[ \begin{aligned}
\psi_{1}^{1}\beta&=A_{1}+\sqrt{(A_{1}^{2}-B_{1}^{2n+1})},\\
\psi_{1}^{1}\beta&=A_{1}-\sqrt{(A_{1}^{2}-B_{1}^{2n+1})},
\end{aligned} \]
\ednote{In the first of these two equations the subscript of the printed $\psi$ is damaged; it should read $\psi^{1}\beta$ without subscript, and the quotient in the preceding display should likewise have $\psi^{1}\beta$ in the numerator.}
therefore
\[ \frac{\psi^{k}(\beta)}{(\psi^{1}\beta)^{k}}.2\sqrt{(A_{1}^{2}-B_{1}^{2n+1})}=Q(\Phi\beta)-[A_{1}-\sqrt{(A_{1}^{2}-B_{1}^{2n+1})}].P(\Phi\beta). \]
Therefore we shall have
\[ \psi^{k}(\beta)=(\psi^{1}\beta)^{k}.\{F_{k}+H_{k}.\sqrt{(A_{1}^{2}-B_{1}^{2n+1})}\}, \]
where $F_{k}$ and $H_{k}$ are rational functions of $\Phi(2n+1)\beta$. On replacing $A_{1}$ and $B_{1}$ and substituting the values of $\psi^{k}(\beta)$ and $(\psi^{1}\beta)^{k}$, there will result:
\[ \sqrt[2n+1]{[A_{k}+\sqrt{(A_{k}^{2}-B_{k}^{2n+1})}]}=[A+\sqrt{(A^{2}-B^{2n+1})}]^{\frac{k}{2n+1}}[F_{k}+H_{k}.\sqrt{(A^{2}-B^{2n+1})}], \]
therefore the value of $\Phi_{1}\beta$ will become:
\[ \begin{aligned}
\Phi_{1}\beta=\Phi(2n+1)\beta+\frac{1}{2n+1}.&\left\{[A+\sqrt{(A^{2}-B^{2n+1})}]^{\frac{1}{2n+1}}\right.\\
&+[F_{2}+H_{2}\sqrt{(A^{2}-B_{2n}^{2n+1})}]\,[A+\sqrt{(A^{2}-B^{2n+1})}]^{\frac{2}{2n+1}}\\
&\left.+....[F_{2n}+H_{2n}\sqrt{(A^{2}-B^{2n+1})}]\,[A+\sqrt{(A^{2}-B^{2n+1})}]^{\frac{2n}{2n+1}}\right\}.
\end{aligned} \tag{87.} \]
\ednote{In the second line the printed $B$ under the radical carries the subscript $2n$, which the first and third lines lack; an apparent misprint.}
By an entirely similar procedure we shall find
\[ \begin{aligned}
\Phi\beta=\frac{1}{2n+1}&\left\{\Phi_{1}\beta+[C+\sqrt{(C^{2}-D^{2n+1})}]^{\frac{1}{2n+1}}\right.\\
&+[K_{2}+L_{2}\sqrt{(C^{2}-D^{2n+1})}]\,[C+\sqrt{(C^{2}-D^{2n+1})}]^{\frac{2}{2n+1}}\\
&\left.+....[K_{2n}+L_{2n}\sqrt{(C^{2}-D^{2n+1})}]\,[C+\sqrt{(C^{2}-D^{2n+1})}]^{\frac{2n}{2n+1}}\right\},
\end{aligned} \tag{88.} \]
where $K_{2}$, $L_{2}$, $K_{3}$, $L_{3}....K_{2n}$, $L_{2n}$ are rational functions of $\Phi_{1}\beta$.

\origpage{137}

These expressions of $\Phi_{1}\beta$ and $\Phi\beta$ have only $2n+1$ different values, which one will obtain by attributing to the radicals their $2n+1$ values. --- It follows from our analysis that one may take $\sqrt{(A^{2}-B^{2n+1})}$ and $\sqrt{(C^{2}-D^{2n+1})}$ with whatever sign one pleases.

\subsection*{18.}

The value which we have found for $\Phi(\beta)$ or $\Phi\left(\frac{\alpha}{2n+1}\right)$ contains, besides the function $\Phi\alpha$, the following:
\[ \begin{gathered}
e,\ c,\ \theta,\\
\Phi\left(\frac{m\omega}{2n+1}\right),\ \Phi\left(\frac{m\varpi i}{2n+1}\right),\ f\left(\frac{m\omega}{2n+1}\right),\\
f\left(\frac{m\varpi i}{2n+1}\right),\ F\left(\frac{m\omega}{2n+1}\right),\ F\left(\frac{m\varpi i}{2n+1}\right),
\end{gathered} \]
for any values whatever of $m$ from $1$ to $2n$. Now whatever the value of $m$ may be, one can always express algebraically $\Phi\left(\frac{m\omega}{2n+1}\right)$, $f\left(\frac{m\omega}{2n+1}\right)$, $F\left(\frac{m\omega}{2n+1}\right)$ in $\Phi\left(\frac{\omega}{2n+1}\right)$, and $\Phi\left(\frac{\mu\varpi i}{2n+1}\right)$, $f\left(\frac{\mu\varpi i}{2n+1}\right)$, $F\left(\frac{\mu\varpi i}{2n+1}\right)$ in $\Phi\left(\frac{\varpi i}{2n+1}\right)$. Everything is therefore known in the expression of $\Phi\left(\frac{\alpha}{2n+1}\right)$, except the two quantities independent of $\alpha$; $\Phi\left(\frac{\omega}{2n+1}\right)$, $\Phi\left(\frac{\varpi i}{2n+1}\right)$. --- These quantities depend only on $c$ and $e$, and they can be found by the resolution of an equation of the degree $(2n+1)^{2}-1$, namely of the equation $\frac{P_{2n+1}}{x}=0$. --- We are going to see in the following paragraph how one can reduce its resolution to that of lower equations.

\section*{§. V.}

C. \emph{On the equation} $P_{2n+1}=0$.

\subsection*{19.}

The expression which we have just found for $\Phi\left(\frac{\alpha}{2n+1}\right)$ will contain, as we have seen, the two constant quantities $\Phi\left(\frac{\omega}{2n+1}\right)$ and $\Phi\left(\frac{\varpi i}{2n+1}\right)$. We shall find these quantities by resolving the equation
\[ P_{2n+1}=0, \]
whose roots will be represented by
\[ x=\Phi\left(\frac{m\omega+\mu\varpi i}{2n+1}\right), \]

\origpage{138}
where $m$ and $\mu$ may be all the integers from $-n$ up to $+n$. One of these roots, which answers to $m=0$, $\mu=0$, is equal to zero. Therefore $P_{2n+1}$ is divisible by $x$. On removing this factor, we shall have an equation
\[ R=0,\ \text{of the degree }(2n+1)^{2}-1. \]
On making $x^{2}=r$, the equation (90.) $R=0$, in $r$, will be of the degree $\frac{(2n+1)^{2}-1}{2}=2.n(n+1)$, and the roots of this equation will be
\[ r=\Phi^{2}\left(\frac{m\omega\pm\mu\varpi i}{2n+1}\right), \tag{91.} \]
$\mu$ and $m$ having all positive values below $n$, setting aside the root zero.

We are now going to see how one can reduce the resolution of the equation $R=0$ to that of two equations, the one of the degree $n$ and the other of the degree $2n+2$.

First, I say that one can represent all the values of $r$ by
\[ \Phi^{2}\left(\frac{m\omega}{2n+1}\right)\ \text{and}\ \Phi^{2}\left(m\frac{\mu\omega+\varpi i}{2n+1}\right)\ldots\ldots, \tag{92.} \]
on giving to $\mu$ all the integer values from zero up to $2n$, and to $m$ all those from $1$ to $n$.

In fact $\Phi^{2}\left(\frac{m\omega}{2n+1}\right)$ represents first a number $n$ of values of $r$; now the others can be represented by $\Phi^{2}\left(m.\frac{\mu\omega+\varpi i}{2n+1}\right)$. Let, to demonstrate it, $m\mu=(2n+1)k+m'$, where $m'$ is an integer lying between the limits $-n$ and $+n$. On substituting, we shall have
\[ \begin{aligned}
\Phi^{2}\left(m.\frac{\mu\omega+\varpi i}{2n+1}\right)&=\Phi^{2}\left(k.\omega+\frac{m'\omega+\mu\varpi i}{2n+1}\right)\\
&=\Phi^{2}\left(\frac{m'\omega+\mu\varpi i}{2n+1}\right)=\Phi^{2}\left(\frac{-m'\omega-\mu\varpi i}{2n+1}\right).
\end{aligned} \]
$\Phi^{2}\left(m.\frac{\mu\omega+\varpi i}{2n+1}\right)$ is therefore a value of $r$; now to each value of $m$ there answers a different value of $m'$. For if we had
\[ m_{1}\mu=(2n+1)k_{1}+m', \]
it would follow
\[ (m_{1}-m)\mu=(2n+1).(k-k_{1}), \]
which is impossible, on remarking that $2n+1$ is a prime number. Therefore $\Phi^{2}\left(m.\frac{\mu\omega+\varpi i}{2n+1}\right)$ combined with $\Phi^{2}\left(\frac{m\omega}{2n+1}\right)$ represents all the values of $r$.

\origpage{139}

This being premised, let
\[ \begin{aligned}
&\left[r-\Phi^{2}\left(\frac{\omega}{2n+1}\right)\right]\left[r-\Phi^{2}\left(\frac{2\omega}{2n+1}\right)\right]\ldots\left[r-\Phi^{2}\left(\frac{n\omega}{2n+1}\right)\right]\\
&=r^{n}+p_{n-1}.r^{n-1}+p_{n-2}.r^{n-2}+\ldots\ldots+p_{1}.r+p_{0}.
\end{aligned} \tag{93.} \]
The quantities $p_{0}$, $p_{1}$, $\ldots p_{n-1}$, will be rational and symmetric functions of $\Phi^{2}\left(\frac{\omega}{2n+1}\right)$, $\Phi^{2}\left(\frac{2\omega}{2n+1}\right)\ldots\Phi^{2}\left(\frac{n\omega}{2n+1}\right)$; now these functions can be found by means of an equation of degree $2n+2$.

Let $p$ be any rational and symmetric function whatever of $\Phi^{2}\left(\frac{\omega'}{2n+1}\right)$, $\Phi^{2}\left(\frac{2\omega'}{2n+1}\right)\ldots\Phi^{2}\left(\frac{n\omega'}{2n+1}\right)$, where $\omega'$ denotes the quantity $m\omega+\mu\varpi i$.

By the formulas which we have given above for expressing $\Phi(n\beta)$ in $\Phi\beta$, it is clear that one can express $\Phi^{2}\left(m'.\frac{\omega'}{2n+1}\right)$ as a rational function of $\Phi^{2}\left(\frac{\omega'}{2n+1}\right)$. Therefore we may make
\[ p=\psi.\left[\Phi^{2}\left(\frac{\omega'}{2n+1}\right)\right]=\theta\left[\Phi^{2}\left(\frac{\omega'}{2n+1}\right),\ \Phi^{2}\left(\frac{2\omega'}{2n+1}\right),\ \ldots\ \Phi^{2}\left(\frac{n\omega'}{2n+1}\right)\right], \tag{94.} \]
$\theta$ denoting a symmetric and rational function. On putting $\nu\omega'$ in place of $\omega'$, there will result
\[ \psi\left[\Phi^{2}\left(\frac{\nu\omega'}{2n+1}\right)\right]=\theta\left[\Phi^{2}\left(\frac{\nu\omega'}{2n+1}\right),\ \Phi^{2}\left(\frac{2\nu\omega'}{2n+1}\right),\ \ldots\ \Phi^{2}\left(\frac{n\nu\omega'}{2n+1}\right)\right], \tag{95.} \]
now, on making:
\[ a.\nu=(2n+1).k_{a}'+k_{a}, \]
where $k_{a}$ is an integer and lies between $-n$ and $+n$, the series
\[ k_{1},\ k_{2}\ \ldots\ldots\ k_{n} \]
will have, apart from the sign, the same terms as this one:
\[ 1,\ 2,\ 3\ \ldots\ldots\ n; \]
therefore it is clear that the second member of the equation (95.) will have the same value as $p$. Therefore:
\[ \psi\left[\Phi^{2}\left(\frac{\nu\omega'}{2n+1}\right)\right]=\psi\left[\Phi^{2}\left(\frac{\omega'}{2n+1}\right)\right], \tag{96.} \]
\ednote{The exponent of the second Phi in (96.) is damaged in the print and appears as a raised dot; read as the exponent 2 of the first member.}
an equation which, on making $\omega'=\omega$ and $\omega'=m\omega+\varpi i$, will give these two:
\[ \left\{\begin{aligned}
&\psi\left[\Phi^{2}\left(\frac{\nu\omega}{2n+1}\right)\right]=\psi.\left[\Phi^{2}\left(\frac{\omega}{2n+1}\right)\right],\\
&\psi\left[\Phi^{2}\left(\nu.\frac{m\omega+\varpi i}{2n+1}\right)\right]=\psi\left[\Phi^{2}\left(\frac{m\omega+\varpi i}{2n+1}\right)\right],
\end{aligned}\right. \tag{97.} \]
or else, on making, for brevity,
\[ \Phi^{2}\left(\frac{\nu\omega}{2n+1}\right)=r_{\nu},\ \Phi^{2}\left(\nu.\frac{m\omega+\varpi i}{2n+1}\right)=r_{\nu,m}, \tag{98.} \]

\origpage{140}
there will result:
\[ \psi.r_{\nu}=\psi r_{1};\ \psi r_{\nu,m}=\psi r_{1,m}. \tag{99.} \]

This being premised, let
\[ \left\{\begin{aligned}
&(p-\psi r_{1})(p-\psi r_{1,0})(p-\psi r_{1,1})(p-\psi r_{1,2})\ldots\ (p-\psi r_{1,2n})\\
&=q_{0}+q_{1}.p+q_{2}.p^{2}+\ldots q_{2n+1}.p^{2n+1}+p^{2n}.
\end{aligned}\right. \tag{100.} \]
\ednote{The exponent of the last term is partly obscured by a smudge and reads $2n$; the degree $2n+2$ is expected from (104.).}

I say that the coefficients $q_{0}$, $q_{1}$ etc. can be expressed rationally in terms of $e$ and $c$.

First, by virtue of the known formulas, these coefficients can be expressed rationally in terms of $t_{1}$, $t_{2}\ldots t_{2n}$, if one makes, for brevity,
\[ t_{k}=(\psi r_{1})^{k}+(\psi r_{1,0})^{k}+(\psi r_{1,1})^{k}+\ldots+(\psi r_{1,2n})^{k}. \tag{101.} \]

It is therefore a question of finding the quantities $t_{1}$, $t_{2}\ldots$, now this can easily be done by means of the relations (99.). Indeed, on making therein successively $\nu=1$, $2\ldots n$, after having raised both members to the $k^{\text{th}}$ power, one will deduce at once:
\[ \left\{\begin{aligned}
(\psi r_{1})^{k}&=\frac{1}{n}\left[(\psi r_{1})^{k}+(\psi r_{2})^{k}+\ldots+(\psi r_{n})^{k}\right],\\
(\psi r_{1,m})^{k}&=\frac{1}{n}\left[(\psi r_{1,m})^{k}+(\psi r_{2,m})^{k}+\ldots+(\psi r_{n,m})^{k}\right].
\end{aligned}\right. \tag{102.} \]
Therefore, putting for $m$ all the integral numbers $0$, $1$, $\ldots$ $2n$, and then substituting in the expression for $t_{k}$, there will result:
\[ \left\{\begin{aligned}
n.t_{k}&=(\psi r_{1})^{k}+(\psi r_{2})^{k}+\ldots\ldots+(\psi r_{n})^{k}\\
&+(\psi r_{1,0})^{k}+(\psi r_{2,0})^{k}+\ldots\ldots+(\psi r_{n,0})^{k}\\
&+(\psi r_{1,1})^{k}+(\psi r_{2,1})^{k}+\ldots\ldots+(\psi r_{n,1})^{k}\\
&+\cdots\cdots\cdots\cdots\cdots\cdots\cdots\\
&+(\psi r_{1,2n})^{k}+(\psi r_{2,2n})^{k}+\ldots\ldots+(\psi r_{n,2n})^{k}.
\end{aligned}\right. \tag{103.} \]

This value of $t_{k}$ is, as one sees, a rational and symmetric function of the $2n+2$ quantities $r_{1}$, $r_{2}$, $\ldots r_{n}$; $r_{1,0}$, $r_{2,0}$, $\ldots r_{n,0}\ldots$ $r_{1,2n}$, $r_{2,2n}\ldots r_{n,2n}$, which are the $2n+2$ roots of the equation $R=0$. Therefore, as is known, $t_{k}$ can be expressed rationally by the coefficients of this equation, and consequently as a rational function of $e$ and $c$. Having thus found the quantities $t_{k}$, one deduces from them the values of $q_{0}$, $q_{1}$, $\ldots$ $\ldots q_{2n+1}$; which will likewise be rational functions of $e$ and $c$.

\subsection*{20.}

This being premised, on supposing
\[ 0=q_{0}+q_{1}.p+q_{2}.p^{2}+\ldots\ldots+q_{2n+1}.p^{2n+1}+p^{2n+2}, \tag{104.} \]
one will have an equation of the $(2n+2)^{\text{th}}$ degree, whose roots will be
\[ \psi r_{1};\ \psi r_{1,0};\ \psi r_{1,1};\ \psi r_{1,2}\ \ldots\ldots\ \psi r_{1,2n}. \]

\origpage{141}

The function $\psi r_{1}$, that is, any rational and symmetric function whatever of the roots $r_{1}$, $r_{2}$, $r_{3}$, $\ldots\ldots r_{n}$, can therefore be found by means of an equation of the degree $2n+2$.

Therefore one will in this manner have the coefficients $p_{0}$, $p_{1}$, $\ldots\ldots p_{n-1}$, by solving a number $n$ of equations, each of the $(2n+2)^{\text{th}}$ degree.

Having determined $p_{0}$, $p_{1}$, $\ldots\ldots$, one will have, by solving the equation
\[ 0=p_{0}+p_{1}r+\ldots\ldots+p_{n-1}.r^{n-1}+r^{n}, \tag{105.} \]
the value of the quantities
\[ r_{1},\ r_{2},\ \ldots\ldots r_{n};\ r_{1,0},\ r_{2,0},\ \ldots\ldots r_{n,0};\ r_{1,1},\ r_{2,1},\ \ldots\ldots r_{n,1}\ \text{etc. etc.,} \]
of which the first is equal to $\Phi^{2}\left(\frac{\omega}{2n+1}\right)$. Therefore the determination of this quantity, or rather the solution of the equation $R=0$, which is of the degree $(2n+2).n$, is reduced to that of equations of the degrees $(2n+2)$ and $n$.

But the preceding procedure can still be simplified. Indeed, as we shall see, in order to have the quantities $p_{0}$, $p_{1}$, $\ldots\ldots$, it suffices to know any one of them, and then the others can be expressed rationally by that one.

Let $p$, $q$ be, in general, two rational and symmetric functions of the quantities $r_{1}$, $r_{2}$, $\ldots\ldots r_{n}$; one can make, as we have seen:
\[ p=\psi r_{1};\ q=\theta r_{1}, \]
$\psi r_{1}$ and $\theta r_{1}$ denote two rational functions of $r_{1}$, which have this property: of remaining the same, if one changes $r_{1}$ into any other of the quantities $r_{1}$, $r_{2}$, $\ldots\ldots r_{n}$.

Suppose now:
\[ s_{k}=(\psi r_{1})^{k}.\theta r_{1}+(\psi r_{1,0})^{k}.\theta r_{1,0}+(\psi r_{1,1})^{k}.\theta r_{1,1}+\ldots\ldots+(\psi r_{1,2n})^{k}.\theta r_{1,2n}, \]
I say that $s_{k}$ can be expressed rationally in terms of $e$ and $c$.

Indeed, one has
\[ (\psi r_{1})^{k}.\theta r_{1}=(\psi r_{\nu})^{k}.\theta r_{\nu}=\frac{1}{n}\left\{(\psi r_{1})^{k}.\theta r_{1}+(\psi r_{2})^{k}.\theta r_{2}+\ldots\ldots+(\psi r_{n})^{k}.\theta r_{n}\right\}, \]
\[ (\psi r_{1,m})^{k}.\theta r_{1,m}=(\psi r_{\nu,m})^{k}.\theta r_{\nu,m}=\frac{1}{n}\left\{(\psi r_{1,m})^{k}.\theta r_{1,m}+(\psi r_{2,m})^{k}.\theta r_{2,m}+\ldots+(\psi r_{n,m})^{k}.\theta r_{n,m}\right\}. \]
On making $m=0$, $1$, $2$, $\ldots\ldots 2n$, and substituting in the expression for $s_{k}$, it will be seen that $s_{k}$ will be a rational and symmetric function of the roots $r_{0}$, $r_{1}$, $\ldots\ldots r_{1,0}$ etc., etc. of the equation $R=0$; therefore $s_{k}$ can be expressed rationally in terms of $e$ and $c$.

Knowing $s_{k}$, one will obtain, on making $k=0$, $1$, $2$, $\ldots\ldots 2n$, $2n+1$ equations, from which one will easily deduce the value of $\theta r_{1}$ as a function

\origpage{142}
rational of $\psi_{1}r$.\ednote{Printed $\psi_{1}r$ with the subscript on the psi; $\psi r_{1}$ is expected, as on the previous page.} Therefore, a function of the form $p$ being given, one can express any other function whatever of the same form as a rational function of $p$. Therefore, as we have said, one can express the coefficients $p_{0}$, $p_{1}$, $\ldots\ldots p_{n-1}$ rationally by any one of them whatever. Therefore, finally, in order to have their values, it suffices to solve a single equation of the degree $2n+2$, and consequently, in order to have the roots of the equation $R=0$, it suffices to solve one equation of the degree $2n+2$, and $2n+2$ equations of the degree $n$.

\subsection*{21.}

Now, among the equations on which the determination of the quantities $\varphi\left(\frac{\omega}{2m+1}\right)$; $\varphi\left(\frac{\varpi i}{2n+1}\right)$,\ednote{Printed $2m+1$ in the first denominator; $2n+1$ is apparently expected.} those of the degree $n$ can be solved algebraically. The procedure by which we shall effect this solution is entirely similar to that which is due to Mr.\ \emph{Gauss} for the solution of the equation
\[ \theta^{2n+1}-1=0. \]
Let the equation be proposed
\[ 0=p_{1}+p_{1}.r+p_{2}.r^{2}+\ldots\ldots+p_{n-1}.r^{n-1}+r^{n}, \tag{106.} \]
\ednote{The first coefficient is printed $p_{1}$; $p_{0}$ is expected.}
whose roots are:
\[ \Phi^{2}\left(\frac{\omega'}{2n+1}\right);\ \Phi^{2}\left(\frac{2\omega'}{2n+1}\right);\ \ldots\ldots\ \Phi^{2}\left(\frac{n\omega'}{2n+1}\right), \]
where $\omega'$ has one of the values of $\omega$, $m\omega+\varpi i$. Let us denote by $\alpha$ one of the primitive roots of the number $2n+1$, that is, an integer such that $\mu=2n$\uncertain{$\,+1$}\ednote{The characters after $2n$ are overstruck by a heavy bar and hard to read; $\mu=2n$ is expected.} is the smallest number which renders $\alpha^{\mu}-1$ divisible by $2n+1$, I say that the roots of the equation (106.) can also be represented by
\[ \Phi^{2}(\varepsilon),\ \Phi^{2}(\alpha\varepsilon),\ \Phi^{2}(\alpha^{2}\varepsilon),\ \Phi^{2}(\alpha^{3}\varepsilon)\ldots\ \Phi^{2}(\alpha^{n-1}.\varepsilon), \tag{107.} \]
where $\varepsilon=\frac{\omega'}{2n+1}$.

Let
\[ \alpha^{m}=(2n+1)k_{m}\pm a_{m}, \]
where $k$ is an integer and $a_{m}$ an integer, positive and less than $n+1$, I say that the terms of the series
\[ 1,\ a_{1},\ a_{2},\ \ldots\ldots\ a_{n-1} \]
will all be different from one another.

Indeed, if one has
\[ a_{m}=a_{\mu}, \]
it follows,

\origpage{143}
\[ \begin{aligned}
\text{or }\ \alpha^{m}-\alpha^{\mu}&=(2n+1)(k_{m}-k_{\mu}),\\
\text{or }\ \alpha^{m}+\alpha^{\mu}&=(2n+1)(k_{m}+k_{\mu}).
\end{aligned} \]
It is therefore necessary that one of the quantities $\alpha^{m}-\alpha^{\mu}$, $\alpha^{m}+\alpha^{\mu}$ be divisible by $2n+1$; now let $m>\mu$ be posited, which is permitted, it is necessary that $\alpha^{m-\mu}-1$ or $\alpha^{m-\mu}+1$ be divisible by $n$,\ednote{Printed $n$ here; $2n+1$ is apparently expected.} now this is impossible, for $m-\mu$ is less than $n$.

Therefore the quantities $1$, $a_{1}$, $a_{2}\ldots a_{n-1}$ are different from one another and consequently they coincide, but in a different order, with the numbers $1$, $2$, $3$, $4\ldots n$.

Therefore, remarking that
\[ \varphi^{2}\{[(2n+1)k_{m}\pm a_{m}]\varepsilon\}=\varphi^{2}(a_{m}\varepsilon), \]
one sees that the quantities (107.) are the same as these:
\[ \Phi^{2}(\varepsilon),\ \Phi^{2}(2\varepsilon)\ \ldots\ \Phi^{2}(n\varepsilon), \]
that is, the roots of the equation (106.) Q.\ E.\ D.

There is still to be remarked, that having
\[ \alpha^{n}=(2n+1)k_{n}-1, \]
one will have
\[ \alpha^{n+m}=(2n+1)k_{n}\alpha^{m}-\alpha^{m}, \]
therefore
\[ a_{n+m}=\pm a_{m} \]
and
\[ \Phi^{2}(\alpha^{n+m}\varepsilon)=\Phi^{2}(\alpha^{m}\varepsilon). \]
This being premised, let $\theta$ be any imaginary root whatever of the equation
\[ \theta^{n}-1=0 \]
and
\[ \psi(\varepsilon)=\Phi^{2}(\varepsilon)+\Phi^{2}(\alpha\varepsilon)\theta+\Phi^{2}(\alpha^{2}\varepsilon)\theta^{2}+\ldots+\Phi^{2}(\alpha^{n-1}\varepsilon)\theta^{n-1}. \tag{108.} \]
By virtue of what we have seen previously, the second member of this equation can be transformed into a \emph{rational} function of $\Phi^{2}(\varepsilon)$. Let us make
\[ \psi(\varepsilon)=\chi[\Phi^{2}(\varepsilon)]. \tag{109.} \]
On putting in the first expression for $\psi(\varepsilon)$, $\alpha^{m}\varepsilon$ in place of $\varepsilon$, there will result:
\[ \begin{aligned}
\psi(\alpha^{m}\varepsilon)&=\varphi^{2}(\alpha^{m}\varepsilon)+\varphi^{2}(\alpha^{m+1}\varepsilon).\theta+\varphi^{2}(\alpha^{m+2}\varepsilon).\theta^{2}+\ldots\\
&\ldots+\Phi^{2}(\alpha^{n-1}\varepsilon).\theta^{n-m-1}+\Phi^{2}(\alpha^{n}\varepsilon).\theta^{n-m}+\ldots+\Phi^{2}(\alpha^{n+m-1}\varepsilon)\theta^{n-1},
\end{aligned} \]
but we have seen that $\Phi^{2}(\alpha^{n+m}\varepsilon)=\Phi^{2}(\alpha^{m}\varepsilon)$; therefore:
\[ \begin{aligned}
\psi(\alpha^{m}\varepsilon)&=\theta^{n-m}.\Phi^{2}(\varepsilon)+\theta^{n-m+1}.\Phi^{2}(\alpha\varepsilon)+\theta^{n-m+2}.\Phi^{2}(\alpha^{2}\varepsilon)+\ldots\\
&\ldots+\theta^{n-1}.\Phi^{2}(\alpha^{m-1}\varepsilon)+\Phi^{2}(\alpha^{m}\varepsilon)+\theta.\Phi^{2}(\alpha^{m+1}\varepsilon)+\ldots+\theta^{n-m-1}.\Phi^{2}(\alpha^{n-1}\varepsilon).
\end{aligned} \]

\origpage{144}
On multiplying by $\theta^{m}$, the second member will become equal to $\psi(\varepsilon)$, therefore:
\[ \psi(\alpha^{m}\varepsilon)=\theta^{-m}.\psi(\varepsilon), \tag{110.} \]
or rather:
\[ \psi(\varepsilon)=\theta^{m}.\chi[\Phi^{2}(\alpha^{m}\varepsilon)], \]
whence, on raising both to the $n^{\text{th}}$ power, one deduces by virtue of the relation $\theta^{+mn}=1$:
\[ [\psi(\varepsilon)]^{n}=\{\chi[\Phi^{2}(\alpha^{m}\varepsilon)]\}^{n}. \tag{111.} \]

This formula gives, on making successively $m=0$, $1$, $2$, $3\ldots\ldots n-1$, $n$ equations, which, added member to member, will give the following:
\[ n[\psi(\varepsilon)]^{n}=\{\chi[\Phi^{2}(\varepsilon)]\}^{n}+\{\chi[\Phi^{2}(\alpha\varepsilon)]\}^{n}+\{\chi[\Phi^{2}(\alpha^{2}\varepsilon)]\}^{n}+\ldots+\{\chi[\Phi^{2}(\alpha^{n-1}\varepsilon)]\}^{n}; \tag{112.} \]
now, the second member of this equation is a rational and symmetric function of the quantities $\Phi^{2}(\varepsilon)$, $\Phi^{2}(\alpha\varepsilon)\ldots\Phi^{2}(\alpha^{n-1}\varepsilon)$, that is, of the roots of the equation (106.); therefore $(\psi\varepsilon)^{n}$ can be expressed as a rational function of $p_{0}$, $p_{1}\ldots p_{n-1}$, consequently as a rational function of any one whatever of these quantities. Let $v$ be the value of $\psi(\varepsilon)$, one will have
\[ \sqrt[n]{v}=\Phi^{2}(\varepsilon)+\theta.\Phi^{2}(\alpha\varepsilon)+\theta^{2}.\Phi^{2}(\alpha^{2}\varepsilon)+\ldots+\theta^{n-1}.\Phi^{2}(\alpha^{n-1}\varepsilon). \tag{113.} \]
This being premised, let $\theta=\cos\frac{2\pi}{n}+i\sin\frac{2\pi}{n}$. The imaginary roots of the equation $\theta^{n}-1$ can be represented by
\[ \theta^{1},\ \theta^{2}\ \ldots\ldots\ \theta^{n-1}. \]
Therefore, on making successively $\theta$ equal to each of these roots and denoting the corresponding values of $v$ by $v_{1}$, $v_{2}\ldots v_{n-1}$, there will result
\[ \begin{aligned}
\sqrt[n]{v_{1}}&=\Phi^{2}(\varepsilon)+\theta.\Phi^{2}(\alpha\varepsilon)+\ldots\ldots+\theta^{n-1}.\Phi^{2}(\alpha^{n-1}\varepsilon),\\
\sqrt[n]{v_{2}}&=\Phi^{2}(\varepsilon)+\theta^{2}.\Phi^{2}(\alpha\varepsilon)+\ldots\ldots+\theta^{2n-2}.\Phi^{2}(\alpha^{n-1}\varepsilon)\\
&\cdots\cdots\cdots\cdots\cdots\cdots\cdots\cdots\cdots\\
\sqrt[n]{v_{n-1}}&=\Phi^{2}(\varepsilon)+\theta^{n-1}.\Phi^{2}(\alpha\varepsilon)+\ldots\ldots+\theta^{(n-1)^{2}}.\Phi^{2}(\alpha^{n-1}\varepsilon),
\end{aligned} \]
On combining these equations with the following:
\[ -p_{n-1}=\Phi^{2}(\varepsilon)+\Phi^{2}(\alpha\varepsilon)+\ldots\ldots+\Phi^{2}(\alpha^{n-1}\varepsilon), \]
one easily deduces from them:
\[ \begin{aligned}
\Phi^{2}(\alpha^{m}\varepsilon)=+\frac{1}{n}&\left\{-p_{n-1}+\theta^{-m}.\sqrt[n]{v_{1}}+\theta^{-2m}.\sqrt[n]{v_{2}}+\theta^{-3m}.\sqrt[n]{v_{3}}+\ldots\right.\\
&\left.\ldots+\theta^{-(n-1)m}\sqrt[n]{v_{n-1}}\right\},
\end{aligned} \tag{114.} \]
and for $m=0$:
\[ \Phi^{2}(\varepsilon)=\frac{1}{n}\left\{-p_{n-1}+\sqrt[n]{v_{1}}+\sqrt[n]{v_{2}}+\ldots+\sqrt[n]{v_{n-1}}\right\}. \tag{115.} \]

\origpage{145}

\subsection*{22.}

All the roots of the equation (106.) are contained in the formula (115.), but since their number is only $n$, it still remains to give to $\Phi^{2}(\varepsilon)$ a form which does not contain roots foreign to the question. Now this is easily done, as follows:

Let
\[ s_{k}=\frac{\sqrt[n]{v_{k}}}{(\sqrt[n]{v_{1}})^{k}}. \]
On putting here $\varepsilon$ in place of $\alpha^{m}\varepsilon$, $\sqrt[n]{v_{k}}$ will be changed into $\theta^{-km}.\sqrt[n]{v_{k}}$, and $v_{1}$ into $\theta^{-m}.v_{1}$, therefore $s_{k}$ will be changed into:
\[ \frac{\theta^{-km}.\sqrt[n]{v_{k}}}{(\theta^{-m}\sqrt[n]{v_{1}})^{k}}=\frac{\sqrt[n]{v_{k}}}{(\sqrt[n]{v_{1}})^{k}}. \]

The function $s_{k}$, as one sees, does not change in value on putting $\alpha^{m}\varepsilon$ in place of $\varepsilon$. Now $s_{k}$ is a rational function of $\Phi^{2}(\varepsilon)$. Therefore, denoting $s_{k}$ by $\lambda[\Phi^{2}(\varepsilon)]$, one will have
\[ s_{k}=\lambda[\Phi^{2}(\alpha^{m}\varepsilon)], \]
whatever the integer $m$ may be. From this one will deduce in the same manner, and as we have found $(\psi\varepsilon)^{n}$, the value of $s_{k}$ as a rational function of one of the quantities $p_{0}$, $p_{1}$, $....p_{n-1}$. Knowing $s_{k}$, one has:
\[ \sqrt[n]{v_{k}}=s_{k}.(\sqrt[n]{v_{1}})^{k}. \]
Therefore, putting $v$ in place of $v_{1}$, the expression for $\Phi^{2}(\alpha^{m}\varepsilon)$ will become:
\[ \Phi^{2}(\alpha^{m}\varepsilon)=\frac{1}{n}\left\{-p_{n-1}+\theta^{-m}.v^{\frac{1}{n}}+s_{2}\theta^{-2m}.v^{\frac{2}{n}}+....+s_{n-1}\theta^{-(n-1)m}.v^{\frac{n-1}{n}}\right\}, \tag{116.} \]
for $m=0$:
\[ \Phi^{2}(\varepsilon)=\frac{1}{n}\left\{-p_{n-1}+v^{\frac{1}{n}}+s_{2}.v^{\frac{2}{n}}+s_{3}.v^{\frac{3}{n}}+....+s_{n-1}.v^{\frac{n-1}{n}}\right\}. \tag{117.} \]
This value has only $n$ different values, which correspond to the $n$ values of $v^{\frac{1}{n}}$. Therefore, in the last place, the solution of the equation $P_{2n+1}=0$ is reduced to that of a single equation of the degree $2n+2$; but this equation does not appear to be algebraically soluble \emph{in general}. Nevertheless it can be completely solved in several particular cases, e.g., when $e=c$, $e=c\sqrt{3}$, $e=c(2\pm\sqrt{3})$ etc. In the course of this memoir I shall occupy myself with these cases, of which the first especially is remarkable, as much by the simplicity of the solution as by its fine application in geometry.

\origpage{146}

Indeed, among other things, I have arrived at this theorem:

“The entire circumference of the lemniscate can be divided into $m$ equal parts by the \emph{ruler and compass alone}, if $m$ is of the form $2^{n}$ or \uncertain{$2^{n}+1$}, the latter number being at the same time prime; or else if $m$ is a product of several numbers of these two forms.”

This theorem is, as one sees, precisely the same as that of Mr.\ \emph{Gauss}, relative to the circle.

\section*{§. VI.}

\emph{Various expressions of the functions} $\Phi(n\beta)$, $f(n\beta)$, $F(n\beta)$.

\subsection*{23.}

On making use of the known formulas which give the values of the coefficients of an algebraic equation as functions of the roots, one can derive several expressions of the functions $\varphi(n\beta)$, $f(n\beta)$, $F(n\beta)$ from the formulas of the preceding section.

I shall consider the most remarkable.

To abridge the formulas, I shall make use of the following notations: I shall denote:

1) By $\overset{k'}{\underset{k}{\Sigma}}_{m}\psi(m)$ the sum, and by $\overset{k'}{\underset{k}{\Pi}}_{m}\psi(m)$ the product of all the quantities of the form $\psi(m)$, which one will obtain by giving to $m$ all the integral values from $k$ to $k'$, the limits $k$ and $k'$ being included.

2) By $\overset{k'}{\underset{k}{\Sigma}}_{m}\overset{\nu'}{\underset{\nu}{\Sigma}}_{\mu}\psi(m,\mu)$ the sum, and by $\overset{k'}{\underset{k}{\Pi}}_{m}\overset{\nu'}{\underset{\nu}{\Pi}}_{\mu}\psi(m,\mu)$ the product of all the quantities of the form $\psi(m,\mu)$, which one will obtain by giving to $m$ all the integral values from $k$ to $k'$, and to $\mu$ the integral values from $\nu$ to $\nu'$, always including the limits therein.

According to this it is clear that one will have:
\[ \overset{k'}{\underset{k}{\Sigma}}_{m}\psi(m)=\psi(k)+\psi(k+1)+.....\psi(k'), \tag{119.} \]
\ednote{No equation number 118. is printed; the numbering passes from 117. to 119.}
\[ \overset{k'}{\underset{k}{\Pi}}_{m}\psi(m)=\psi(k).\psi(k+1)+.....\psi(k'), \tag{120.} \]
\ednote{A plus sign is printed before the row of dots in the product; a multiplication dot is expected.}
\[ \overset{k'}{\underset{k}{\Sigma}}_{m}\overset{\nu'}{\underset{\nu}{\Sigma}}_{\mu}\psi(m,\mu)=\overset{\nu'}{\underset{\nu}{\Sigma}}_{\mu}\psi(k,\mu)+\overset{\nu'}{\underset{\nu}{\Sigma}}_{\mu}\psi(k+1,\mu)+.....+\overset{\nu'}{\underset{\nu}{\Sigma}}_{\mu}\psi(k',\mu), \tag{121.} \]
\[ \overset{k'}{\underset{k}{\Pi}}_{m}\overset{\nu'}{\underset{\nu}{\Pi}}_{\mu}\psi(m,\mu)=\overset{\nu'}{\underset{\nu}{\Pi}}_{\mu}\psi(k,\mu).\overset{\nu'}{\underset{\nu}{\Pi}}_{\mu}\psi(k+1,\mu).........\overset{\nu'}{\underset{\nu}{\Pi}}_{\mu}\psi(k',\mu). \tag{122.} \]

\origpage{147}

This being premised, let us consider the equations
\[ \left\{\begin{aligned}
\Phi(2n+1)\beta&=\frac{P_{2n+1}}{Q_{2n+1}},\\
f(2n+1)\beta&=\frac{P'_{2n+1}}{Q_{2n+1}},\\
F(2n+1)\beta&=\frac{P''_{2n+1}}{Q_{2n+1}}.
\end{aligned}\right. \tag{123.} \]

We have seen that $P_{2n+1}$ is a rational function of $x$ of the degree $(2n+1)^{2}$ and of the form $x.\psi(x^{2})$. Likewise $P'_{2n+1}$ and $F''_{2n+1}$ are functions of this same form, the first of $y$ and the second of $z$.\ednote{Printed $F''_{2n+1}$ where $P''_{2n+1}$ is expected.} Finally $Q_{2n+1}$ is a function which, expressed indifferently in $x$, $y$ or $z$, will be of the degree $(2n+1)^{2}-1$, and will contain only even powers.
Therefore one will have
\[ \begin{aligned}
P_{2n+1}&=A.x^{(2n+1)^{2}}+......+B.x;&\quad P'_{2n+1}&=A'.y^{(2n+1)^{2}}+......+B'.y;\\
P''_{2n+1}&=A''.z^{(2n+1)^{2}}+......+B''.z;&\quad Q_{2n+1}&=C.x^{(2n+1)^{2}-1}+....+D;\\
Q_{2n+1}&=C'.y^{(2n+1)^{2}-1}+....+D';&\quad Q_{2n+1}&=C''.z^{(2n+1)^{2}-1}+....+D''.
\end{aligned} \]

On substituting these values in (123.), there will result
\[ \begin{aligned}
&\{A.x^{(2n+1)^{2}}+.....+B.x\}-\varphi(2n+1)\beta.\{C.x^{(2n+1)^{2}-1}+....+D\},\\
&\{A'.y^{(2n+1)^{2}}+.....+B'.y\}-f(2n+1)\beta.\{C'.y^{(2n+1)^{2}-1}+....+D'\},\\
&\{A''.z^{(2n+1)^{2}}+.....+B''.z\}-F(2n+1)\beta.\{C''.z^{(2n+1)^{2}-1}+....+D''\}.
\end{aligned} \]

In the first of these equations $A$ is the coefficient of the first term, $-\varphi(2n+1)\beta.C$ that of the second and $-\varphi(2n+1)\beta.D$ the last term. Therefore $\frac{C}{A}\varphi(2n+1)\beta$ is equal to the sum and $\pm\frac{D}{A}\varphi(2n+1)\beta$ equal to the product of the roots of the equation in question, which equation is the same as this:
\[ \Phi(2n+1)\beta=\frac{P_{2n+1}}{Q_{2n+1}}. \tag{124.} \]

Therefore, remarking that $A$, $C$ and $D$ (and in general all the coefficients) are independent of $\beta$, one sees that $\Phi(2n+1)\beta$ is (to within a constant coefficient) equal to the sum and to the product of all the roots of the equation (124.).

In the same manner one sees that $f(2n+1)\beta$ and $F(2n+1)\beta$ are respectively equal to the product or to the sum of the roots of the equations
\[ f(2n+1)\beta=\frac{P'_{2n+1}}{Q_{2n+1}},\ F(2n+1)\beta=\frac{P''_{2n+1}}{Q_{2n+1}}, \]

\origpage{148}
taking care to multiply the result by a constant coefficient, suitably chosen.

Now the roots of the equations (123.), according to No.\ 11., are respectively:
\[ \begin{aligned}
x&=(-1)^{m+\mu}.\Phi\left(\beta+\frac{m}{2n+1}\omega+\frac{\mu}{2n+1}\varpi i\right),\\
y&=(-1)^{m}.f\left(\beta+\frac{m}{2n+1}\omega+\frac{\mu}{2n+1}\varpi i\right),\\
z&=(-1)^{\mu}.F\left(\beta+\frac{m}{2n+1}\omega+\frac{\mu}{2n+1}\varpi i\right),
\end{aligned} \]
where the limits of $m$ and $\mu$ are $-n$ and $+n$.

Therefore, by virtue of what has just been seen, and making use of the notations adopted, one will have the following formulas:
\[ \left\{\begin{aligned}
\Phi(2n+1)\beta&=A.\overset{+n}{\underset{-n}{\Sigma}}_{m}\overset{+n}{\underset{-n}{\Sigma}}_{\mu}(-1)^{m+\mu}.\Phi\left(\beta+\frac{m.\omega+\mu\varpi i}{2n+1}\right),\\
f(2n+1)\beta&=A'.\overset{+n}{\underset{-n}{\Sigma}}_{m}\overset{+n}{\underset{-n}{\Sigma}}_{\mu}(-1)^{m}.f\left(\beta+\frac{m\omega+\mu\varpi i}{2n+1}\right),\\
F(2n+1)\beta&=A''.\overset{+n}{\underset{-n}{\Sigma}}_{m}\overset{+n}{\underset{-n}{\Sigma}}_{\mu}(-1)^{\mu}.F\left(\beta+\frac{m\omega+\mu\varpi i}{2n+1}\right);\\
\Phi(2n+1)\beta&=B.\overset{+n}{\underset{-n}{\Pi}}_{m}\overset{+n}{\underset{-n}{\Pi}}_{\mu}.\Phi\left(\beta+\frac{m\omega+\mu\varpi i}{2n+1}\right),\\
f(2n+1)\beta&=B'.\overset{+n}{\underset{-n}{\Pi}}_{m}\overset{+n}{\underset{-n}{\Pi}}_{\mu}.f\left(\beta+\frac{m\omega+\mu\varpi i}{2n+1}\right),\\
F(2n+1)\beta&=B''.\overset{+n}{\underset{-n}{\Pi}}_{m}\overset{+n}{\underset{-n}{\Pi}}_{\mu}.F\left(\beta+\frac{m\omega+\mu\varpi i}{2n+1}\right).
\end{aligned}\right. \tag{125.} \]

To determine the constant quantities $A$, $A'$, $A''$, $B$, $B'$, $B''$, it will be necessary to give $\beta$ a particular value. Thus on making in the first three formulas $\beta=\frac{\omega}{2}+\frac{\varpi}{2}i$, after having divided both members by $\Phi\beta$, there will result, remarking that $\Phi\left(\frac{\omega}{2}+\frac{\varpi}{2}i\right)=\frac{1}{0}$:
\[ \left.\begin{aligned}
A&=\frac{\varphi(2n+1)\beta}{\varphi\beta}\\
A'&=\frac{f(2n+1)\beta}{f\beta}\\
A''&=\frac{F(2n+1)\beta}{f\beta}
\end{aligned}\right\}\ \text{for }\beta=\frac{\omega}{2}+\frac{\varpi}{2}i. \]
\ednote{The denominator of the third fraction is printed $f\beta$; $F\beta$ is apparently expected.}

\origpage{149}

Let $\beta=\frac{\omega}{2}+\frac{\varpi}{2}i+\alpha$, one has:
\[ \left.\begin{aligned}
A&=\frac{\varphi\left((2n+1)\alpha+n\omega+n\varpi i+\frac{\omega}{2}+\frac{\varpi}{2}i\right)}{\varphi\left(\alpha+\frac{\omega}{2}+\frac{\varpi}{2}i\right)}\\
&=\frac{\varphi\left((2n+1)\alpha+\frac{\omega}{2}+\frac{\varpi}{2}i\right)}{\varphi\left(\alpha+\frac{\omega}{2}+\frac{\varpi}{2}i\right)}=\frac{\varphi\alpha}{\varphi(2n+1)\alpha},\\
A'&=\frac{f\left((2n+1)\alpha+n\omega+n\varpi i+\frac{\omega}{2}+\frac{\varpi}{2}i\right)}{f\left(\alpha+\frac{\omega}{2}+\frac{\varpi}{2}i\right)}\\
&=(-1)^{n}.\frac{f\left((2n+1)\alpha+\frac{\omega}{2}+\frac{\varpi}{2}i\right)}{f\left(\alpha+\frac{\omega}{2}+\frac{\varpi}{2}i\right)}=(-1)^{n}.\frac{f\left(\alpha+\frac{\omega}{2}\right)}{f\left((2n+1)\alpha+\frac{\omega}{2}\right)},\\
A''&=\frac{F\left((2n+1)\alpha+n\omega+n\varpi i+\frac{\omega}{2}+\frac{\varpi}{2}i\right)}{F\left(\alpha+\frac{\omega}{2}+\frac{\varpi}{2}i\right)}\\
&=(-1)^{n}.\frac{F\left((2n+1)\alpha+\frac{\omega}{2}+\frac{\varpi}{2}i\right)}{F\left(\alpha+\frac{\omega}{2}+\frac{\varpi}{2}i\right)}=(-1)^{n}.\frac{F\left(\alpha+\frac{\varpi}{2}i\right)}{F\left((2n+1)\alpha+\frac{\varpi}{2}i\right)}.
\end{aligned}\right\}\ \text{for }\alpha=0. \]

These expressions for $A$, $A'$, $A''$ will become of the form $\frac{0}{0}$, on making $\alpha=0$, therefore one will find according to the known rules:
\[ A=\frac{1}{2n+1},\quad A'=A''=\frac{(-1)^{n}}{2n+1}. \]
According to this the first three formulas will become:
\[ \left\{\begin{aligned}
\Phi(2n+1)\beta&=\frac{1}{2n+1}.\overset{+n}{\underset{-n}{\Sigma}}_{m}\overset{+n}{\underset{-n}{\Sigma}}_{\mu}(-1)^{m+\mu}.\varphi\left(\beta+\frac{m\omega+\mu\varpi i}{2n+1}\right),\\
f(2n+1)\beta&=\frac{(-1)^{n}}{2n+1}.\overset{+n}{\underset{-n}{\Sigma}}_{m}\overset{+n}{\underset{-n}{\Sigma}}_{\mu}(-1)^{m}.f\left(\beta+\frac{m\omega+\mu\varpi i}{2n+1}\right),\\
F(2n+1)\beta&=\frac{(-1)^{n}}{2n+1}.\overset{+n}{\underset{-n}{\Sigma}}_{m}\overset{+n}{\underset{-n}{\Sigma}}_{\mu}(-1)^{\mu}.F\left(\beta+\frac{m\omega+\mu\varpi i}{2n+1}\right).
\end{aligned}\right. \tag{126.} \]

\origpage{150}

To obtain the value of the constants $B$, $B'$, $B''$, I remark that there will be:
\[ \begin{aligned}
&\overset{+n}{\underset{-n}{\Pi}}_{m}\overset{+n}{\underset{-n}{\Pi}}_{\mu}\psi(m,\mu)=\psi(0,0).\overset{n}{\underset{1}{\Pi}}_{m}\psi(m,0).\overset{n}{\underset{1}{\Pi}}_{\mu}\psi(0,\mu)\\
&\quad\times\overset{n}{\underset{1}{\Pi}}_{m}\overset{n}{\underset{1}{\Pi}}_{\mu}\psi(m,\mu).\psi(-m,-\mu).\overset{n}{\underset{1}{\Pi}}_{m}\overset{n}{\underset{1}{\Pi}}_{\mu}\psi(m,-\mu).\psi(-m,\mu).
\end{aligned} \tag{127.} \]

Applying this transformation to the formulas (126.), dividing the first by $\Phi\beta$, the second by $f\beta$ and the third by $F\beta$, then making in the first $\beta=0$, in the second $\beta=\frac{\omega}{2}$ and in the third $\beta=\frac{\varpi}{2}i$, and remarking that $\frac{\varphi(2n+1)\beta}{\varphi\beta}=2n+1$, for $\beta=0$, that $\frac{f(2n+1)\beta}{f\beta}=2n+1$, for $\beta=\frac{\omega}{2}$, and that $\frac{F(2n+1)\beta}{F\beta}=2n+1$, for $\beta=\frac{\varpi}{2}i$: we shall find:
\[ \left\{\begin{aligned}
(2n+1)&=B.\overset{n}{\underset{1}{\Pi}}_{m}\varphi^{2}\left(\frac{m\omega}{2n+1}\right).\overset{n}{\underset{1}{\Pi}}_{\mu}\varphi^{2}\left(\frac{\mu\varpi i}{2n+1}\right)\\
&\quad\times\overset{n}{\underset{1}{\Pi}}_{m}\overset{n}{\underset{1}{\Pi}}_{\mu}\varphi^{2}\left(\frac{m\omega+\mu\varpi i}{2n+1}\right).\varphi^{2}\left(\frac{m\omega-\mu\varpi i}{2n+1}\right),\\
(2n+1)&=B'.\overset{n}{\underset{1}{\Pi}}_{m}f^{2}\left(\frac{\omega}{2}+\frac{m\omega}{2n+1}\right).\overset{n}{\underset{1}{\Pi}}_{\mu}f^{2}\left(\frac{\omega}{2}+\frac{\mu\varpi i}{2n+1}\right)\\
&\quad\times\overset{n}{\underset{1}{\Pi}}_{m}\overset{n}{\underset{1}{\Pi}}_{\mu}f^{2}\left(\frac{\omega}{2}+\frac{m\omega+\mu\varpi i}{2n+1}\right).f^{2}\left(\frac{\omega}{2}+\frac{m\omega-\mu\varpi i}{2n+1}\right),\\
(2n+1)&=B''.\overset{n}{\underset{1}{\Pi}}_{m}F^{2}\left(\frac{\varpi}{2}i+\frac{m\omega}{2n+1}\right).\overset{n}{\underset{1}{\Pi}}_{\mu}F^{2}\left(\frac{\varpi}{2}i+\frac{\mu\varpi i}{2n+1}\right)\\
&\quad\times\overset{n}{\underset{1}{\Pi}}_{m}\overset{n}{\underset{1}{\Pi}}_{\mu}F^{2}\left(\frac{\varpi}{2}i+\frac{m\omega+\mu\varpi i}{2n+1}\right).F^{2}\left(\frac{\varpi}{2}i+\frac{m\omega-\mu\varpi i}{2n+1}\right).
\end{aligned}\right. \tag{128.} \]

Drawing from these equations the values of $B$, $B'$, $B''$, and then substituting them in the transformed formulas, there will result:

\origpage{151}
\[ \left\{\begin{aligned}
&\Phi(2n+1)\beta=\\
&(2n+1)\Phi\beta.\overset{n}{\underset{1}{\Pi}}_{m}\frac{\varphi\left(\beta+\frac{m\omega}{2n+1}\right).\varphi\left(\beta-\frac{m\omega}{2n+1}\right)}{\varphi^{2}\left(\frac{m\omega}{2n+1}\right)}.\overset{n}{\underset{1}{\Pi}}_{\mu}\frac{\varphi\left(\beta+\frac{\mu\varpi i}{2n+1}\right).\varphi\left(\beta-\frac{\mu\varpi i}{2n+1}\right)}{\varphi^{2}\left(\frac{\mu\varpi i}{2n+1}\right)}\\
&\quad\times\overset{n}{\underset{1}{\Pi}}_{m}\overset{n}{\underset{1}{\Pi}}_{\mu}\frac{\varphi\left(\beta+\frac{m\omega+\mu\varpi i}{2n+1}\right).\varphi\left(\beta-\frac{m\omega+\mu\varpi i}{2n+1}\right)}{\varphi^{2}\left(\frac{m\omega+\mu\varpi i}{2n+1}\right)}.\frac{\varphi\left(\beta+\frac{m\omega-\mu\varpi i}{2n+1}\right).\varphi\left(\beta-\frac{m\omega-\mu\varpi i}{2n+1}\right)}{\varphi^{2}\left(\frac{m\omega-\mu\varpi i}{2n+1}\right)},\\
&f(2n+1)\beta=\\
&(2n+1)f\beta.\overset{n}{\underset{1}{\Pi}}_{m}\frac{f\left(\beta+\frac{m\omega}{2n+1}\right).f\left(\beta-\frac{m\omega}{2n+1}\right)}{f^{2}\left(\frac{\omega}{2}+\frac{m\omega}{2n+1}\right)}.\overset{n}{\underset{1}{\Pi}}_{\mu}\frac{f\left(\beta+\frac{\mu\varpi i}{2n+1}\right).f\left(\beta-\frac{\mu\varpi i}{2n+1}\right)}{f^{2}\left(\frac{\omega}{2}+\frac{\mu\varpi i}{2n+1}\right)}\\
&\quad\times\overset{n}{\underset{1}{\Pi}}_{m}\overset{n}{\underset{1}{\Pi}}_{\mu}\frac{f\left(\beta+\frac{m\omega+\mu\varpi i}{2n+1}\right).f\left(\beta-\frac{m\omega+\mu\varpi i}{2n+1}\right)}{f^{2}\left(\frac{\omega}{2}+\frac{m\omega+\mu\varpi i}{2n+1}\right)}.\frac{f\left(\beta+\frac{m\omega-\mu\varpi i}{2n+1}\right).f\left(\beta-\frac{m\omega-\mu\varpi i}{2n+1}\right)}{f^{2}\left(\frac{\omega}{2}+\frac{m\omega-\mu\varpi i}{2n+1}\right)},\\
&F(2n+1)\beta=\\
&(2n+1)F\beta.\overset{n}{\underset{1}{\Pi}}_{m}\frac{F\left(\beta+\frac{m\omega}{2n+1}\right).F\left(\beta-\frac{m\omega}{2n+1}\right)}{F^{2}\left(\frac{\varpi}{2}i+\frac{m\omega}{2n+1}\right)}.\overset{n}{\underset{1}{\Pi}}_{\mu}\frac{F\left(\beta+\frac{\mu\varpi i}{2n+1}\right).F\left(\beta-\frac{\mu\varpi i}{2n+1}\right)}{F^{2}\left(\frac{\varpi}{2}i+\frac{\mu\varpi i}{2n+1}\right)}\\
&\quad\times\overset{n}{\underset{1}{\Pi}}_{m}\overset{n}{\underset{1}{\Pi}}_{\mu}\frac{F\left(\beta+\frac{m\omega+\mu\varpi i}{2n+1}\right).F\left(\beta-\frac{m\omega+\mu\varpi i}{2n+1}\right)}{F^{2}\left(\frac{\varpi}{2}i+\frac{m\omega+\mu\varpi i}{2n+1}\right)}.\frac{F\left(\beta+\frac{m\omega-\mu\varpi i}{2n+1}\right).F\left(\beta-\frac{m\omega-\mu\varpi i}{2n+1}\right)}{F^{2}\left(\frac{\varpi}{2}i+\frac{m\omega-\mu\varpi i}{2n+1}\right)}.
\end{aligned}\right. \tag{129.} \]

These formulas may be given simpler forms, by making use of the following formulas:
\[ \frac{\varphi(\beta+\alpha).\varphi(\beta-\alpha)}{\varphi^{2}\alpha}=-\frac{1-\frac{\varphi^{2}\beta}{\varphi^{2}\alpha}}{1-\frac{\varphi^{2}\beta}{\varphi^{2}\left(\alpha+\frac{\omega}{2}+\frac{\varpi}{2}i\right)}}, \]
\[ \frac{f(\beta+\alpha).f(\beta-\alpha)}{f^{2}\left(\frac{\omega}{2}+\alpha\right)}=-\frac{1-\frac{f^{2}\beta}{f^{2}\left(\frac{\omega}{2}+\alpha\right)}}{1-\frac{f^{2}\beta}{f^{2}\left(\alpha+\frac{\omega}{2}+\frac{\varpi}{2}i\right)}}, \]
\[ \frac{F^{2}(\beta+\alpha).F(\beta-\alpha)}{F^{2}\left(\frac{\varpi}{2}i+\alpha\right)}=-\frac{1-\frac{F^{2}\beta}{F^{2}\left(\frac{\varpi}{2}i+\alpha\right)}}{1-\frac{F^{2}\beta}{F^{2}\left(\frac{\omega}{2}+\frac{\varpi}{2}i+\alpha\right)}}, \]
\ednote{The numerator of the third left-hand fraction is printed $F^{2}(\beta+\alpha).F(\beta-\alpha)$; $F(\beta+\alpha).F(\beta-\alpha)$ is apparently expected.}

\origpage{152}
which one will easily verify by means of the formulas (13.), (16.), (18.).

By virtue of these formulas it is clear that the equations (129.) may be put under the form:
\[ \left\{\begin{aligned}
&\Phi(2n+1)\beta=\\
&(2n+1)\Phi\beta.\overset{n}{\underset{1}{\Pi}}_{m}\frac{1-\frac{\varphi^{2}\beta}{\varphi^{2}\left(\frac{m\omega}{2n+1}\right)}}{1-\frac{\varphi^{2}\beta}{\varphi^{2}\left(\frac{\omega}{2}+\frac{\varpi}{2}i+\frac{m\omega}{2n+1}\right)}}.\overset{n}{\underset{1}{\Pi}}_{\mu}\frac{1-\frac{\varphi^{2}\beta}{\varphi^{2}\left(\frac{\mu\varpi i}{2n+1}\right)}}{1-\frac{\varphi^{2}\beta}{\varphi^{2}\left(\frac{\omega}{2}+\frac{\varpi}{2}i+\frac{\mu\varpi i}{2n+1}\right)}}\\
&\quad\times\overset{n}{\underset{1}{\Pi}}_{m}\overset{n}{\underset{1}{\Pi}}_{\mu}\frac{1-\frac{\varphi^{2}\beta}{\varphi^{2}\left(\frac{m\omega+\mu\varpi i}{2n+1}\right)}}{1-\frac{\varphi^{2}\beta}{\varphi^{2}\left(\frac{\omega}{2}+\frac{\varpi}{2}i+\frac{m\omega+\mu\varpi i}{2n+1}\right)}}.\frac{1-\frac{\varphi^{2}\beta}{\varphi^{2}\left(\frac{m\omega-\mu\varpi i}{2n+1}\right)}}{1-\frac{\varphi^{2}\beta}{\varphi^{2}\left(\frac{\omega}{2}+\frac{\varpi}{2}i+\frac{m\omega-\mu\varpi i}{2n+1}\right)}},\\
&f(2n+1)\beta=\\
&(2n+1)f\beta.\overset{n}{\underset{1}{\Pi}}_{m}\frac{1-\frac{f^{2}\beta}{f^{2}\left(\frac{m\omega}{2n+1}+\frac{\omega}{2}\right)}}{1-\frac{f^{2}\beta}{f^{2}\left(\frac{\omega}{2}+\frac{\varpi}{2}i+\frac{m\omega}{2n+1}\right)}}.\overset{n}{\underset{1}{\Pi}}_{\mu}\frac{1-\frac{f^{2}\beta}{f^{2}\left(\frac{\omega}{2}+\frac{\mu\varpi i}{2n+1}\right)}}{1-\frac{f^{2}\beta}{f^{2}\left(\frac{\omega}{2}+\frac{\varpi}{2}i+\frac{\mu\varpi i}{2n+1}\right)}}\\
&\quad\times\overset{n}{\underset{1}{\Pi}}_{m}\overset{n}{\underset{1}{\Pi}}_{\mu}\frac{1-\frac{f^{2}\beta}{f^{2}\left(\frac{\omega}{2}+\frac{m\omega+\mu\varpi i}{2n+1}\right)}}{1-\frac{f^{2}\beta}{f^{2}\left(\frac{\omega}{2}+\frac{\varpi}{2}i+\frac{m\omega+\mu\varpi i}{2n+1}\right)}}.\frac{1-\frac{f^{2}\beta}{f^{2}\left(\frac{\omega}{2}+\frac{m\omega-\mu\varpi i}{2n+1}\right)}}{1-\frac{f^{2}\beta}{f^{2}\left(\frac{\omega}{2}+\frac{\varpi}{2}i+\frac{m\omega-\mu\varpi i}{2n+1}\right)}},\\
&F(2n+1)\beta=\\
&(2n+1)F\beta.\overset{n}{\underset{1}{\Pi}}_{m}\frac{1-\frac{F^{2}\beta}{F^{2}\left(\frac{\varpi}{2}i+\frac{m\omega}{2n+1}\right)}}{1-\frac{F^{2}\beta}{F^{2}\left(\frac{\omega}{2}+\frac{\varpi}{2}i+\frac{m\omega}{2n+1}\right)}}.\overset{n}{\underset{1}{\Pi}}_{\mu}\frac{1-\frac{F^{2}\beta}{F^{2}\left(\frac{\varpi}{2}i+\frac{\mu\varpi i}{2n+1}\right)}}{1-\frac{F^{2}\beta}{F^{2}\left(\frac{\omega}{2}+\frac{\varpi}{2}i+\frac{\mu\varpi i}{2n+1}\right)}}\\
&\quad\times\overset{n}{\underset{1}{\Pi}}_{m}\overset{n}{\underset{1}{\Pi}}_{\mu}\frac{1-\frac{F^{2}\beta}{F^{2}\left(\frac{\varpi}{2}i+\frac{m\omega+\mu\varpi i}{2n+1}\right)}}{1-\frac{F^{2}\beta}{F^{2}\left(\frac{\omega}{2}+\frac{\varpi}{2}i+\frac{m\omega+\mu\varpi i}{2n+1}\right)}}.\frac{1-\frac{F^{2}\beta}{F^{2}\left(\frac{\varpi}{2}i+\frac{m\omega-\mu\varpi i}{2n+1}\right)}}{1-\frac{F^{2}\beta}{F^{2}\left(\frac{\omega}{2}+\frac{\varpi}{2}i+\frac{m\omega-\mu\varpi i}{2n+1}\right)}}.
\end{aligned}\right. \tag{130.} \]

\origpage{153}

These formulas give, as one sees, the values of $\Phi(2n+1)\beta$, $f(2n+1)\beta$ and $F(2n+1)\beta$, expressed respectively as rational functions of $\Phi\beta$, $f\beta$ and $F\beta$, in the form of products.

We shall further give the values of $f(2n+1)\beta$, $F(2n+1)\beta$ under another form, which will be useful in what follows.

We have $f^{2}\beta=1-c^{2}\varphi^{2}\beta$, hence
\[ 1-\frac{f^{2}\beta}{f^{2}\alpha}=\frac{c^{2}(\varphi^{2}\beta-\varphi^{2}\alpha)}{f^{2}\alpha}=\frac{c^{2}}{f^{2}\alpha}\Phi^{2}\beta-\frac{c^{2}\varphi^{2}\alpha}{f^{2}\alpha} \]
and
\[ 1-\frac{f^{2}\beta}{f^{2}\left(\frac{\varpi}{2}i+\alpha\right)}=\frac{c^{2}\left[\varphi^{2}\beta-\varphi^{2}\left(\frac{\varpi}{2}i+\alpha\right)\right]}{f^{2}\left(\frac{\varpi}{2}i+\alpha\right)}, \]
now by virtue of (18.) we have:
\[ f^{2}\left(\frac{\varpi}{2}i+\alpha\right)=\frac{e^{2}+c^{2}}{e^{2}}.\frac{1}{f^{2}\alpha}; \]
hence:
\[ \frac{1-\frac{f^{2}\beta}{f^{2}\alpha}}{1-\frac{f^{2}\beta}{f^{2}\left(\frac{\varpi}{2}i+\alpha\right)}}=\frac{1}{f^{4}\alpha}.\frac{e^{2}+c^{2}}{e^{2}}.\frac{1-\frac{\varphi^{2}\beta}{\varphi^{2}\alpha}}{1-\frac{\varphi^{2}\beta}{\varphi^{2}\left(\frac{\varpi}{2}i+\alpha\right)}}. \]
We shall find in the same way:
\[ \frac{1-\frac{F^{2}\beta}{F^{2}\alpha}}{1-\frac{F^{2}\beta}{F^{2}\left(\frac{\omega}{2}+\alpha\right)}}=\frac{1}{F^{4}\alpha}.\frac{e^{2}+c^{2}}{e^{2}}.\frac{1-\frac{\varphi^{2}\beta}{\varphi^{2}\alpha}}{1-\frac{\varphi^{2}\beta}{\varphi^{2}\left(\frac{\omega}{2}+\alpha\right)}}. \]

By virtue of these formulas, and making $\beta=0$ in order to determine the constant factor, it is clear that the expressions of $f(2n+1)\beta$, $F(2n+1)\beta$ may be written as follows:

\origpage{154}
\[ \left\{\begin{aligned}
&f(2n+1)\beta=f\beta.\overset{n}{\underset{1}{\Pi}}_{m}\frac{1-\frac{\varphi^{2}\beta}{\varphi^{2}\left(\frac{\omega}{2}+\frac{m\omega}{2n+1}\right)}}{1-\frac{\varphi^{2}\beta}{\varphi^{2}\left(\frac{\omega}{2}+\frac{\varpi}{2}i+\frac{m\omega}{2n+1}\right)}}.\overset{n}{\underset{1}{\Pi}}_{\mu}\frac{1-\frac{\varphi^{2}\beta}{\varphi^{2}\left(\frac{\omega}{2}i+\frac{\mu\varpi i}{2n+1}\right)}}{1-\frac{\varphi^{2}\beta}{\varphi^{2}\left(\frac{\omega}{2}+\frac{\varpi}{2}i+\frac{\mu\varpi i}{2n+1}\right)}}\\
&\quad\times\overset{n}{\underset{1}{\Pi}}_{m}\overset{n}{\underset{1}{\Pi}}_{\mu}\frac{1-\frac{\varphi^{2}\beta}{\varphi^{2}\left(\frac{\omega}{2}+\frac{m\omega+\mu\varpi i}{2n+1}\right)}}{1-\frac{\varphi^{2}\beta}{\varphi^{2}\left(\frac{\omega}{2}+\frac{\varpi}{2}i+\frac{m\omega+\mu\varpi i}{2n+1}\right)}}.\frac{1-\frac{\varphi^{2}\beta}{\varphi^{2}\left(\frac{\omega}{2}+\frac{m\omega-\mu\varpi i}{2n+1}\right)}}{1-\frac{\varphi^{2}\beta}{\varphi^{2}\left(\frac{\omega}{2}+\frac{\varpi}{2}i+\frac{m\omega-\mu\varpi i}{2n+1}\right)}},\\
&F(2n+1)\beta=F\beta.\overset{n}{\underset{1}{\Pi}}_{m}\frac{1-\frac{\varphi^{2}\beta}{\varphi^{2}\left(\frac{\varpi}{2}i+\frac{m\omega}{2n+1}\right)}}{1-\frac{\varphi^{2}\beta}{\varphi^{2}\left(\frac{\omega}{2}+\frac{\varpi}{2}i+\frac{m\omega}{2n+1}\right)}}.\overset{n}{\underset{1}{\Pi}}_{\mu}\frac{1-\frac{\varphi^{2}\beta}{\varphi^{2}\left(\frac{\varpi}{2}i+\frac{\mu\varpi i}{2n+1}\right)}}{1-\frac{\varphi^{2}\beta}{\varphi^{2}\left(\frac{\omega}{2}+\frac{\varpi}{2}i+\frac{\mu\varpi i}{2n+1}\right)}}\\
&\quad\times\overset{n}{\underset{1}{\Pi}}_{m}\overset{n}{\underset{1}{\Pi}}_{\mu}\frac{1-\frac{\varphi^{2}\beta}{\varphi^{2}\left(\frac{\varpi}{2}i+\frac{m\omega+\mu\varpi i}{2n+1}\right)}}{1-\frac{\varphi^{2}\beta}{\varphi^{2}\left(\frac{\omega}{2}+\frac{\varpi}{2}i+\frac{m\omega+\mu\varpi i}{2n+1}\right)}}.\frac{1-\frac{\varphi^{2}\beta}{\varphi^{2}\left(\frac{\varpi}{2}i+\frac{m\omega-\mu\varpi i}{2n+1}\right)}}{1-\frac{\varphi^{2}\beta}{\varphi^{2}\left(\frac{\omega}{2}+\frac{\varpi}{2}i+\frac{m\omega-\mu\varpi i}{2n+1}\right)}}.
\end{aligned}\right. \tag{130$'$.} \]
\ednote{In the numerator of the second product of the first line, the print sets an $i$ after $\omega/2$; there is apparently no $i$ there, by analogy with the first product. Reproduced as printed.}

In this paragraph we have considered the functions $\Phi(n\beta)$, $f(n\beta)$, $F(n\beta)$ only in the case of odd values of $n$. One might find analogous expressions of these functions for even values of $n$; but as there is no difficulty in this, and as moreover the formulas to which we have arrived are those which will be most useful to us in what follows, I shall not occupy myself with it.

\section*{§. VII.}

\emph{Development of the functions} $\Phi\alpha$, $f\alpha$, $F\alpha$ \emph{in series and in infinite products.}

\subsection*{24.}

Making in the formulas of the preceding paragraph $\beta=\frac{\alpha}{2n+1}$, we shall obtain expressions of the functions $\Phi\alpha$, $f\alpha$, $F\alpha$ which, on account of the indeterminate number $n$, may be varied in an infinity of ways.

\origpage{155}
Among all the formulas that one will obtain in this way, those which result from the supposition of $n$ infinite are the most remarkable. Then the functions $\Phi$, $f$, $F$ will disappear from the values of $\Phi\alpha$, $f\alpha$, $F\alpha$, and one will obtain for these functions algebraic expressions, but composed of an infinity of terms. To have these expressions, one must make in the formulas (126.), (130.) $\beta=\frac{\alpha}{2n+1}$, and then seek the limit of the second member of these equations for ever increasing values of $n$. For brevity, let $v$ be a quantity whose limit is zero for ever increasing values of $n$. This being premised, let us consider successively the three formulas (126.).

Making in the first of the formulas (126.) $\beta=\frac{\alpha}{2n+1}$, and remarking that
\[ \begin{aligned}
&\overset{+n}{\underset{-n}{\Sigma}}_{m}\overset{+n}{\underset{-n}{\Sigma}}_{\mu}\theta(m,\mu)=\theta(0,0)+\overset{n}{\underset{1}{\Sigma}}_{m}\theta(m,0)+\overset{n}{\underset{1}{\Sigma}}_{\mu}\theta(0,\mu)\\
&\quad+\overset{n}{\underset{1}{\Sigma}}_{m}\overset{n}{\underset{1}{\Sigma}}_{\mu}\{\theta(m,\mu)+\theta(-m,-\mu)+\theta(m,-\mu)+\theta(-m,\mu)\},
\end{aligned} \tag{131.} \]
it is clear that one may put the formula in question under the form:
\[ \begin{aligned}
\Phi\alpha={}&\frac{1}{2n+1}.\varphi\left(\frac{\alpha}{2n+1}\right)+\frac{1}{2n+1}\overset{n}{\underset{1}{\Sigma}}_{m}(-1)^{m}\left\{\Phi\left(\frac{\alpha+m\omega}{2n+1}\right)+\Phi\left(\frac{\alpha-m\omega}{2n+1}\right)\right\}\\
&+\frac{1}{2n+1}\overset{n}{\underset{1}{\Sigma}}_{\mu}(-1)^{\mu}\left\{\Phi\left(\frac{\alpha+\mu\varpi i}{2n+1}\right)+\Phi\left(\frac{\alpha-\mu\varpi i}{2n+1}\right)\right\}-\frac{i}{ec}.\overset{n}{\underset{1}{\Sigma}}_{m}\overset{n}{\underset{1}{\Sigma}}_{\mu}(-1)^{m+\mu}.\psi(n-m,n-\mu)\\
&+\frac{i}{ec}\overset{n}{\underset{1}{\Sigma}}_{m}\overset{n}{\underset{1}{\Sigma}}_{\mu}(-1)^{m+\mu}.\psi_{1}(n-m,n-\mu),
\end{aligned} \tag{132.} \]
where one has made, for brevity,
\[ \left\{\begin{aligned}
\psi(m,\mu)&=\frac{1}{2n+1}.\left\{\frac{1}{\varphi\left(\frac{\alpha+(m+\tfrac{1}{2})\omega+(\mu+\tfrac{1}{2})\varpi i}{2n+1}\right)}+\frac{1}{\varphi\left(\frac{\alpha-(m+\tfrac{1}{2})\omega-(\mu+\tfrac{1}{2})\varpi i}{2n+1}\right)}\right\},\\
\psi_{1}(m,\mu)&=\frac{1}{2n+1}.\left\{\frac{1}{\varphi\left(\frac{\alpha+(m+\tfrac{1}{2})\omega-(\mu+\tfrac{1}{2})\varpi i}{2n+1}\right)}+\frac{1}{\varphi\left(\frac{\alpha-(m+\tfrac{1}{2})\omega+(\mu+\tfrac{1}{2})\varpi i}{2n+1}\right)}\right\}.
\end{aligned}\right. \tag{133.} \]

Now, remarking that
\[ \Phi\left(\frac{\alpha+m\omega}{2n+1}\right)+\Phi\left(\frac{\alpha-m\omega}{2n+1}\right)=\frac{2\varphi\left(\frac{\alpha}{2n+1}\right).f\left(\frac{m\omega}{2n+1}\right).F\left(\frac{m\omega}{2n+1}\right)}{1+e^{2}c^{2}.\varphi^{2}\left(\frac{m\omega}{2n+1}\right).\varphi^{2}\left(\frac{\alpha}{2n+1}\right)}=\frac{A_{m}}{2n+1}, \]
\[ \Phi\left(\frac{\alpha+\mu\varpi i}{2n+1}\right)+\Phi\left(\frac{\alpha-\mu\varpi i}{2n+1}\right)=\frac{2\varphi\left(\frac{\alpha}{2n+1}\right).f\left(\frac{\mu\varpi i}{2n+1}\right).F\left(\frac{\mu\varpi i}{2n+1}\right)}{1+e^{2}c^{2}.\varphi^{2}\left(\frac{\mu\varpi i}{2n+1}\right).\varphi^{2}\left(\frac{\alpha}{2n+1}\right)}=\frac{B_{\mu}}{2n+1}, \]

\origpage{156}
where $A_{m}$ and $B_{\mu}$ are finite quantities, the part of the equation (132.) as far as the member which has the sign $-$ will take the form:
\[ \frac{1}{2n+1}.\varphi\left(\frac{\alpha}{2n+1}\right)+\frac{1}{(2n+1)^{2}}.\overset{n}{\underset{1}{\Sigma}}_{m}(-1)^{m}(A_{m}+B_{\mu}); \]
now the limit of this quantity is evidently zero; hence, taking the limit of the formula (132.), we shall have:
\[ \begin{aligned}
\Phi\alpha&=-\frac{i}{ec}\ \text{lim.}\ \overset{n}{\underset{1}{\Sigma}}_{m}\overset{n}{\underset{1}{\Sigma}}_{\mu}(-1)^{m+\mu}.\psi(n-m,n-\mu)\\
&\quad+\frac{i}{ec}\ \text{lim.}\ \overset{n}{\underset{1}{\Sigma}}_{m}\overset{n}{\underset{1}{\Sigma}}_{\mu}(-1)^{m+\mu}.\psi_{1}(n-m,n-\mu),
\end{aligned} \]
or else:
\[ \begin{aligned}
\Phi\alpha&=-\frac{i}{ec}\ \text{lim.}\ \overset{n-1}{\underset{0}{\Sigma}}_{m}\overset{n-1}{\underset{0}{\Sigma}}_{\mu}(-1)^{m+\mu}.\psi(m,\mu)\\
&\quad+\frac{i}{ec}\ \text{lim.}\ \overset{n-1}{\underset{0}{\Sigma}}_{m}\overset{n-1}{\underset{0}{\Sigma}}_{\mu}(-1)^{m+\mu}.\psi_{1}(m,\mu).
\end{aligned} \tag{134.} \]

It suffices to know one of these limits, for we shall have the other by merely changing the sign of $i$.

Let us seek the limit of
\[ \overset{n-1}{\underset{0}{\Sigma}}_{m}\overset{n-1}{\underset{0}{\Sigma}}_{\mu}(-1)^{m+\mu}.\psi(m,\mu). \]
For this, one must try to put the preceding quantity under the form
\[ P+v, \]
where $P$ is independent of $n$, and $v$ is a quantity which has zero for its limit; for then the quantity $P$ will be precisely the limit in question.

\subsection*{25.}

Let us consider first the expression:
\[ \overset{n-1}{\underset{0}{\Sigma}}_{\mu}(-1)^{\mu}.\psi(m,\mu). \]

Let
\[ \theta(m,\mu)=\frac{2\alpha}{\alpha^{2}-(m\omega+\mu\varpi i)^{2}}, \tag{135.} \]
and let us make
\[ \psi(m,\mu)-\theta(m,\mu)=\frac{2\alpha}{(2n+1)^{2}}.R_{\mu}, \tag{136.} \]
we shall have:
\[ \overset{n-1}{\underset{0}{\Sigma}}_{\mu}(-1)^{\mu}.\psi(m,\mu)-\overset{n-1}{\underset{0}{\Sigma}}_{\mu}(-1)^{\mu}.\theta(m,\mu)=2\alpha.\overset{n-1}{\underset{0}{\Sigma}}_{\mu}(-1)^{n}.\frac{R_{\mu}}{(2n+1)^{2}}. \tag{137.} \]
\ednote{The exponent of the factor $-1$ in the last sum is apparently printed $n$; $\mu$ is expected, as in the two sums on the left.}

This being premised, I say that the second member of this equation is a quantity of the form $\frac{v}{2n+1}$.

\origpage{157}

According to (12.), (13.) we shall have
\[ \frac{1}{\varphi(\beta+\varepsilon)}+\frac{1}{\varphi(\beta-\varepsilon)}=\frac{\varphi(\beta+\varepsilon)+\varphi(\beta-\varepsilon)}{\varphi(\beta+\varepsilon).\varphi(\beta-\varepsilon)}=\frac{2\varphi\beta.f\varepsilon.F\varepsilon}{\varphi^{2}\beta-\varphi^{2}\varepsilon}, \]
hence, making $\beta=\frac{\alpha}{2n+1}$ and $\varepsilon=\frac{m\omega+\mu\varpi i}{2n+1}=\frac{\varepsilon_{\mu}}{2n+1}$ and $f\varepsilon.F\varepsilon=\theta\varepsilon$, we have:
\[ \psi(m,\mu)=\frac{1}{2n+1}.\frac{2\varphi\left(\frac{\alpha}{2n+1}\right).\theta\left(\frac{\varepsilon_{\mu}}{2n+1}\right)}{\varphi^{2}\left(\frac{\alpha}{2n+1}\right)-\varphi^{2}\left(\frac{\varepsilon_{\mu}}{2n+1}\right)}. \]
Now we have:
\[ \varphi\left(\frac{\alpha}{2n+1}\right)=\frac{\alpha}{2n+1}+\frac{A\alpha^{3}}{(2n+1)^{3}}, \]
hence:
\[ \psi(m,\mu)=\frac{\theta\left(\frac{\varepsilon_{\mu}}{2n+1}\right)}{\varphi^{2}\left(\frac{\alpha}{2n+1}\right)-\varphi^{2}\left(\frac{\varepsilon_{\mu}}{2n+1}\right)}.\left\{\frac{2A\alpha^{3}}{(2n+1)^{4}}+\frac{2\alpha}{(2n+1)^{2}}\right\}, \]
and consequently:
\[ \begin{aligned}
\psi(m,\mu)-\theta(m,\mu)&=\frac{2\alpha}{(2n+1)^{2}}\left\{\frac{\theta\left(\frac{\varepsilon_{\mu}}{2n+1}\right)}{\varphi^{2}\left(\frac{\alpha}{2n+1}\right)-\varphi^{2}\left(\frac{\varepsilon_{\mu}}{2n+1}\right)}-\frac{1}{\left(\frac{\alpha}{2n+1}\right)^{2}-\left(\frac{\varepsilon_{\mu}}{2n+1}\right)^{2}}\right\}\\
&\quad+\frac{2A\alpha^{3}}{(2n+1)^{4}}.\frac{\theta\left(\frac{\varepsilon_{\mu}}{2n+1}\right)}{\varphi^{2}\left(\frac{\alpha}{2n+1}\right)-\varphi^{2}\left(\frac{\varepsilon_{\mu}}{2n+1}\right)}.
\end{aligned} \]
Hence the value of $R_{\mu}$ will become
\[ R_{\mu}=\frac{\theta\left(\frac{\varepsilon_{\mu}}{2n+1}\right)}{\varphi^{2}\left(\frac{\alpha}{2n+1}\right)-\varphi^{2}\left(\frac{\varepsilon_{\mu}}{2n+1}\right)}.\left\{1+\frac{A\alpha^{2}}{(2n+1)^{2}}\right\}-\frac{1}{\left(\frac{\alpha}{2n+1}\right)^{2}-\left(\frac{\varepsilon_{\mu}}{2n+1}\right)^{2}}. \tag{138.} \]

This being premised, there are two cases to consider, namely whether $\frac{\varepsilon_{\mu}}{2n+1}$ has zero for its limit or not.

$a)$ If $\frac{\varepsilon_{\mu}}{2n+1}$ has zero for its limit, we shall have:
\[ \begin{aligned}
\varphi^{2}\left(\frac{\varepsilon_{\mu}}{2n+1}\right)&=\frac{\varepsilon_{\mu}^{2}}{(2n+1)^{2}}+\frac{B_{\mu}.\varepsilon_{\mu}^{4}}{(2n+1)^{4}},\\
\theta\left(\frac{\varepsilon_{\mu}}{2n+1}\right)&=\sqrt{\left[1-c^{2}\varphi^{2}\left(\frac{\varepsilon_{\mu}}{2n+1}\right)\right]}.\sqrt{\left[1+e^{2}\varphi^{2}\left(\frac{\varepsilon_{\mu}}{2n+1}\right)\right]}=1+\frac{C_{\mu}.\varepsilon_{\mu}^{2}}{(2n+1)^{2}},\\
\varphi^{2}\left(\frac{\alpha}{2n+1}\right)&=\frac{\alpha^{2}}{(2n+1)^{2}}+\frac{D.\alpha^{4}}{(2n+1)^{4}},
\end{aligned} \]
where $B_{\mu}$, $C_{\mu}$, $D$ have finite limits; hence, substituting:

\origpage{158}
\[ \begin{aligned}
R_{\mu}&=A\alpha^{2}.\frac{\frac{1}{\varepsilon_{\mu}^{2}}+\frac{C_{\mu}}{(2n+1)^{2}}}{\frac{\alpha^{2}}{\varepsilon_{\mu}^{2}}-1+\frac{D\alpha^{4}}{(2n+1)^{2}.\varepsilon_{\mu}^{2}}-B_{\mu}\frac{\varepsilon_{\mu}^{2}}{(2n+1)^{2}}}\\
&\quad+\frac{C_{\mu}.\frac{\alpha^{2}}{\varepsilon_{\mu}^{2}}-\frac{D\alpha^{4}}{\varepsilon_{\mu}^{4}}-C_{\mu}+B_{\mu}}{\left(1-\frac{\alpha^{2}}{\varepsilon_{\mu}^{2}}\right)^{2}-\left(1-\frac{\alpha^{2}}{\varepsilon_{\mu}^{2}}\right)\left\{\frac{D\alpha^{4}}{(2n+1)^{2}.\varepsilon_{\mu}^{2}}-B_{\mu}.\frac{\varepsilon_{\mu}^{2}}{(2n+1)^{2}}\right\}},
\end{aligned} \tag{139.} \]
now whether $\varepsilon_{\mu}$ be finite or infinite, it is clear that this quantity will always converge towards a finite quantity for ever increasing values of $n$. Hence we shall have
\[ R_{\mu}=r_{\mu}+v_{\mu}, \tag{140.} \]
where $r_{\mu}$ is a finite quantity independent of $n$.

$b)$ If $\frac{\varepsilon_{\mu}}{2n+1}$ has for its limit a finite quantity, it is clear that, naming this limit $\delta_{\mu}$, we shall have:
\[ R_{\mu}=-\frac{\theta(\delta_{\mu})}{\varphi^{2}(\delta_{\mu})}+\frac{1}{\delta_{\mu}^{2}}+v_{\mu}^{1}. \tag{141.} \]
This being premised, let us consider the expression $\overset{n}{\underset{1}{\Sigma}}_{\mu}(-1)^{\mu}.\frac{R_{\mu}}{(2n+1)^{2}}$. We have
\[ \begin{aligned}
\overset{n}{\underset{1}{\Sigma}}_{\mu}(-1)^{\mu}.\frac{R_{\mu}}{(2n+1)^{2}}&=\frac{1}{(2n+1)^{2}}.\{R_{0}-R_{1}+R_{2}-R_{3}+....\\
&\quad....+(-1)^{\nu-1}.R_{\nu-1}+(-1)^{\nu+1}[R_{\nu+1}-R_{\nu+1}+R_{\nu+2}-R_{\nu+3}...+(-1)^{n-\nu-1}.R_{n-1}]\}.
\end{aligned} \tag{142.} \]
\ednote{The first term inside the square bracket is printed $R_{\nu+1}$; $R_{\nu}$ is apparently expected.}
Let us suppose first that $\frac{\varepsilon_{\mu}}{2n+1}$ has for its limit a finite quantity, whatever be the value of $\mu$. Then, remarking that
\[ \delta_{\mu+1}=\delta_{\mu}, \]
we shall have
\[ R_{\mu}-R_{\mu+1}=v_{\mu}^{1}-v_{\mu+1}^{1}, \]
hence:
\[ \overset{n-1}{\underset{0}{\Sigma}}_{\mu}(-1)^{\mu}.\frac{R_{\mu}}{(2n+1)^{2}}=\frac{1}{(2n+1)^{2}}(v_{0}^{1}-v_{1}^{1}+v_{2}^{1}-v_{3}^{1}+....+v_{k-2}^{1}-v_{k-1}^{1})+\frac{B}{(2n+1)^{2}}, \]
where $k=n$ or $n-1$, according as $n$ is even or odd. The quantity $B$ always has for its limit a finite quantity, namely $B=0$ if $n$ is even, and $B=R_{n-1}$ if $n$ is odd.

Now it is known that a sum such as
\[ v_{0}^{1}-v_{1}^{1}+v_{2}^{1}-....+v_{k-2}^{1}-v_{k-1}^{1}, \]
may be put under the form $k.v$, $v$ having zero for its limit. Hence, substituting:

\origpage{159}
\[ \overset{n-1}{\underset{0}{\Sigma}}_{\mu}(-1)^{\mu}.\frac{R_{\mu}}{(2n+1)^{2}}=\frac{k.v+B}{(2n+1)^{2}}; \]
now, $k$ being equal to $n$ or to $n-1$, and $B$ finite, the limit of $\frac{k.v+B}{2n+1}$ will be zero, hence:
\[ \overset{n-1}{\underset{0}{\Sigma}}_{\mu}(-1)^{\mu}.\frac{R_{\mu}}{(2n+1)^{2}}=\frac{v}{2n+1}. \tag{143.} \]

Let us now suppose that $\frac{m}{2n+1}$ has zero for its limit. Then $\frac{\varepsilon_{\mu}}{2n+1}$ likewise has zero for its limit, unless at the same time $\frac{\mu}{2n+1}$ has for its limit a finite quantity. Let $\nu$ in this case be the whole number immediately below $\sqrt{n}$, and let us consider the sum
\[ R_{0}-R_{1}+R_{2}-R_{3}+....+(-1)^{\nu-1}R_{\nu-1}. \]
Supposing that $\mu$ is one of the numbers $0$, $1$, $....\nu$, it is clear that $\frac{\varepsilon_{\mu}}{(2n+1)}=\frac{m\omega+\mu\varpi i}{2n+1}$ has zero for its limit; hence, according to what we have seen, $R_{\mu}$ will be a finite quantity, and consequently
\[ R_{0}-R_{1}+R_{2}-....+(-1)^{\nu-1}R_{\nu-1}=\nu.R, \]
where $R$ is likewise a finite quantity.

Let us now consider the sum
\[ (-1)^{\nu}\{R_{\nu}-R_{\nu+1}+R_{\nu+2}-....+(-1)^{n-\nu-1}.R_{n-1}\}. \]

If $\frac{\varepsilon_{\mu}}{2n+1}$ has for its limit a quantity different from zero, we have, as will be seen:
\[ R_{\mu}-R_{\mu+1}=v_{\mu}^{1}-v_{\mu+1}^{1}; \]
if on the contrary $\frac{\varepsilon_{\mu}}{2n+1}$ has zero for its limit, we have:
\[ R_{\mu}=v_{\mu}+v_{\mu}^{1}; \]
now, if at the same time $\mu>\sqrt{n}$, it is clear that, by virtue of the value of $R_{\mu}$,
\[ v_{\mu}=B_{\mu}-C_{\mu}; \]
now it is clear that $B_{\mu}$ and $C_{\mu}$ both have for limits quantities independent of $\mu$, hence, naming these limits $B$ and $C$, we shall have:
\[ R_{\mu}=B-C+v_{\mu}, \]
and consequently, in this case also,
\[ R_{\mu}-R_{\mu+1}=v_{\mu}-v_{\mu+1}. \]
\ednote{The last subscript is printed indistinctly, apparently $v_{\mu+1}$ with the $\mu$ damaged in print.}

Hence, as in the case where $\frac{\varepsilon_{\mu}}{2n+1}$ would have a limit different from zero for all values of $\mu$, we shall demonstrate that
\[ \frac{(-1)^{\nu}}{(2n+1)^{2}}\{R_{\nu}-R_{\nu+1}+....+(-1)^{n-\nu-1}R_{n-1}\}=\frac{v}{(2n+1)}. \]

\origpage{160}
Now, combining the equations above, we shall derive from them
\[ \overset{n-1}{\underset{0}{\Sigma}}_{\mu}(-1)^{\mu}.\frac{R_{\mu}}{(2n+1)^{2}}=\frac{1}{(2n+1)^{2}}.(\nu R)+\frac{v}{2n+1}; \]
but $\frac{v}{2n+1}$ has zero for its limit, therefore
\[ \overset{n-1}{\underset{0}{\Sigma}}_{\mu}(-1)^{\mu}.\frac{R_{\mu}}{(2n+1)^{2}}=\frac{v}{2n+1}. \]
Therefore this formula always holds, and consequently the formula (137.) will become:
\[ \overset{n-1}{\underset{0}{\Sigma}}_{\mu}(-1)^{\mu}\psi(m,\mu)-\overset{n-1}{\underset{0}{\Sigma}}_{\mu}(-1)^{\mu}.\theta(m,\mu)=\frac{v}{2n+1}. \tag{144.} \]
This being premised, the point is to put $\overset{n-1}{\underset{0}{\Sigma}}_{\mu}(-1)^{\mu}\theta(m,\mu)$ under the form $P+\frac{v}{2n+1}$.
Now this can be done as follows. We have:
\[ \left\{\begin{aligned}
&\overset{n-1}{\underset{0}{\Sigma}}_{\mu}(-1)^{\mu}\theta(m,\mu)=\overset{\frac{1}{0}}{\underset{0}{\Sigma}}_{\mu}(-1)^{\mu}\theta(m,\mu)-\overset{\frac{1}{0}}{\underset{n}{\Sigma}}_{\mu}(-1)^{\mu}\theta(m,\mu)\ \text{and}\\
&\overset{\frac{1}{0}}{\underset{n}{\Sigma}}_{\mu}(-1)^{\mu}.\theta(m,\mu)=(-1)^{n}\{\theta(m,n)-\theta(m,n+1)+\theta(m,n+2)-....\text{etc.}\}.
\end{aligned}\right. \tag{145.} \]
Now, by a known formula, we have:
\[ \begin{aligned}
&\theta(m,n)-\theta(m,n+1)+\theta(m,n+2)-......\\
&=\tfrac{1}{2}\theta(m,n)+A\frac{\partial\theta(m,n)}{\partial n}+B.\frac{\partial^{3}\theta(m,n)}{\partial^{3}n}+......,
\end{aligned} \]
where $A$, $B$$....$ are numbers; now
\[ \theta(m,n)=\frac{2\alpha}{\alpha^{2}-(m\omega+n\varpi i)^{2}}, \]
therefore, on substituting:
\[ \theta(m,n)-\theta(m,n+1)+....=\frac{\alpha}{\alpha^{2}-(m\omega+n\varpi i)^{2}}+\frac{4A\alpha\varpi i.n}{\{\alpha^{2}-(m\omega+n\varpi i)^{2}\}^{2}}+.... \]
\ednote{The exponent outside the braces in the denominator of the second term is printed indistinctly; $2$ is read here.}
From this it follows that
\[ \theta(m,n)-\theta(m,n+1)+....=\frac{\alpha}{\varpi^{2}.n^{2}}+\frac{v}{n^{2}}=\frac{v}{2n+1}. \]
Therefore, by virtue of the equations (145.)
\[ \overset{n-1}{\underset{0}{\Sigma}}_{\mu}(-1)^{\mu}\theta(m,\mu)=\overset{\frac{1}{0}}{\underset{0}{\Sigma}}_{\mu}(-1)^{\mu}.\theta(m,\mu)+\frac{v}{2n+1}, \]
and consequently
\[ \overset{n-1}{\underset{0}{\Sigma}}_{\mu}(-1)^{\mu}\psi(m,\mu)=\overset{\infty}{\underset{0}{\Sigma}}_{\mu}(-1)^{\mu}.\theta(m,\mu)+\frac{v}{2n+1}. \tag{146.} \]

\subsection*{26.}

Having transformed in this way the quantity $\overset{n-1}{\underset{0}{\Sigma}}_{\mu}(-1)^{\mu}.\psi(m,\mu)$, we derive from the equation (146.):

\origpage{161}
\[ \overset{n-1}{\underset{0}{\Sigma}}_{m}\overset{n-1}{\underset{0}{\Sigma}}_{\mu}(-1)^{m+\mu}\psi(m,\mu)=\overset{n-1}{\underset{0}{\Sigma}}(-1)^{m}.\varrho_{m}+\overset{n-1}{\underset{0}{\Sigma}}.\frac{v_{m}}{2n+1}, \tag{147.} \]
on putting
\[ \varrho_{m}=\overset{\infty}{\underset{0}{\Sigma}}_{\mu}(-1)^{\mu}.\theta(m,\mu); \tag{148.} \]
now
\[ \overset{n-1}{\underset{0}{\Sigma}}.\frac{v_{m}}{2n+1}=\frac{v_{0}+v_{1}+v_{2}\ldots v_{n-1}}{2n+1}=\frac{nv}{2n+1}=\frac{v}{2}, \]
$v$ having zero for its limit. Therefore the equation (147.) will give, on making $n$ infinite:
\[ \text{lim.}\overset{n-1}{\underset{0}{\Sigma}}_{m}\overset{n-1}{\underset{0}{\Sigma}}_{\mu}(-1)^{m+\mu}.\psi(m,\mu)=\overset{\infty}{\underset{0}{\Sigma}}(-1)^{m}.\varrho_{m}. \tag{149.} \]
In the same way, if one puts, for brevity:
\[ \left\{\begin{aligned}\theta_{1}(m,\mu)&=\frac{2\alpha}{\alpha^{2}-(m\omega-\mu\varpi i)^{2}},\\ \varrho'_{m}&=\overset{\infty}{\underset{0}{\Sigma}}_{\mu}(-1)^{\mu}.\theta_{1}(m,\mu),\end{aligned}\right. \tag{150.} \]
we shall have:
\[ \text{lim.}\overset{n-1}{\underset{0}{\Sigma}}_{m}\overset{n-1}{\underset{0}{\Sigma}}_{\mu}(-1)^{m+\mu}.\psi_{1}(m,\mu)=\overset{\infty}{\underset{0}{\Sigma}}(-1)^{m}.\varrho'_{m}. \tag{151.} \]
Having found these two quantities, of which the expression of $\Phi\alpha$ is composed, we shall have, on substituting:
\[ \Phi\alpha=-\frac{i}{ec}.\overset{\infty}{\underset{0}{\Sigma}}(-1)^{m}.\varrho_{m}+\frac{i}{ec}.\overset{\infty}{\underset{0}{\Sigma}}(-1)^{m}.\varrho'_{m}=\frac{i}{ec}.\overset{\infty}{\underset{0}{\Sigma}}(-1)^{m}\{\varrho'_{m}-\varrho_{m}\}, \]
or else, on restoring the values of $\varrho'_{m}$ and $\varrho_{m}$,
\[ \Phi\alpha=\frac{i}{ec}.\overset{\infty}{\underset{0}{\Sigma}}_{m}(-1)^{m}\left\{\overset{\infty}{\underset{0}{\Sigma}}_{\mu}\left\{\frac{2\alpha}{\alpha^{2}-[(m+\tfrac{1}{2})\omega-(\mu+\tfrac{1}{2})\varpi i]^{2}}-\frac{2\alpha}{\alpha^{2}-[(m+\tfrac{1}{2})\omega+(\mu+\tfrac{1}{2})\varpi i]^{2}}\right\}\right\}. \tag{152.} \]
Now
\[ \frac{2\alpha}{\alpha^{2}-[(m+\tfrac{1}{2})\omega\pm(\mu+\tfrac{1}{2})\varpi i]^{2}}=\frac{1}{\alpha-(m+\tfrac{1}{2})\omega\mp(\mu+\tfrac{1}{2})\varpi i}+\frac{1}{\alpha+(m+\tfrac{1}{2})\omega\pm(\mu+\tfrac{1}{2})\varpi i}, \]
therefore:
\[ \begin{aligned}&\frac{2\alpha}{\alpha^{2}-[(m+\tfrac{1}{2})\omega-(\mu+\tfrac{1}{2})\varpi i]^{2}}-\frac{2\alpha}{\alpha^{2}-[(m+\tfrac{1}{2})\omega+(\mu+\tfrac{1}{2})\varpi i]^{2}}\\ &=+\frac{(2\mu+1)\varpi i}{[\alpha+(m+\tfrac{1}{2})\omega]^{2}+(\mu+\tfrac{1}{2})^{2}\varpi^{2}}-\frac{(2\mu+1)\varpi i}{[\alpha-(m+\tfrac{1}{2})\omega]^{2}+(\mu+\tfrac{1}{2})^{2}\varpi^{2}},\end{aligned} \]
therefore the expression of $\Phi\alpha$ will become, under a real form:
\[ \Phi\alpha=\frac{1}{ec}.\overset{\infty}{\underset{0}{\Sigma}}_{m}(-1)^{m}.\overset{\infty}{\underset{0}{\Sigma}}_{\mu}(-1)^{\mu}.\left\{\frac{(2\mu+1)\varpi}{[\alpha-(m+\tfrac{1}{2})\omega]^{2}+(\mu+\tfrac{1}{2})^{2}\varpi^{2}}-\frac{(2\mu+1)\varpi}{[\alpha+(m+\tfrac{1}{2})\omega]^{2}+(\mu+\tfrac{1}{2})^{2}\varpi^{2}}\right\}, \tag{153.} \]
that is, we shall have:

\origpage{162}
\[ \begin{aligned}\Phi\alpha&=\frac{\varpi}{ec}(\delta_{0}-\delta_{1}+\delta_{2}-\delta_{3}+....+(-1)^{m}\delta_{m}-....)\\ &-\frac{\varpi}{ec}(\delta'_{0}-\delta'_{1}+\delta'_{2}-\delta'_{3}+....+(-1)^{m}\delta'_{m}-....),\end{aligned} \tag{154.} \]
or else:
\[ \left\{\begin{aligned}\delta_{m}&=\frac{1}{[\alpha-(m+\tfrac{1}{2})\omega]^{2}+\frac{\varpi^{2}}{4}}-\frac{3}{[\alpha-(m+\tfrac{1}{2})\omega]^{2}+\frac{9\varpi^{2}}{4}}+\frac{5}{[\alpha-(m+\tfrac{1}{2})\omega]^{2}+\frac{25\varpi^{2}}{4}}-\ldots\\ \delta'_{\mu}&=\frac{1}{[\alpha+(m+\tfrac{1}{2})\omega]^{2}+\frac{\varpi^{2}}{4}}-\frac{3}{[\alpha+(m+\tfrac{1}{2})\omega]^{2}+\frac{9\varpi^{2}}{4}}+\frac{5}{[\alpha+(m+\tfrac{1}{2})\omega]^{2}+\frac{25\varpi^{2}}{4}}-\ldots\end{aligned}\right. \tag{155.} \]

If one begins the search for the limit of the function $\overset{n-1}{\underset{0}{\Sigma}}_{m}\overset{n-1}{\underset{0}{\Sigma}}_{\mu}(-1)^{m+\mu}.\psi(m,\mu)$ with that of $\overset{n-1}{\underset{0}{\Sigma}}(-1)^{m}\psi(m,\mu)$ instead of that of $\overset{n-1}{\underset{0}{\Sigma}}(-1)^{\mu}\psi(m,\mu)$, as we have done, one will find in place of the formula (53.)\ednote{Printed (53.); the formula meant is presumably (153.), an apparent misprint.} the following:
\[ \Phi\alpha=\frac{1}{ec}.\overset{\infty}{\underset{0}{\Sigma}}_{\mu}(-1)^{\mu}.\overset{\infty}{\underset{0}{\Sigma}}_{m}(-1)^{m}.\left\{\frac{(2\mu+1)\varpi}{[\alpha-(m+\tfrac{1}{2})\omega]^{2}+(\mu+\tfrac{1}{2})^{2}\varpi^{2}}-\frac{(2\mu+1)\varpi}{[\alpha+(m+\tfrac{1}{2})\omega]^{2}+(\mu+\tfrac{1}{2})^{2}\varpi^{2}}\right\}, \tag{156.} \]
that is:
\[ \begin{aligned}\Phi\alpha&=\frac{\varpi}{ec}(\varepsilon_{0}-3\varepsilon_{1}+5\varepsilon_{2}-7\varepsilon_{3}+....+(-1)^{\mu}.(2\mu+1)\varepsilon_{\mu}-....)\\ &-\frac{\varpi}{ec}(\varepsilon'_{0}-3\varepsilon'_{1}+5\varepsilon'_{2}-7\varepsilon'_{3}+....+(-1)^{\mu}.(2\mu+1)\varepsilon_{\mu}-....),\end{aligned} \tag{157.} \]
\ednote{In the second row the last term is printed with $\varepsilon_{\mu}$ without the prime; $\varepsilon'_{\mu}$ is expected, an apparent misprint.}
where
\[ \left\{\begin{aligned}\varepsilon_{\mu}&=\frac{1}{\left(\alpha-\frac{\omega}{2}\right)^{2}+(\mu+\tfrac{1}{2})^{2}\varpi^{2}}-\frac{1}{\left(\alpha-\frac{3\omega}{2}\right)^{2}+(\mu+\tfrac{1}{2})^{2}\varpi^{2}}+\frac{1}{\left(\alpha-\frac{5\omega}{2}\right)^{2}+(\mu+\tfrac{1}{2})^{2}\varpi^{2}}-....\\ \varepsilon'_{\mu}&=\frac{1}{\left(\alpha+\frac{\omega}{2}\right)^{2}+(\mu+\tfrac{1}{2})^{2}\varpi^{2}}-\frac{1}{\left(\alpha+\frac{3\omega}{2}\right)^{2}+(\mu+\tfrac{1}{2})^{2}\varpi^{2}}+\frac{1}{\left(\alpha+\frac{5\omega}{2}\right)^{2}+(\mu+\tfrac{1}{2})^{2}\varpi^{2}}-....\end{aligned}\right. \tag{158.} \]

\subsection*{27.}

Let us now seek the expression of $f\alpha$, by means of the second of the formulas (126.). On making therein $\beta=\frac{\alpha}{2n+1}$, and remarking that the limit of the quantities contained in the first two lines then becomes equal to zero, we shall have:
\[ \begin{aligned}f\alpha&=\text{lim.}(-1)^{n}\overset{n}{\underset{1}{\Sigma}}_{m}\overset{n}{\underset{1}{\Sigma}}_{\mu}(-1)^{m}.\psi(n-m,n-\mu)\\ &+\text{lim.}(-1)^{n}\overset{n}{\underset{1}{\Sigma}}_{m}\overset{n}{\underset{1}{\Sigma}}_{\mu}(-1)^{m}.\psi_{1}(n-m,n-\mu),\end{aligned} \tag{159.} \]

\origpage{163}
where we have put, for brevity:
\[ \begin{aligned}\psi(n-m,n-\mu)&=\frac{1}{2n+1}\left\{f\left(\frac{\alpha+m\omega+\mu\varpi i}{2n+1}\right)+f\left(\frac{\alpha-m\omega-\mu\varpi i}{2n+1}\right)\right\},\\ \psi_{1}(n-m,n-\mu)&=\frac{1}{2n+1}\left\{f\left(\frac{\alpha+m\omega-\mu\varpi i}{2n+1}\right)+f\left(\frac{\alpha-m\omega+\mu\varpi i}{2n+1}\right)\right\}.\end{aligned} \]
Now we have:
\[ f(\beta-\varepsilon)+f(\beta-\varepsilon)=\frac{2f\beta.f\varepsilon}{1+e^{2}c^{2}\varphi^{2}\varepsilon.\varphi^{2}\beta}=\frac{f\varepsilon}{e^{2}c^{2}\varphi^{2}\varepsilon}.\frac{f\beta}{\varphi^{2}\beta+\frac{1}{e^{2}c^{2}\varphi^{2}\varepsilon}}. \]
\ednote{Printed $f(\beta-\varepsilon)$ in both terms; the first term is presumably $f(\beta+\varepsilon)$, an apparent misprint.}

Let
\[ \varepsilon=\frac{m\omega+\mu\varpi i}{2n+1}, \]
we shall have:
\[ \begin{aligned}\frac{1}{\varphi\varepsilon}&=-iec.\Phi\left(\frac{\omega}{2}+\frac{\varpi}{2}i-\varepsilon\right)=-iec.\Phi\left(\frac{(n-m+\tfrac{1}{2})\omega+(n-\mu+\tfrac{1}{2})\varpi i}{2n+1}\right),\\ \frac{f\varepsilon}{\varphi\varepsilon}&=-\frac{i\sqrt{(e^{2}+c^{2})}}{F\left(\varepsilon-\frac{\varpi}{2}i\right)}=-i\sqrt{(e^{2}+c^{2})}.\frac{c}{\sqrt{(e^{2}+c^{2})}}.F\left(\varepsilon-\frac{\omega}{2}-\frac{\varpi}{2}i\right)\\ &=-ci.F.\left(\frac{(n-m+\tfrac{1}{2})\omega+(n-\mu+\tfrac{1}{2})\varpi i}{2n+1}\right).\end{aligned} \]
Therefore we shall have, on substituting and putting $m$ and $\mu$ respectively in place of $n-m$ and $n-\mu$:
\[ \psi(m,\mu)=\frac{1}{e}.\frac{2\varphi\left(\frac{(m+\tfrac{1}{2})\omega+(\mu+\tfrac{1}{2})\varpi i}{2n+1}\right).F\left(\frac{(m+\tfrac{1}{2})\omega+(\mu+\tfrac{1}{2})\varpi i}{2n+1}\right)}{(2n+1)\left[\varphi^{2}\left(\frac{\alpha}{2n+1}\right)-\varphi^{2}\left(\frac{(m+\tfrac{1}{2})\omega+(\mu+\tfrac{1}{2})\varpi i}{2n+1}\right)\right]}. \]
We shall have the value of $\psi_{1}(m,\mu)$ by changing only the sign of $i$. Making now
\[ \theta(m,\mu)=\frac{(2m+1)\omega+(2\mu+1)\varpi i}{\alpha^{2}-[(m+\tfrac{1}{2})\omega+(\mu+\tfrac{1}{2})\varpi i]^{2}} \]
and
\[ \theta_{1}(m,\mu)=\frac{(2m+1)\omega-(2\mu+1)\varpi i}{\alpha^{2}-[(m+\tfrac{1}{2})\omega-(\mu+\tfrac{1}{2})\varpi i]^{2}}, \]
and then seeking the limit of the function
\[ \overset{n-1}{\underset{0}{\Sigma}}_{m}\overset{n-1}{\underset{0}{\Sigma}}_{\mu}(-1)^{m}.\psi(m,\mu), \]
in the same manner as before, we shall find:
\[ \text{lim.}\overset{n-1}{\underset{0}{\Sigma}}_{m}\overset{n-1}{\underset{0}{\Sigma}}_{\mu}(-1)^{m}.\psi(m,\mu)=\frac{1}{e}.\overset{\infty}{\underset{0}{\Sigma}}_{\mu}\left\{\overset{\infty}{\underset{0}{\Sigma}}_{m}(-1)^{m}.\theta(m,\mu)\right\} \]
and
\[ \text{lim.}\overset{n-1}{\underset{0}{\Sigma}}_{m}\overset{n-1}{\underset{0}{\Sigma}}_{\mu}(-1)^{m}.\psi_{1}(m,\mu)=\frac{1}{e}.\overset{\infty}{\underset{0}{\Sigma}}_{\mu}\left\{\overset{\infty}{\underset{0}{\Sigma}}_{m}(-1)^{m}.\theta_{1}(m,\mu)\right\}; \]

\origpage{164}
therefore, on substituting in (159.), and restoring the values of $\theta(m,\mu)$ and $\theta_{1}(m,\mu)$, we have:
\[ f\alpha=\frac{1}{e}.\overset{\infty}{\underset{0}{\Sigma}}_{\mu}\overset{\infty}{\underset{0}{\Sigma}}_{m}(-1)^{m}\left\{\frac{(2m+1)\omega+(2\mu+1)\varpi i}{\alpha^{2}-[(m+\tfrac{1}{2})\omega+(\mu+\tfrac{1}{2})\varpi i]^{2}}+\frac{(2m+1)\omega-(2\mu+1)\varpi i}{\alpha^{2}-[(m+\tfrac{1}{2})\omega-(\mu+\tfrac{1}{2})\varpi i]^{2}}\right\}. \tag{160.} \]

The quantity enclosed between the brackets may also be put under the form:
\[ \frac{2[\alpha-(m+\tfrac{1}{2})\omega]}{[\alpha-(m+\tfrac{1}{2})\omega]^{2}+(\mu+\tfrac{1}{2})^{2}\varpi^{2}}-\frac{2[\alpha+(m+\tfrac{1}{2})\omega]}{[\alpha+(m+\tfrac{1}{2})\omega]^{2}+(\mu+\tfrac{1}{2})^{2}\varpi^{2}}, \]
therefore also:
\[ f\alpha=\frac{1}{e}.\overset{\infty}{\underset{0}{\Sigma}}_{\mu}\left\{\overset{\infty}{\underset{0}{\Sigma}}_{m}(-1)^{m}.\frac{2[\alpha-(m+\tfrac{1}{2})\omega]}{[\alpha-(m+\tfrac{1}{2})\omega]^{2}+(\mu+\tfrac{1}{2})^{2}\varpi^{2}}-\overset{\infty}{\underset{0}{\Sigma}}_{m}(-1)^{m}\frac{2[\alpha+(m+\tfrac{1}{2})\omega]}{[\alpha+(m+\tfrac{1}{2})\omega]^{2}+(\mu+\tfrac{1}{2})^{2}\varpi^{2}}\right\}. \tag{161.} \]

We shall have in the same manner:
\[ F(\alpha)=\frac{1}{c}.\overset{\infty}{\underset{0}{\Sigma}}_{m}\left\{\overset{\infty}{\underset{0}{\Sigma}}_{\mu}(-1)^{\mu}.\frac{(2\mu+1)\varpi}{[\alpha-(m+\tfrac{1}{2})\omega]^{2}+(\mu+\tfrac{1}{2})\varpi^{2}}+\overset{\infty}{\underset{0}{\Sigma}}(-1)^{\mu}.\frac{(2\mu+1)\varpi}{[\alpha+(m+\tfrac{1}{2})\omega]^{2}+(\mu+\tfrac{1}{2})\varpi^{2}}\right\}. \tag{162.} \]
\ednote{Both denominators are printed with $(\mu+\tfrac{1}{2})\varpi^{2}$; the square on the bracket, $(\mu+\tfrac{1}{2})^{2}\varpi^{2}$, is expected as in (161.), an apparent misprint.}

\subsection*{28.}

Let us now come to the formulas (130.). To find the value of the second member, after having made $\beta=\frac{\alpha}{2n+1}$, and supposed $n$ infinite, we shall first seek the limit of the following expression:
\[ t=\overset{n}{\underset{1}{\Pi}}_{m}\overset{n}{\underset{1}{\Pi}}_{\mu}\left\{\frac{1-\frac{\varphi^{2}\left(\frac{\alpha}{2n+1}\right)}{\varphi^{2}\left(\frac{m\omega+\mu\varpi i+k}{2n+1}\right)}}{1-\frac{\varphi^{2}\left(\frac{\alpha}{2n+1}\right)}{\varphi^{2}\left(\frac{m\omega+\mu\varpi i+l}{2n+1}\right)}}\right\}, \tag{163.} \]
where $k$ and $l$ are two quantities independent of $n$, $m$, $\mu$.

On taking the logarithm, and putting, for brevity:
\[ \psi(m,\mu)=\log\left\{\frac{1-\frac{\varphi^{2}\left(\frac{\alpha}{2n+1}\right)}{\varphi^{2}\left(\frac{m\omega+\mu\varpi i+k}{2n+1}\right)}}{1-\frac{\varphi^{2}\left(\frac{\alpha}{2n+1}\right)}{\varphi^{2}\left(\frac{m\omega+\mu\varpi i+l}{2n+1}\right)}}\right\}, \tag{164.} \]

\origpage{165}
we shall have:
\[ \log t=\overset{n}{\underset{1}{\Sigma}}_{m}\overset{n}{\underset{1}{\Sigma}}_{\mu}.\psi(m,\mu). \tag{165.} \]

Let us first consider the expression: $\overset{n}{\underset{1}{\Sigma}}_{\mu}\psi(m,\mu)$. Let
\[ \theta(m,\mu)=\log\left\{\frac{1-\frac{\alpha^{2}}{(m\omega+\mu\varpi i+k)^{2}}}{1-\frac{\alpha^{2}}{(m\omega+\mu\varpi i+l)^{2}}}\right\}, \tag{166.} \]
we shall have:
\[ \psi(m,\mu)-\theta(m,\mu)=\log\left\{\frac{1-\frac{\varphi^{2}\left(\frac{\alpha}{2n+1}\right)}{\varphi^{2}\left(\frac{m\omega+\mu\varpi i+k}{2n+1}\right)}}{1-\frac{\alpha^{2}}{(m\omega+\mu\varpi i+k)^{2}}}.\frac{1-\frac{\alpha^{2}}{(m\omega+\mu\varpi i+l)^{2}}}{1-\frac{\varphi^{2}\left(\frac{\alpha}{2n+1}\right)}{\varphi^{2}\left(\frac{m\omega+\mu\varpi i+l}{2n+1}\right)}}\right\}. \]
This being premised, I say that the second member of this equation is, for every value of $m$ and $\mu$, of the form:
\[ \psi(m,\mu)-\theta(m,\mu)=\frac{v}{(2n+1)^{2}}. \]

To demonstrate this, it is necessary to distinguish two cases: whether the limit of $\frac{m\omega+\mu\varpi i}{2n+1}$ is a quantity different from zero, or whether it is equal to zero.

$a)$ In the first case we shall have, on naming $a$ the limit in question:
\[ \begin{aligned}\varphi^{2}\left(\frac{m\omega+\mu\varpi i+k}{2n+1}\right)&=\Phi^{2}a+v,\\ \varphi^{2}\left(\frac{m\omega+\mu\varpi i+l}{2n+1}\right)&=\Phi^{2}a+v',\\ \varphi^{2}\left(\frac{\alpha}{2n+1}\right)&=\frac{\alpha^{2}}{(2n+1)^{2}}+\frac{v'}{(2n+1)^{2}},\end{aligned} \]
therefore:
\[ \begin{aligned}1-\frac{\varphi^{2}\left(\frac{\alpha}{2n+1}\right)}{\varphi^{2}\left(\frac{m\omega+\mu\varpi i+k}{2n+1}\right)}&=1-\frac{\alpha^{2}}{(2n+1)^{2}\varphi^{2}a}+\frac{v}{(2n+1)^{2}},\\ 1-\frac{\varphi^{2}\left(\frac{\alpha}{2n+1}\right)}{\varphi^{2}\left(\frac{m\omega+\mu\varpi i+l}{2n+1}\right)}&=1-\frac{\alpha^{2}}{(2n+1)^{2}.\varphi^{2}a}+\frac{v'}{(2n+1)^{2}}.\end{aligned} \]

\origpage{166}
We have likewise
\[ \begin{aligned}1-\frac{\alpha^{2}}{(m\omega+\mu\varpi i+k)^{2}}&=1-\frac{\alpha^{2}}{(2n+1)^{2}\left\{\frac{m\omega+\mu\varpi i+k}{(2n+1)}\right\}}=1-\frac{\alpha^{2}}{(2n+1)^{2}a^{2}}+\frac{v}{(2n+1)^{2}},\\ 1-\frac{\alpha^{2}}{(m\omega+\mu\varpi i+l)^{2}}&=1-\frac{\alpha^{2}}{(2n+1)^{2}a^{2}}+\frac{v'}{(2n+1)^{2}}.\end{aligned} \]
\ednote{In the first line the braced quotient in the middle denominator is printed without a square; a square is expected, as in the neighbouring terms.}

On substituting these values, the expression of $\psi(m,\mu)-\theta(m,\mu)$ will take the form:
\[ \psi(m,\mu)-\theta(m,\mu)=\log\left\{\frac{1-\frac{v}{(2n+1)^{2}}}{1-\frac{v'}{(2n+1)^{2}}}.\frac{1-\frac{v_{1}}{(2n+1)^{2}}}{1-\frac{v'_{1}}{(2n+1)^{2}}}\right\}, \]
the quantities $v$, $v'$, $v_{1}$, $v'_{1}$ all having zero for their limit.

We have therefore:
\[ \log\left(1-\frac{v}{(2n+1)^{2}}\right)=\frac{v}{(2n+1)^{2}}\ \text{etc.}, \]
and consequently:
\[ \psi(m,\mu)-\theta(m,\mu)=\frac{v}{(2n+1)^{2}}. \]

$b)$ If the limit of the quantity $\frac{m\omega+\mu\varpi i}{2n+1}$ is equal to zero, we shall have:
\[ \begin{aligned}\varphi^{2}\left(\frac{m\omega+\mu\varpi i+k}{2n+1}\right)&=\frac{(m\omega+\mu\varpi i+k)^{2}}{(2n+1)^{2}}+A.\frac{(m\omega+\mu\varpi i+k)^{4}}{(2n+1)^{4}}+....,\\ \varphi^{2}\left(\frac{\alpha}{2n+1}\right)&=\frac{\alpha^{2}}{(2n+1)^{2}}+A'.\frac{\alpha^{4}}{(2n+1)^{4}}+....,\end{aligned} \]
therefore:
\[ 1-\frac{\varphi^{2}\left(\frac{\alpha}{2n+1}\right)}{\varphi^{2}\left(\frac{m\omega+\mu\varpi i+k}{2n+1}\right)}=1-\frac{\alpha^{2}+A'.\frac{\alpha^{4}}{(2n+1)^{2}}+.....}{(m\omega+\mu\varpi i+k)^{2}+A.\frac{(m\omega+\mu\varpi i+k)^{4}}{(2n+1)^{2}}}. \]

If now $m\omega+\mu\varpi i$ does not go on increasing indefinitely with $n$, we shall have:
\[ 1-\frac{\varphi^{2}\left(\frac{\alpha}{2n+1}\right)}{\varphi^{2}\left(\frac{m\omega+m\varpi i+k}{2n+1}\right)}=1-\frac{\alpha^{2}}{(m\omega+\mu\varpi i+k)^{2}}+\frac{B}{(2n+1)^{2}}, \]
\ednote{Printed $m\varpi i$ in the argument of the denominator function instead of $\mu\varpi i$; an apparent misprint.}
likewise:
\[ 1-\frac{\varphi^{2}\left(\frac{\alpha}{2n+1}\right)}{\varphi^{2}\left(\frac{m\omega+\mu\varpi i+k}{2n+1}\right)}=1-\frac{\alpha^{2}}{(m\omega+\mu\varpi i+l)^{2}}+\frac{C}{(2n+1)^{2}}, \]
\ednote{The argument of the denominator function is printed with $+k$ (the glyph is blotted) where $+l$ is expected, as in the right-hand side.}
therefore in this case:

\origpage{167}
\[ \psi(m,\mu)-\theta(m,\mu)=\log\left\{\frac{1-\frac{B'}{(2n+1)^{2}}}{1-\frac{C'}{(2n+1)^{2}}}\right\}, \]
$B'$ and $C'$ having finite limits, or else:
\[ \psi(m,\mu)-\theta(m,\mu)=\frac{D}{(2n+1)^{2}}, \]
the limit of $D$ being likewise a finite quantity.

If on the contrary the quantity $m\omega+\mu\varpi i$ increases indefinitely with $n$, we have:
\[ \begin{aligned}\frac{1-\frac{\varphi^{2}\left(\frac{\alpha}{2n+1}\right)}{\varphi^{2}\left(\frac{m\omega+\mu\varpi i+k}{2n+1}\right)}}{1-\frac{\alpha^{2}}{(m\omega+\mu\varpi i+k)^{2}}}&=1+\frac{\alpha^{2}}{(2n+1)^{2}}.\left\{\frac{A-A'\frac{\alpha^{2}}{(m\omega+\mu\varpi i+k)^{2}}}{1+A\left(\frac{m\omega+\mu\varpi i+k}{2n+1}\right)^{2}}\right\}\\ &\times\frac{1}{1-\frac{\alpha^{2}}{(m\omega+\mu\varpi i+k)^{2}}};\end{aligned} \]
now the quantities $\frac{1}{m\omega+\mu\varpi i+k}$, $\frac{m\omega+\mu\varpi i+k}{2n+1}$ have zero for their limit; therefore the preceding quantity will be of the form:
\[ 1+\frac{\alpha^{2}}{(2n+1)^{2}}.A'', \]
$A''$ having a finite quantity for its limit. On changing $k$ into $l$, and designating the corresponding value of $A''$ by $A''_{1}$, the value of $\psi(m,\mu)-\theta(m,\mu)$ will become:
\[ \psi(m,\mu)-\theta(m,\mu)=\log\left\{\frac{1+\frac{\alpha^{2}}{(2n+1)^{2}}A''}{1+\frac{\alpha^{2}}{(2n+1)^{2}}A''_{1}}\right\}=\frac{\alpha^{2}(A''-A''_{1})}{(2n+1)^{2}}+\frac{v}{(2n+1)^{2}}. \]
Now the limit of $A''$ is the same as that of $A$; but it is clear that this latter limit is independent of $k$, $m$, $\mu$ (it is in fact equal to the coefficient of $\alpha^{4}$ in the development of $\Phi^{2}\alpha$). Therefore we shall have:
\[ A''=M+v, \]
and on changing $k$ into $l$:
\[ A''_{1}=M+v', \]
whence $A''-A''_{1}=v-v'=v$. Therefore $A''-A''_{1}$ has zero for its limit, and consequently we have:
\[ \psi(m,\mu)-\theta(m,\mu)=\frac{v}{(2n+1)^{2}}. \]

\origpage{168}

Therefore we have demonstrated that, on making
\[ \psi(m,\mu)-\theta(m,\mu)=\frac{A_{m,\mu}}{(2n+1)^{2}}, \tag{167.} \]
the limit of $A_{m,\mu}$ will be equal to zero every time that $m\omega+\mu\varpi i$ increases indefinitely with $n$, and that it will be equal to a finite quantity in the contrary case.

\subsection*{29.}

This being premised, let us consider the quantity
\[ \overset{n}{\underset{1}{\Sigma}}_{\mu}\psi(m,\mu). \]
On substituting the value of $\psi(m,\mu)$, there will result:
\[ \overset{n}{\underset{1}{\Sigma}}_{\mu}\psi(m,\mu)=\overset{n}{\underset{1}{\Sigma}}_{\mu}\theta(m,\mu)+\frac{1}{(2n+1)^{2}}.\overset{n}{\underset{1}{\Sigma}}_{\mu}(A_{m,\mu}). \tag{168.} \]
Let $\nu$ be the greatest whole number contained in $\sqrt{n}$; we may put:
\[ \begin{aligned}\overset{n}{\underset{1}{\Sigma}}_{\mu}A_{m,\mu}&=A_{m,1}+A_{m,2}+.....+A_{m,\nu}\\ &+A_{m,\nu+1}+A_{m,\nu+2}+.....+A_{m,n}.\end{aligned} \]

Now, from the nature of the quantities $A_{m,\mu}$, the sum contained in the first line will be equal to $\nu.A_{m}$, and the second equal to $A'_{m}(n-\nu-1)$, where $A_{m}$ is a finite quantity and $A'_{m}$ a quantity which has zero for its limit, therefore:
\[ \overset{n}{\underset{1}{\Sigma}}_{\mu}A_{m,\mu}=\nu.A_{m}+(n-\nu-1)A'_{m}=(2n+1).B_{m}, \]
where
\[ B_{m}=\frac{\nu}{2n+1}A_{m}+\frac{n-\nu-1}{2n+1}A'_{m}. \]
Therefore the quantity $B_{m}$ has zero for its limit, remarking that $\nu$ does not exceed $\sqrt{n}$.

Hereby the expression of $\overset{n}{\underset{1}{\Sigma}}_{\mu}\psi(m,\mu)$ is changed into:
\[ \overset{n}{\underset{1}{\Sigma}}_{\mu}\psi(m,\mu)=\overset{n}{\underset{1}{\Sigma}}_{\mu}\theta(m,\mu)+\frac{B_{m}}{2n+1}. \tag{169.} \]

To have the limit of $\overset{n}{\underset{1}{\Sigma}}_{\mu}\theta(m,\mu)$, I write
\[ \overset{n}{\underset{1}{\Sigma}}_{\mu}\theta(m,\mu)=\overset{\infty}{\underset{1}{\Sigma}}_{\mu}\theta(m,\mu)-\overset{\infty}{\underset{n+1}{\Sigma}}_{\mu}\theta(m,\mu)=\overset{\infty}{\underset{1}{\Sigma}}_{\mu}\theta(m,\mu)-\overset{\infty}{\underset{1}{\Sigma}}_{\mu}\theta(m,\mu+n). \]
Now the value of $\overset{\infty}{\underset{1}{\Sigma}}_{\mu}\theta(m,\mu+n)$ may be found as follows.

\origpage{169}

We have:
\[ \begin{aligned}\theta(m.\,\mu+n)&=\log\left\{\frac{1-\frac{\alpha^{2}}{[m\omega+(\mu+n)\varpi i+k]^{2}}}{1-\frac{\alpha^{2}}{[m\omega+(\mu+n)\varpi i+l]^{2}}}\right\}\\ &=\alpha^{2}\left\{\frac{1}{[m\omega+(\mu+n)\varpi i+l]^{2}}-\frac{1}{[m\omega+(\mu+n)\varpi i+k]^{2}}\right\}\\ &-\tfrac{1}{2}\alpha^{4}\left\{\frac{1}{[m\omega+(\mu+n)\varpi i+l]^{2}}-\frac{1}{[m\omega+(\mu+n)\varpi i+k]^{2}}\right\}\\ &+\ \text{etc.}\end{aligned} \]
\ednote{Printed $\theta(m.\,\mu+n)$ with a full stop instead of a comma between the arguments; an apparent misprint.}
From this we derive:
\[ \overset{\mu}{\underset{1}{\Sigma}}_{\mu}\theta(m,\mu+n)=\frac{\alpha^{2}}{n}.\Sigma\frac{1}{n}.\theta\left(\frac{\mu}{n}\right)-\frac{\alpha^{4}}{2n^{3}}.\Sigma\frac{1}{n}.\theta_{1}\left(\frac{\mu}{n}\right)+...., \]
where
\[ \begin{gathered}\theta\left(\frac{\mu}{n}\right)=\frac{1}{\left(\frac{m\omega+l}{n}+\varpi i+\frac{\mu}{n}\varpi i\right)^{2}}-\frac{1}{\left(\frac{m\omega+k}{n}+\varpi i+\frac{\mu}{n}\varpi i\right)^{2}},\\ \theta_{1}\left(\frac{\mu}{n}\right)=\frac{1}{\left(\frac{m\omega+l}{n}+\varpi i+\frac{\mu}{n}\varpi i\right)^{4}}-\frac{1}{\left(\frac{m\omega+k}{n}+\varpi i+\frac{\mu}{n}\varpi i\right)^{4}},\\ \text{etc.}\end{gathered} \]
Now it is known that the limit of
\[ \overset{\mu}{\underset{1}{\Sigma}}_{\mu}\frac{1}{n}\theta\left(\frac{\mu}{n}\right)\ \text{is equal to}\ \int_{0}^{x}\theta(x)\,\partial x, \]
therefore:
\[ \begin{gathered}\overset{\mu}{\underset{1}{\Sigma}}_{\mu}\frac{1}{n}\theta\left(\frac{\mu}{n}\right)=\int_{0}^{x}\theta(x)\,\partial x+v,\\ \overset{\mu}{\underset{1}{\Sigma}}_{\mu}\frac{1}{n}\theta_{1}\left(\frac{\mu}{n}\right)=\int_{0}^{x}\theta_{1}(x)\,\partial x+v_{1},\\ \text{etc},\end{gathered} \]
and consequently, on substituting:
\[ \overset{\mu}{\underset{1}{\Sigma}}_{\mu}\theta(m,\mu+n)=\frac{\alpha^{2}}{n}.\int_{0}^{x}\theta(x)-\frac{\alpha^{4}}{2n^{3}}.\int_{0}^{x}\theta_{1}(x)+....+\frac{v\alpha^{2}}{n}-\frac{v_{1}\alpha^{3}}{2n^{3}}+.... \]
\ednote{The integrals are printed without the differential $\partial x$, and the last fraction has $\alpha^{3}$ where $\alpha^{4}$ is expected; both apparent misprints, reproduced as printed.}
where
\[ \begin{gathered}\theta(x)=\frac{1}{\left(\frac{m\omega+l}{n}+\varpi i+x\varpi i\right)^{2}}-\frac{1}{\left(\frac{m\omega+k}{n}+\varpi i+x\varpi i\right)^{2}}\\ \text{etc.}\end{gathered} \]
therefore we shall have:

\origpage{170}
\[ \begin{aligned}\int_{0}^{x}\theta(x)\,\partial x&=\frac{1}{\varpi i}.\left\{\frac{1}{\left(\frac{m\omega+l}{n}+\varpi i+x\varpi i\right)}-\frac{1}{\left(\frac{m\omega+k}{n}+\varpi i+x\varpi i\right)}\right\}\\ &-\frac{1}{\varpi i}.\left\{\frac{1}{\left(\frac{m\omega+l}{n}+\varpi i\right)}-\frac{1}{\left(\frac{m\omega+k}{n}+\varpi i\right)}\right\}\\ &=\frac{1}{\varpi i}.\frac{k-l}{n}.\frac{1}{\left(\frac{m\omega+l}{n}+\varpi i+x\varpi i\right)\left(\frac{m\omega+k}{n}+\varpi i+x\varpi i\right)}\\ &-\frac{1}{\varpi i}.\frac{k-l}{n}.\frac{1}{\left(\frac{m\omega+l}{n}+\varpi i\right)\left(\frac{m\omega+k}{n}+\varpi i\right)}.\end{aligned} \]

The limit of this expression of $\displaystyle\int_{0}^{x}\theta(x)\,\partial x$ is zero for any value whatever of $x$. In the same way it will be found that the limit of $\displaystyle\int_{0}^{x}\theta_{1}(x)\,\partial x$ is zero, therefore:
\[ \begin{aligned}\overset{\mu}{\underset{1}{\Sigma}}_{\mu}\theta(m,\mu+n)&=\frac{\alpha^{2}}{n}.v-\frac{\alpha^{4}}{2n^{3}}.v'+\frac{\alpha^{6}}{3n^{5}}.v''-....\\ &=\frac{\alpha^{2}}{2n+1}.\left\{v-\frac{\alpha^{2}}{2n^{2}}v'+\frac{\alpha^{4}}{3n^{4}}v''-....\right\}\frac{2n+1}{n}=\frac{v}{2n+1},\end{aligned} \]
therefore also, on making $\mu=\infty$:
\[ \overset{\infty}{\underset{1}{\Sigma}}_{\mu}\theta(m,\mu+n)=\frac{v}{2n+1}, \]
whence:
\[ \overset{n}{\underset{1}{\Sigma}}_{\mu}\psi(m,\mu+n)=\overset{\infty}{\underset{1}{\Sigma}}_{\mu}\theta(m,\mu)-\frac{v}{2n+1},\ \text{and} \]
\[ \overset{n}{\underset{1}{\Sigma}}_{\mu}\psi(m,\mu)=\overset{\infty}{\underset{1}{\Sigma}}_{\mu}\theta(m,\mu)+\frac{v_{m}}{2n+1}, \tag{170.} \]
$v_{m}$ having zero for its limit. From this we draw
\[ \overset{n}{\underset{1}{\Sigma}}_{m}\overset{n}{\underset{1}{\Sigma}}_{\mu}\psi(m,\mu)=\overset{n}{\underset{1}{\Sigma}}_{m}\left\{\overset{\infty}{\underset{1}{\Sigma}}_{\mu}\theta(m,\mu)\right\}+\overset{n}{\underset{1}{\Sigma}}\frac{v_{m}}{2n+1}. \]
On taking the limit of the two members and remarking that
\[ \overset{n}{\underset{1}{\Sigma}}\frac{v_{m}}{2n+1}=\frac{v_{1}+v_{2}+.....v_{m}}{2n+1}=v, \]
there will be:
\[ \text{lim.}\overset{n}{\underset{1}{\Sigma}}_{m}\overset{n}{\underset{1}{\Sigma}}_{\mu}\psi(m,\mu)=\overset{\infty}{\underset{1}{\Sigma}}_{m}\left\{\overset{\infty}{\underset{1}{\Sigma}}_{\mu}\theta(m,\mu)\right\}. \tag{171.} \]
On restoring the values of $\psi(m,\mu)$ and $\theta(m,\mu)$, and passing from logarithms to numbers, we draw from it:

\origpage{171}
\[ \text{lim.}\overset{n}{\underset{1}{\Pi}}_{m}\overset{n}{\underset{1}{\Pi}}_{\mu}\left\{\frac{1-\frac{\varphi^{2}\left(\frac{\alpha}{2n+1}\right)}{\varphi^{2}\left(\frac{m\omega+\mu\varpi i+k}{2n+1}\right)}}{1-\frac{\varphi^{2}\left(\frac{\alpha}{2n+1}\right)}{\varphi^{2}\left(\frac{m\omega+\mu\varpi i+l}{2n+1}\right)}}\right\}=\overset{\infty}{\underset{1}{\Pi}}_{m}\left\{\overset{\infty}{\underset{1}{\Pi}}_{\mu}\left(\frac{1-\frac{\alpha^{2}}{(m\omega+\mu\varpi i+k)^{2}}}{1-\frac{\alpha^{2}}{(m\omega+\mu\varpi i+l)^{2}}}\right)\right\}. \tag{172.} \]

By an analysis quite similar to the preceding, but simpler, we shall find in the same way:
\[ \text{lim.}\overset{n}{\underset{1}{\Pi}}_{m}\left\{\frac{1-\frac{\varphi^{2}\left(\frac{\alpha}{2n+1}\right)}{\varphi^{2}\left(\frac{m\omega+k}{2n+1}\right)}}{1-\frac{\varphi^{2}\left(\frac{\alpha}{2n+1}\right)}{\varphi^{2}\left(\frac{m\omega+l}{2n+1}\right)}}\right\}=\overset{\infty}{\underset{1}{\Pi}}_{m}\left\{\frac{1-\frac{\alpha^{2}}{(m\omega+k)^{2}}}{1-\frac{\alpha^{2}}{(m\omega+l)^{2}}}\right\}, \tag{173.} \]
\[ \text{lim.}\overset{n}{\underset{1}{\Pi}}_{\mu}\left\{\frac{1-\frac{\varphi^{2}\left(\frac{\alpha}{2n+1}\right)}{\varphi^{2}\left(\frac{\mu\varpi i+k}{2n+1}\right)}}{1-\frac{\varphi^{2}\left(\frac{\alpha}{2n+1}\right)}{\varphi^{2}\left(\frac{\mu\varpi i+k}{2n+1}\right)}}\right\}=\overset{\infty}{\underset{1}{\Pi}}_{\mu}\left\{\frac{1-\frac{\alpha^{2}}{(\mu\varpi i+k)^{2}}}{1-\frac{\alpha^{2}}{(\mu\varpi i+l)^{2}}}\right\}. \tag{174.} \]
\ednote{In the left side of the last equation the lower fraction is printed with $k$ in the argument of the denominator function, as the upper one; presumably $l$ was meant, as in the right side. An apparent misprint.}

\subsection*{30.}

Now nothing is easier than to find the values of $\Phi\alpha$, $f\alpha$, $F\alpha$.

Let us consider first the first formula (130.). We have:
\[ \begin{aligned}e^{2}c^{2}\Phi^{2}\left(\frac{m\omega+\mu\varpi i}{2n+1}\right)&=-\frac{1}{\varphi^{2}\left(\frac{m\omega+\mu\varpi i}{2n+1}-\frac{\omega}{2}-\frac{\varpi}{2}i\right)}\\ &=-\frac{1}{\varphi^{2}\left(\frac{(n-m+\tfrac{1}{2})\omega+(n-\mu+\tfrac{1}{2})\varpi i}{2n+1}\right)},\end{aligned} \]
therefore:

\origpage{172}
\[ \begin{aligned}\overset{n}{\underset{1}{\Pi}}_{m}\overset{\mu}{\underset{1}{\Pi}}\left\{\frac{1-\frac{\varphi^{2}\beta}{\varphi^{2}\left(\frac{m\omega+\mu\varpi i}{2n+1}\right)}}{1+e^{2}c^{2}.\varphi^{2}\left(\frac{m\omega+\mu\varpi i}{2n+1}\right)\varphi^{2}\beta}\right\}&=\frac{\overset{n}{\underset{1}{\Pi}}_{m}\overset{\mu}{\underset{1}{\Pi}}.\left\{1-\frac{\varphi^{2}\beta}{\varphi^{2}\left(\frac{m\omega+\mu\varpi i}{2n+1}\right)}\right\}}{\overset{n}{\underset{1}{\Pi}}_{m}\overset{\mu}{\underset{1}{\Pi}}.\left\{1-\frac{\varphi^{2}\beta}{\varphi^{2}\left(\frac{(n-m+\tfrac{1}{2})\omega+(n-\mu+\tfrac{1}{2})\varpi i}{2n+1}\right)}\right\}}\\ &=\overset{n}{\underset{1}{\Pi}}_{m}\overset{\mu}{\underset{1}{\Pi}}_{\mu}\left\{\frac{1-\frac{\varphi^{2}\beta}{\varphi^{2}\left(\frac{m\omega+\mu\varpi i}{2n+1}\right)}}{1-\frac{\varphi^{2}\beta}{\varphi^{2}\left(\frac{m\omega+\mu\varpi i+\frac{\omega}{2}+\frac{\varpi}{2}i}{2n+1}\right)}}\right\}.\end{aligned} \]

This being premised, if we make $\beta=\frac{\alpha}{2n+1}$ and suppose $n$ infinite, there will result, on making use of the formulas (172.), (173.), (174.), and remarking that $\alpha$ is the limit of $(2n+1)\varphi\left(\frac{\alpha}{2n+1}\right)$ and equal to $\alpha$:
\[ \begin{aligned}\Phi\alpha&=\alpha\overset{\infty}{\underset{1}{\Pi}}_{m}\left\{1-\frac{\alpha^{2}}{(m\omega)^{2}}\right\}.\overset{\infty}{\underset{1}{\Pi}}\left\{1+\frac{\alpha^{2}}{(\mu\varpi)^{2}}\right\}\\ &\times\overset{\infty}{\underset{1}{\Pi}}_{m}\left\{\overset{\infty}{\underset{1}{\Pi}}_{\mu}\left(\frac{1-\frac{\alpha^{2}}{(m\omega+\mu\varpi i)^{2}}}{1-\frac{\alpha^{2}}{[(m-\tfrac{1}{2})\omega+(\mu-\tfrac{1}{2})\varpi i]^{2}}}\right).\overset{\infty}{\underset{1}{\Pi}}_{\mu}\left(\frac{1-\frac{\alpha^{2}}{(m\omega-\mu\varpi i)^{2}}}{1-\frac{\alpha^{2}}{[(m-\tfrac{1}{2})\omega-(\mu-\tfrac{1}{2})\varpi i]^{2}}}\right)\right\}.\end{aligned} \tag{175.} \]

The two formulas (130.) will give in the same manner, on making $\beta=\frac{\alpha}{2n+1}$, and remarking that $f(0)=1$, $F(0)=1$:
\[ \begin{aligned}f\alpha&=\overset{\infty}{\underset{1}{\Pi}}_{m}\left\{\overset{\infty}{\underset{1}{\Pi}}_{\mu}\left(\frac{1-\frac{\alpha^{2}}{[(m-\tfrac{1}{2})\omega+\mu\varpi i]^{2}}}{1-\frac{\alpha^{2}}{[(m-\tfrac{1}{2})\omega+(\mu-\tfrac{1}{2})\varpi i]^{2}}}\right)\left(\frac{1-\frac{\alpha^{2}}{[(m-\tfrac{1}{2})\omega-\mu\varpi i]^{2}}}{1-\frac{\alpha^{2}}{[(m-\tfrac{1}{2})\omega-(\mu-\tfrac{1}{2})\varpi i]^{2}}}\right)\right\}\\ &\times\overset{\infty}{\underset{0}{\Pi}}_{\mu}\left\{1-\frac{\alpha^{2}}{(m+\tfrac{1}{2})^{2}\omega^{2}}\right\},\end{aligned} \tag{176.} \]
\[ \begin{aligned}F\alpha&=\overset{\infty}{\underset{0}{\Pi}}_{\mu}\left\{1+\frac{\alpha^{2}}{(\mu+\tfrac{1}{2})\varpi^{2}}\right\}.\overset{\infty}{\underset{1}{\Pi}}_{m}\left\{\overset{\infty}{\underset{1}{\Pi}}_{\mu}\left(\frac{1-\frac{\alpha^{2}}{[m\omega+(\mu-\tfrac{1}{2})\varpi i]^{2}}}{1-\frac{\alpha^{2}}{[(m-\tfrac{1}{2})\omega+(\mu-\tfrac{1}{2})\varpi i]^{2}}}\right)\left(\frac{1-\frac{\alpha^{2}}{[m\omega+(\mu-\tfrac{1}{2})\varpi i]^{2}}}{1-\frac{\alpha^{2}}{[(m-\tfrac{1}{2})\omega+(\mu-\tfrac{1}{2})\varpi i]^{2}}}\right)\right\}.\end{aligned} \tag{177.} \]
\ednote{The two factors of the last product are printed identically, where the earlier formulas have opposite signs; an apparent misprint, reproduced as printed. In the preceding formulas the factor $(m+\tfrac{1}{2})^{2}\omega^{2}$ of (176.) carries its square, and the product index there is printed $\mu$ although $m$ occurs; the factor $(\mu+\tfrac{1}{2})\varpi^{2}$ of (177.) is printed without a square on the bracket. All reproduced as printed.}

We can also give a real form to the preceding expressions as follows:

\origpage{173}
\[
\begin{aligned}
\Phi\alpha &= \alpha.\overset{\infty}{\underset{1}{\Pi}}_{\mu}\left(1+\frac{\alpha^{2}}{\mu^{2}\varpi^{2}}\right).\overset{\infty}{\underset{1}{\Pi}}_{m}\left(1-\frac{\alpha^{2}}{m^{2}\omega^{2}}\right)\\
&\quad\times\overset{\infty}{\underset{1}{\Pi}}_{m}\overset{\infty}{\underset{1}{\Pi}}_{\mu}\frac{1+\frac{(\alpha+m\omega)^{2}}{\mu^{2}\varpi^{2}}}{1+\frac{[\alpha+(m-\tfrac{1}{2})\omega]^{2}}{\mu^{2}\varpi^{2}}}.\frac{1+\frac{(\alpha-m\omega)^{2}}{\mu^{2}\varpi^{2}}}{1+\frac{[\alpha-(m-\tfrac{1}{2})\omega]^{2}}{\mu^{2}\varpi^{2}}}.\left(\frac{1+\frac{(m-\tfrac{1}{2})^{2}\omega^{2}}{\mu^{2}\varpi^{2}}}{1+\frac{m^{2}\omega^{2}}{\mu^{2}\varpi^{2}}}\right)^{2},
\end{aligned}
\tag{178.}
\]
\[
\begin{aligned}
f\alpha &= \overset{\infty}{\underset{1}{\Pi}}_{m}\left(1-\frac{\alpha^{2}}{(m-\tfrac{1}{2})^{2}\omega^{2}}\right)\\
&\quad\times\overset{\infty}{\underset{1}{\Pi}}_{m}\overset{\infty}{\underset{1}{\Pi}}_{\mu}\frac{1+\frac{[\alpha+(m-\tfrac{1}{2})\omega]^{2}}{\mu^{2}\varpi^{2}}}{1+\frac{[\alpha+(m-\tfrac{1}{2})\omega]^{2}}{(\mu-\tfrac{1}{2})^{2}\varpi^{2}}}.\frac{1+\frac{[\alpha-(m-\tfrac{1}{2})\omega]^{2}}{\mu^{2}\varpi^{2}}}{1+\frac{[\alpha-(m-\tfrac{1}{2})\omega]^{2}}{(\mu-\tfrac{1}{2})^{2}\varpi^{2}}}.\left(\frac{1+\frac{(m-\tfrac{1}{2})^{2}\omega^{2}}{(\mu-\tfrac{1}{2})^{2}\varpi^{2}}}{1+\frac{(m-\tfrac{1}{2})^{2}\omega^{2}}{\mu^{2}\varpi^{2}}}\right)^{2},
\end{aligned}
\tag{179.}
\]
\[
\begin{aligned}
F\alpha &= \overset{\infty}{\underset{0}{\Pi}}_{\mu}\left(1+\frac{\alpha^{2}}{(\mu+\tfrac{1}{2})^{2}\varpi^{2}}\right)\\
&\quad\times\overset{\infty}{\underset{1}{\Pi}}_{m}\overset{\infty}{\underset{1}{\Pi}}_{\mu}\frac{1+\frac{(\alpha+m\omega)^{2}}{(\mu-\tfrac{1}{2})^{2}\varpi^{2}}}{1+\frac{[\alpha+(m-\tfrac{1}{2})\omega]^{2}}{(\mu-\tfrac{1}{2})^{2}\varpi^{2}}}.\frac{1+\frac{(\alpha-m\omega)^{2}}{(\mu-\tfrac{1}{2})^{2}\varpi^{2}}}{1+\frac{[\alpha-(m-\tfrac{1}{2})\omega]^{2}}{(\mu-\tfrac{1}{2})^{2}\varpi^{2}}}.\left(\frac{1+\frac{(m-\tfrac{1}{2})^{2}\omega^{2}}{(\mu-\tfrac{1}{2})^{2}\varpi^{2}}}{1+\frac{m^{2}\omega^{2}}{(\mu-\tfrac{1}{2})^{2}\varpi^{2}}}\right)^{2}.
\end{aligned}
\tag{180.}
\]

These transformations are easily effected by means of the formula:
\[
\begin{aligned}
\left(1-\frac{\alpha^{2}}{(a+bi)^{2}}\right)\left(1-\frac{\alpha^{2}}{(a-bi)^{2}}\right) &= \left(1+\frac{\alpha}{a+bi}\right)\left(1+\frac{\alpha}{a-bi}\right)\left(1-\frac{\alpha}{a+bi}\right)\left(1-\frac{\alpha}{a-bi}\right)\\
&= \frac{(\alpha+a)^{2}+b^{2}}{a^{2}+b^{2}}.\frac{(\alpha-a)^{2}+b^{2}}{a^{2}+b^{2}} = \left(1+\frac{(\alpha+a)^{2}}{b^{2}}\right)\left(1+\frac{(\alpha-a)^{2}}{b^{2}}\right).\frac{1}{\left(1+\frac{a^{2}}{b^{2}}\right)^{2}}.
\end{aligned}
\]

\subsection*{31.}

In what precedes we have arrived at two kinds of expressions of the functions $\Phi\alpha$, $f\alpha$, $F\alpha$: the one kind gives these functions decomposed into partial fractions, the totality of which forms infinite double series, the other gives these same functions decomposed into an infinite number of factors, each of which is in its turn composed of an infinity of factors.

Now the preceding formulas can be much simplified by means of the exponential and circular functions. This is what we shall see in what follows:

Let us consider first the equations (178.), (179.), (180.). By virtue of the known formulas, we have:

\origpage{174}
\[
\frac{\sin y}{y} = \overset{\infty}{\underset{1}{\Pi}}_{\mu}\left(1-\frac{y^{2}}{\mu^{2}\pi^{2}}\right);\ \cos y = \overset{\infty}{\underset{1}{\Pi}}_{\mu}\left(1-\frac{y^{2}}{(\mu-\tfrac{1}{2})^{2}\pi^{2}}\right);
\]
therefore:
\[
\overset{\infty}{\underset{1}{\Pi}}_{\mu}\frac{\left(1-\frac{z^{2}}{\mu^{2}\pi^{2}}\right)}{\left(1-\frac{y^{2}}{(\mu-\tfrac{1}{2})\pi^{2}}\right)} = \frac{\sin z}{z\cos y};\ \overset{\infty}{\underset{1}{\Pi}}_{\mu}\left(\frac{1-\frac{z^{2}}{(\mu-\tfrac{1}{2})\pi^{2}}}{1-\frac{y^{2}}{(\mu-\tfrac{1}{2})\pi^{2}}}\right) = \frac{\cos z}{\cos y}.
\]
\ednote{In both products the bracket $(\mu-\tfrac{1}{2})$ under $\pi^{2}$ is printed without the square (the denominator of the first product, and the numerator and denominator of the second); an apparent misprint for $(\mu-\tfrac{1}{2})^{2}\pi^{2}$, reproduced as printed.}

By virtue of this formula it is clear that the expressions of $\Phi\alpha$, $f\alpha$, $F\alpha$ may be put under the form:
\[
\begin{aligned}
\Phi\alpha &= \frac{\varpi}{\pi}.\frac{\sin\left(\alpha\frac{\pi}{\varpi}i\right)}{i}.\overset{\infty}{\underset{1}{\Pi}}_{m}\left(1-\frac{\alpha^{2}}{m^{2}\omega^{2}}\right)\\
&\quad\times\overset{\infty}{\underset{1}{\Pi}}_{m}\left(\frac{\sin(\alpha+m\omega)\frac{\pi}{\varpi}i.\sin(\alpha-m\omega)\frac{\pi}{\varpi}i.\cos^{2}(m-\tfrac{1}{2})\omega\frac{\pi}{\varpi}i}{\cos[\alpha+(m-\tfrac{1}{2})\omega]\frac{\pi}{\varpi}i.\cos[\alpha-(m-\tfrac{1}{2})\omega]\frac{\pi}{\varpi}i.\sin^{2}m\omega\frac{\pi}{\varpi}i}.\frac{\left(m\omega\frac{\pi}{\varpi}i\right)^{2}}{(\alpha+m\omega)(\alpha-m\omega).\frac{\pi^{2}}{\omega^{2}}i^{2}}\right),
\end{aligned}
\]
\ednote{In the denominator of the second fraction the factor $\frac{\pi^{2}}{\omega^{2}}$ is printed with $\omega$, where $\varpi$ is expected (compare the factor $\frac{\pi}{\varpi}$ elsewhere in the line); an apparent misprint, reproduced as printed.}
\[
\begin{aligned}
f\alpha &= \overset{\alpha}{\underset{1}{\Pi}}_{m}\left(1-\frac{\alpha^{2}}{(m-\tfrac{1}{2})^{2}\omega^{2}}\right)\\
&\quad\times\overset{\infty}{\underset{1}{\Pi}}_{m}\left(\text{tang}[\alpha+(m-\tfrac{1}{2})\omega]\frac{\pi}{\varpi}i.\text{tang}[\alpha-(m-\tfrac{1}{2})\omega]\frac{\pi}{\varpi}i.\cot^{2}(m-\tfrac{1}{2})\omega\frac{\pi}{\varpi}i.\frac{(m-\tfrac{1}{2})^{2}\omega^{2}}{\alpha^{2}-(m-\tfrac{1}{2})^{2}\omega^{2}}\right),
\end{aligned}
\]
\ednote{In the first line of the expression for $f\alpha$ the upper limit of the product is printed as $\alpha$, where $\infty$ is expected; an apparent misprint, reproduced as printed.}
\[
F\alpha = \cos\left(\alpha\frac{\pi}{\varpi}i\right).\overset{\infty}{\underset{1}{\Pi}}_{m}\left(\frac{\cos(\alpha+m\omega)\frac{\pi}{\varpi}i.\cos(\alpha-m\omega)\frac{\pi}{\varpi}i.\cos^{2}(m-\tfrac{1}{2})\omega\frac{\pi}{\varpi}i}{\cos[\alpha+(m-\tfrac{1}{2})\omega]\frac{\pi}{\varpi}i.\cos[\alpha-(m-\tfrac{1}{2})\omega]\frac{\pi}{\varpi}i.\cos^{2}m\omega.\frac{\pi}{\varpi}i}\right).
\]

Real expressions will be found by substituting for the circular functions their expressions in exponential functions.

We have:
\[
\begin{gathered}
\sin(a-b).\sin(a+b) = \sin^{2}a-\sin^{2}b,\\
\cos(a+b).\cos(a-b) = \cos^{2}a-\sin^{2}b,
\end{gathered}
\]
therefore:
\[
\frac{\sin(\alpha+m\omega)\frac{\pi}{\varpi}i.\sin(\alpha-m\omega)\frac{\pi}{\varpi}i}{\sin^{2}m\omega\frac{\pi}{\varpi}i} = -\left\{1-\frac{\sin^{2}\alpha\frac{\pi}{\varpi}i}{\sin^{2}m\omega\frac{\pi}{\varpi}i}\right\},
\]
\[
\frac{\cos[\alpha+(m-\tfrac{1}{2})\omega]\frac{\pi}{\varpi}i.\cos[\alpha-(m-\tfrac{1}{2})\omega]\frac{\pi}{\varpi}i}{\cos^{2}(m-\tfrac{1}{2})\omega\frac{\pi}{\varpi}i} = 1-\frac{\sin^{2}\alpha\frac{\pi}{\varpi}i}{\cos^{2}(m-\tfrac{1}{2})\omega\frac{\pi}{\varpi}i},
\]

\origpage{175}
\[
\text{tang}[\alpha+(m-\tfrac{1}{2})\omega]\frac{\pi}{\varpi}i.\text{tang}[\alpha-(m-\tfrac{1}{2})\omega]\frac{\pi}{\varpi}i.\cot^{2}(m-\tfrac{1}{2})\omega\frac{\pi}{\varpi}i = -\frac{1-\frac{\sin^{2}\alpha\frac{\pi}{\varpi}i}{\sin^{2}(m-\tfrac{1}{2})\omega\frac{\pi}{\varpi}i}}{1-\frac{\sin^{2}\alpha\frac{\pi}{\varpi}i}{\cos^{2}(m-\tfrac{1}{2})\omega\frac{\pi}{\varpi}i}}.
\]

According to this, and remarking that
\[
\frac{m^{2}\omega^{2}}{\alpha^{2}-m^{2}\omega^{2}} = -\frac{1}{1-\frac{\alpha^{2}}{m^{2}\omega^{2}}}\ \text{and}
\]
\[
\frac{(m-\tfrac{1}{2})^{2}\omega^{2}}{\alpha^{2}-(m-\tfrac{1}{2})^{2}\omega^{2}} = -\frac{1}{1-\frac{\alpha^{2}}{(m-\tfrac{1}{2})^{2}\omega^{2}}},
\]
it is clear that we shall have:
\[
\Phi\alpha = \frac{\pi}{\varpi}.\frac{\sin\left(\frac{\alpha}{\varpi}\pi i\right)}{i}.\overset{\infty}{\underset{1}{\Pi}}_{m}.\frac{1-\frac{\sin^{2}\alpha\frac{\pi}{\varpi}i}{\sin^{2}m\omega\frac{\pi}{\varpi}i}}{1-\frac{\sin^{2}\alpha\frac{\pi}{\varpi}i}{\cos^{2}(m-\tfrac{1}{2})\omega\frac{\pi}{\varpi}i}},
\tag{181.}
\]
\ednote{The factor before the sine is printed as the fraction with $\pi$ over $\varpi$, whereas the preceding expression has $\varpi$ over $\pi$; an apparent misprint, reproduced as printed.}
\[
f\alpha = \overset{\infty}{\underset{0}{\Pi}}_{m}.\frac{1-\frac{\sin^{2}\alpha\frac{\pi}{\varpi}i}{\sin^{2}(m+\tfrac{1}{2})\omega\frac{\pi}{\varpi}i}}{1-\frac{\sin^{2}\alpha\frac{\pi}{\varpi}i}{\cos^{2}(m+\tfrac{1}{2})\omega\frac{\pi}{\varpi}i}},
\tag{182.}
\]
\[
F\alpha = \cos\left(\frac{\alpha}{\varpi}\pi i\right).\overset{\infty}{\underset{1}{\Pi}}_{m}.\frac{1-\frac{\sin^{2}\frac{\alpha}{\varpi}\pi i}{\cos^{2}m\frac{\omega}{\varpi}\pi i}}{1-\frac{\sin^{2}\frac{\alpha}{\varpi}\pi i}{\cos^{2}(m-\tfrac{1}{2})\frac{\omega}{\varpi}\pi i}}.
\tag{183.}
\]

\origpage{176}

On substituting for the cosines and sines of imaginary arcs their values in exponential quantities, these formulas will become:
\[
\Phi\alpha = \tfrac{1}{2}.\frac{\varpi}{\pi}\left(h^{\frac{\alpha}{\varpi}\pi}-h^{-\frac{\alpha}{\varpi}\pi}\right).\overset{\infty}{\underset{1}{\Pi}}_{m}.\frac{1-\left(\frac{h^{\frac{\alpha}{\varpi}\pi}-h^{-\frac{\alpha}{\varpi}\pi}}{h^{m\frac{\omega}{\varpi}\pi}-h^{-m\frac{\omega}{\varpi}\pi}}\right)^{2}}{1+\left(\frac{h^{\frac{\alpha}{\varpi}\pi}-h^{-\frac{\alpha}{\varpi}\pi}}{h^{(m-\frac{1}{2})\frac{\omega}{\varpi}\pi}+h^{-(m-\frac{1}{2})\frac{\omega}{\varpi}\pi}}\right)^{2}},
\tag{184.}
\]
\[
f\alpha = \overset{\infty}{\underset{0}{\Pi}}_{m}.\frac{1-\left(\frac{h^{\frac{\alpha}{\varpi}\pi}-h^{-\frac{\alpha}{\varpi}\pi}}{h^{(m+\frac{1}{2})\frac{\omega}{\varpi}\pi}-h^{-(m+\frac{1}{2})\frac{\omega}{\varpi}\pi}}\right)^{2}}{1+\left(\frac{h^{\frac{\alpha}{\varpi}\pi}-h^{-\frac{\alpha}{\varpi}\pi}}{h^{(m+\frac{1}{2})\frac{\omega}{\varpi}\pi}+h^{-(m+\frac{1}{2})\frac{\omega}{\varpi}\pi}}\right)^{2}},
\tag{185.}
\]
\[
F\alpha = \tfrac{1}{2}\left(e^{\frac{\alpha}{\varpi}\pi}+e^{-\frac{\alpha}{\varpi}\pi}\right).\overset{\infty}{\underset{1}{\Pi}}_{m}\frac{1+\left(\frac{h^{\frac{\alpha}{\varpi}\pi}-h^{-\frac{\alpha}{\varpi}\pi}}{h^{m\frac{\omega}{\varpi}\pi}+h^{-m\frac{\omega}{\varpi}\pi}}\right)^{2}}{1+\left(\frac{h^{\frac{\alpha}{\varpi}\pi}-h^{-\frac{\alpha}{\varpi}\pi}}{h^{(m-\frac{1}{2})\frac{\omega}{\varpi}\pi}+h^{-(m-\frac{1}{2})\frac{\omega}{\varpi}\pi}}\right)^{2}},
\tag{186.}
\]
where $h$ is the number 2,1718218.....
\ednote{The value of $h$ is printed as 2,1718218; the base of the natural logarithms is 2,7182818; an apparent misprint, reproduced as printed.} \ednote{In the first factor of the last formula the base is printed as $e$ in the sum, where the other exponentials of the page use $h$; an apparent misprint, reproduced as printed.}

We can further transform this formula in the following manner:

If we replace $\alpha$ by $\alpha i$, we shall have the values of $\Phi(\alpha i)$, $f(\alpha i)$, $F(\alpha i)$. On now changing $c$ into $e$ and $e$ into $c$, the quantities:
\[
\omega;\ \varpi;\ \Phi(\alpha i);\ f(\alpha i);\ F(\alpha i)
\]
will be changed respectively into:
\[
\varpi;\ \omega;\ i\Phi\alpha;\ F\alpha;\ f\alpha,
\]
therefore the preceding formulas will give:
\[
\Phi\alpha = \frac{\omega}{\pi}.\sin\frac{\alpha\pi}{\omega}.\overset{\infty}{\underset{1}{\Pi}}_{m}.\frac{1+\frac{4\sin^{2}\left(\frac{\alpha\pi}{\omega}\right)}{\left(h^{\frac{m\varpi\pi}{\omega}}-h^{-\frac{m\varpi\pi}{\omega}}\right)^{2}}}{1-\frac{4\sin^{2}\left(\frac{\alpha\pi}{\omega}\right)}{\left(h^{\frac{(2m-1)\varpi\pi}{2\omega}}+h^{-\frac{(2m-1)\varpi\pi}{2\omega}}\right)^{2}}},
\tag{187.}
\]

\origpage{177}
\[
F\alpha = \overset{\infty}{\underset{0}{\Pi}}_{m}\frac{1+\frac{4\sin^{2}\left(\frac{\alpha\pi}{\omega}\right)}{\left(h^{\frac{(2m+1)\varpi\pi}{2\omega}}-h^{-\frac{(2m+1)\varpi\pi}{2\omega}}\right)^{2}}}{1-\frac{4\sin^{2}\left(\frac{\alpha\pi}{\omega}\right)}{\left(h^{\frac{(2m+1)\varpi\pi}{2\omega}}+h^{-\frac{(2m+1)\varpi\pi}{2\omega}}\right)^{2}}},
\tag{188.}
\]
\[
f\alpha = \cos\left(\frac{\alpha\pi}{\omega}\right).\overset{\infty}{\underset{1}{\Pi}}_{m}.\frac{1-\frac{4\sin^{2}\left(\frac{\alpha\pi}{\omega}\right)}{\left(h^{\frac{m\varpi\pi}{\omega}}+h^{-\frac{m\varpi\pi}{\omega}}\right)^{2}}}{1-\frac{4\sin^{2}\left(\frac{\alpha\pi}{\omega}\right)}{\left(h^{\frac{(2m+1)\varpi\pi}{2\omega}}+h^{-\frac{(2m+1)\varpi\pi}{2\omega}}\right)^{2}}},
\tag{189.}
\]

\subsection*{32.}

Let us now consider the formulas (160.), (161.), (162.).

We have:
\[
\overset{\infty}{\underset{0}{\Sigma}}_{\mu}(-1)^{\mu}.\frac{(2\mu+1)\pi}{y^{2}+(\mu+\tfrac{1}{2})^{2}\pi^{2}} = \frac{2}{h^{y}+h^{-y}},
\]
therefore, on making
\[
y = [\alpha\pm(m+\tfrac{1}{2})\omega]\frac{\pi}{\varpi}:
\]
\[
\overset{\infty}{\underset{0}{\Sigma}}_{\mu}(-1)^{\mu}.\frac{(2\mu+1)\varpi}{[\alpha\pm(m+\tfrac{1}{2})\omega]^{2}+(\mu+\tfrac{1}{2})^{2}\varpi^{2}} = \frac{2\pi}{\varpi}.\frac{1}{h^{[\alpha\pm(m+\frac{1}{2})\omega]\frac{\pi}{\varpi}}+h^{-[\alpha\pm(m+\frac{1}{2})\omega]\frac{\pi}{\varpi}}}.
\]
By virtue of this formula it is easy to see that the expressions (153.), (162.) of $\Phi\alpha$ and $F\alpha$ will become:
\[
\Phi\alpha = \frac{2}{ec}.\frac{\pi}{\varpi}.\overset{\infty}{\underset{0}{\Sigma}}_{m}(-1)^{m}.\left\{\frac{1}{h^{[\alpha-(m+\frac{1}{2})\omega]\frac{\pi}{\varpi}}+h^{-[\alpha-(m-\frac{1}{2})\omega]\frac{\pi}{\varpi}}}-\frac{1}{h^{[\alpha+(m+\frac{1}{2})\omega]\frac{\pi}{\varpi}}+h^{-[\alpha+(m+\frac{1}{2})\omega]\frac{\pi}{\varpi}}}\right\},
\tag{190.}
\]
\ednote{In the first denominator the second exponent is printed with $m-\tfrac{1}{2}$ where the first has $m+\tfrac{1}{2}$; an apparent misprint, reproduced as printed.}
\[
F\alpha = \frac{2}{c}.\frac{\pi}{\varpi}.\overset{\infty}{\underset{0}{\Sigma}}_{m}\left\{\frac{1}{h^{[\alpha-(m+\frac{1}{2})\omega]\frac{\pi}{\varpi}}+h^{-[\alpha-(m+\frac{1}{2})\omega]\frac{\pi}{\varpi}}}+\frac{1}{h^{[\alpha+(m+\frac{1}{2})\omega]\frac{\pi}{\varpi}}+h^{-[\alpha+(m+\frac{1}{2})\omega]\frac{\pi}{\varpi}}}\right\}.
\tag{191.}
\]

The preceding expressions of $\Phi\alpha$, $F\alpha$ may be put yet under many other forms: I shall recall the most remarkable. First, on uniting the terms of the second member, we shall find:

\origpage{178}
\[
\Phi\alpha = \frac{2}{ec}.\frac{\pi}{\varpi}.\overset{\infty}{\underset{0}{\Sigma}}_{m}(-1)^{m}.\left\{\frac{\left(h^{\frac{\alpha\pi}{\varpi}}-h^{-\frac{\alpha\pi}{\varpi}}\right)\left(h^{(m+\frac{1}{2})\frac{\omega\pi}{\varpi}}-h^{-(m+\frac{1}{2})\frac{\omega\pi}{\varpi}}\right)}{h^{\frac{2\alpha\pi}{\varpi}}+h^{-\frac{2\alpha\pi}{\varpi}}+h^{(2m+1)\frac{\omega\pi}{\varpi}}+h^{-(2m+1)\frac{\omega\pi}{\varpi}}}\right\},
\tag{192.}
\]
\[
F\alpha = \frac{2}{c}.\frac{\pi}{\varpi}.\overset{\infty}{\underset{0}{\Sigma}}_{m}.\left\{\frac{\left(h^{\frac{\alpha\pi}{\varpi}}+h^{-\frac{\alpha\pi}{\varpi}}\right)\left(h^{(m+\frac{1}{2})\frac{\omega\pi}{\varpi}}+h^{-(m+\frac{1}{2})\frac{\omega\pi}{\varpi}}\right)}{h^{\frac{2\alpha\pi}{\varpi}}+h^{-\frac{2\alpha\pi}{\varpi}}+h^{(2m+1)\frac{\omega\pi}{\varpi}}+h^{-(2m+1)\frac{\omega\pi}{\varpi}}}\right\}.
\tag{193.}
\]
If, for brevity, we suppose:
\[
h^{\frac{\alpha\pi}{\varpi}} = \varepsilon\ \text{and}\ h^{\frac{\omega\pi}{\varpi}} = r^{2},
\tag{194.}
\]
these formulas, on developing the second member, will become:
\[
\Phi\alpha = \frac{2}{ec}.\frac{\pi}{\varpi}\left(\varepsilon-\frac{1}{\varepsilon}\right).\left\{\frac{r-\frac{1}{r}}{r^{2}+\varepsilon^{2}+\frac{1}{\varepsilon^{2}}+\frac{1}{r^{2}}}-\frac{r^{3}-\frac{1}{r^{3}}}{r^{6}+\varepsilon^{2}+\frac{1}{\varepsilon^{2}}+\frac{1}{r^{6}}}+\frac{r^{5}-\frac{1}{r^{5}}}{r^{10}+\varepsilon^{2}+\frac{1}{\varepsilon^{2}}+\frac{1}{r^{10}}}-\ldots\right\},
\tag{195.}
\]
\[
F\alpha = \frac{2}{c}.\frac{\pi}{\varpi}.\left(\varepsilon+\frac{1}{\varepsilon}\right).\left\{\frac{r+\frac{1}{r}}{r^{2}+\varepsilon^{2}+\frac{1}{\varepsilon^{2}}+\frac{1}{r^{2}}}+\frac{r^{3}+\frac{1}{r^{3}}}{r^{6}+\varepsilon^{2}+\frac{1}{\varepsilon^{2}}+\frac{1}{r^{6}}}+\frac{r^{5}+\frac{1}{r^{5}}}{r^{10}+\varepsilon^{2}+\frac{1}{\varepsilon^{2}}+\frac{1}{r^{10}}}-\ldots\right\}.
\tag{196.}
\]

On putting $\alpha i$ in place of $\alpha$ in the formulas (192.), (193.), then changing $c$ into $e$ and $e$ into $c$, and remarking that the quantities
\[
\omega,\ \varpi,\ \Phi(\alpha i),\ F(\alpha i),\ h^{\alpha i\frac{\pi}{\varpi}}-h^{-\alpha i\frac{\pi}{\varpi}},\ h^{\alpha i\frac{\pi}{\varpi}}+h^{\alpha i\frac{\pi}{\varpi}},
\]
\ednote{In the last quantity the sign of the second exponent is printed without a minus sign; an apparent misprint for $h^{-\alpha i\frac{\pi}{\varpi}}$, reproduced as printed.}
will be changed respectively into:
\[
\varpi,\ \omega,\ i\Phi(\alpha),\ f\alpha,\ 2i.\sin\alpha\frac{\pi}{\omega},\ 2\cos\alpha\frac{\pi}{\omega},
\]
there will result:
\[
\Phi\alpha = \frac{4}{ec}.\frac{\pi}{\omega}.\overset{\infty}{\underset{0}{\Sigma}}_{m}(-1)^{m}.\left\{\frac{\sin\frac{\alpha\pi}{\omega}.\left(h^{(m+\frac{1}{2})\frac{\varpi\pi}{\omega}}-h^{-(m+\frac{1}{2})\frac{\varpi\pi}{\omega}}\right)}{h^{(2m+1)\frac{\varpi\pi}{\omega}}+2\cos2\alpha\frac{\pi}{\omega}+h^{-(2m+1)\frac{\varpi\pi}{\omega}}}\right\},
\tag{197.}
\]
\[
f\alpha = \frac{4}{e}.\frac{\pi}{\omega}.\overset{\infty}{\underset{0}{\Sigma}}_{m}.\left\{\frac{\cos\frac{\alpha\pi}{\omega}.\left(h^{(m+\frac{1}{2})\frac{\varpi\pi}{\omega}}+h^{-(m+\frac{1}{2})\frac{\varpi\pi}{\omega}}\right)}{h^{(2m+1)\frac{\varpi\pi}{\omega}}+2\cos2\alpha\frac{\pi}{\omega}+h^{-(2m+1)\frac{\varpi\pi}{\omega}}}\right\}.
\tag{198.}
\]
On making for brevity
\[
h^{\frac{\varpi\pi}{2\omega}} = \varrho,
\tag{199.}
\]

\origpage{179}
and developing, we shall obtain:
\[
\varphi\left(\alpha\frac{\omega}{2}\right) = \frac{4}{ec}.\frac{\pi}{\omega}.\sin\left(\alpha\frac{\pi}{2}\right).\left\{\frac{\varrho-\frac{1}{\varrho}}{\varrho^{2}+2\cos(\alpha\pi)+\frac{1}{\varrho^{2}}}-\frac{\varrho^{3}-\frac{1}{\varrho^{3}}}{\varrho^{6}+2\cos(\alpha\pi)+\frac{1}{\varrho^{6}}}+\frac{\varrho^{5}-\frac{1}{\varrho^{5}}}{\varrho^{10}+2\cos(\alpha\pi)+\frac{1}{\varrho^{10}}}-....\right\},
\tag{200.}
\]
\[
f\left(\alpha\frac{\omega}{2}\right) = \frac{4}{e}.\frac{\pi}{\omega}.\cos\left(\alpha\frac{\pi}{2}\right).\left\{\frac{\varrho+\frac{1}{\varrho}}{\varrho^{2}+2\cos(\alpha\pi)+\frac{1}{\varrho^{2}}}+\frac{\varrho^{3}+\frac{1}{\varrho^{3}}}{\varrho^{6}+2\cos(\alpha\pi)+\frac{1}{\varrho^{6}}}+\frac{\varrho^{5}+\frac{1}{\varrho^{5}}}{\varrho^{10}+2\cos(\alpha\pi)+\frac{1}{\varrho^{10}}}+....\right\}.
\tag{201.}
\]

On substituting in the formulas (190.), (191.) for $h^{\alpha\frac{\pi}{\varpi}}$ and $h^{\frac{\omega\pi}{2\varpi}}$ their values $\varepsilon$ and $r$, there will result:
\[
\Phi\alpha = \frac{2}{ec}.\frac{\pi}{\varpi}.\overset{\infty}{\underset{0}{\Sigma}}_{m}(-m).\left\{\frac{1}{\varepsilon.r^{-(2m+1)}+\varepsilon^{-1}.r^{2m+1}}-\frac{1}{\varepsilon.r^{2m+1}+\varepsilon^{-1}.r^{-(2m+1)}}\right\},
\tag{202.}
\]
\ednote{The factor is printed as $(-m)$ where $(-1)^{m}$ is expected; an apparent misprint, reproduced as printed.}
\[
F\alpha = \frac{2}{c}.\frac{\pi}{\varpi}.\overset{\infty}{\underset{0}{\Sigma}}_{m}\left\{\frac{1}{\varepsilon.r^{-2m-1}+\varepsilon^{-1}.r^{2m+1}}+\frac{1}{\varepsilon.r^{2m+1}+\varepsilon^{-1}.r^{-2m-1}}\right\}.
\tag{203.}
\]

Supposing now $\alpha<\frac{\omega}{2}$, we shall have:
\[
\frac{1}{\varepsilon.r^{-2m-1}+\varepsilon^{-1}.r^{2m+1}} = \frac{\varepsilon^{-1}.r^{-2m-1}}{1+\varepsilon^{2}.r^{-4m-2}} = \varepsilon.r^{-2m-1}-\varepsilon^{3}.r^{-6m-3}+\varepsilon^{5}.r^{-10m-5}-....,
\]
\ednote{In the first fraction the numerator is printed $\varepsilon^{-1}.r^{-2m-1}$ where $\varepsilon.r^{-2m-1}$ is expected (the expansion beside it begins with $\varepsilon.r^{-2m-1}$); an apparent misprint, reproduced as printed.}
\[
\frac{1}{\varepsilon.r^{2m+1}+\varepsilon^{-1}.r^{-2m-1}} = \frac{\varepsilon^{-1}.r^{-2m-1}}{1+\varepsilon^{-2}.r^{-4m-2}} = \varepsilon^{-1}.r^{-2m-1}-\varepsilon^{-3}.r^{-6m-3}+\varepsilon^{-5}.r^{-10m-5}-....,
\]
therefore:
\[
\begin{aligned}
\Phi\alpha &= \frac{2}{ec}.\frac{\pi}{\varpi}.\left\{\overset{\infty}{\underset{0}{\Sigma}}_{m}(-1)^{m}\left\{(\varepsilon-\varepsilon^{-1}).r^{-2m-1}-(\varepsilon^{3}-\varepsilon^{-3}).r^{-6m-3}+(\varepsilon^{5}-\varepsilon^{-5})\,r^{-10m-5}-....\right\}\right\}\\
&= \frac{2}{ec}.\frac{\pi}{\varpi}.\left\{(\varepsilon-\varepsilon^{-1}).\overset{\infty}{\underset{0}{\Sigma}}_{m}(-1)^{m}.r^{-2m-1}-(\varepsilon^{3}-\varepsilon^{-3}).\overset{\infty}{\underset{0}{\Sigma}}_{m}(-1)^{m}.r^{-6m-3}+....\right\}.
\end{aligned}
\]
Now:
\[
\overset{\infty}{\underset{0}{\Sigma}}_{m}(-1)^{m}.r^{-2m-1} = r^{-1}-r^{-3}+r^{-5}-.... = \frac{r^{-1}}{1+r^{-2}} = \frac{r}{r^{2}+1},
\]
\[
\overset{\infty}{\underset{0}{\Sigma}}_{m}(-1)^{m}.r^{3m-3} = r^{-3}-r^{-6}+r^{-9}-.... = \frac{r^{-3}}{1+r^{-6}} = \frac{r^{3}}{r^{6}+1}
\]
\ednote{The exponent in the summand is printed $3m-3$ and the series exponents as $-3$, $-6$, $-9$; the sum is evidently meant to be that of $(-1)^{m}.r^{-6m-3}$, with series $r^{-3}-r^{-9}+r^{-15}-\ldots$; apparent misprints, reproduced as printed.}
etc., therefore:
\[
\Phi\alpha = \frac{2}{ec}.\frac{\pi}{\varpi}.\left\{\frac{\varepsilon-\varepsilon^{-1}}{r+r^{-1}}-\frac{\varepsilon^{3}-\varepsilon^{-3}}{r^{3}+r^{-3}}+\frac{\varepsilon^{5}-\varepsilon^{-5}}{r^{5}+r^{-5}}-.....\right\}.
\tag{204.}
\]
In the same manner we shall find:
\[
F\alpha = \frac{2}{c}.\frac{\pi}{\varpi}.\left\{\frac{\varepsilon+\varepsilon^{-1}}{r-r^{-1}}-\frac{\varepsilon^{3}+\varepsilon^{-3}}{r^{3}-r^{-3}}+\frac{\varepsilon^{5}+\varepsilon^{-5}}{r^{5}-r^{-5}}-.....\right\}.
\tag{205.}
\]

\origpage{180}
On putting $\alpha\frac{\varpi}{2}i$ in place of $\alpha$, and then changing $e$ into $c$ and $c$ into $e$,
\[
\varepsilon;\ \varpi;\ \Phi\left(\alpha\frac{\varpi}{2}i\right);\ F\left(\alpha\frac{\varpi}{2}i\right);\ r;\ \varepsilon^{m}+\varepsilon^{-m};\ \varepsilon^{m}-\varepsilon^{-m}
\]
are changed into:
\[
\varpi;\ \omega;\ i\varphi\left(\alpha\frac{\omega}{2}\right);\ f\left(\alpha\frac{\omega}{2}\right);\ \varrho;\ 2\cos\left(m\alpha\frac{\pi}{2}\right);\ 2i\sin\left(m\alpha\frac{\pi}{2}\right);
\]
therefore:
\[
\varphi\left(\alpha\frac{\omega}{2}\right) = \frac{4}{ec}.\frac{\pi}{\omega}.\left\{\frac{\sin\left(\alpha\frac{\pi}{2}\right)}{\varrho+\frac{1}{\varrho}}-\frac{\sin\left(3\alpha\frac{\pi}{2}\right)}{\varrho^{3}+\frac{1}{\varrho^{3}}}+\frac{\sin\left(5\alpha\frac{\pi}{2}\right)}{\varrho^{5}+\frac{1}{\varrho^{5}}}-....\right\},
\tag{206.}
\]
\[
f\left(\alpha\frac{\omega}{2}\right) = \frac{4}{e}.\frac{\pi}{\omega}.\left\{\frac{\cos\left(\alpha\frac{\pi}{2}\right)}{\varrho-\frac{1}{\varrho}}-\frac{\cos\left(3\alpha\frac{\pi}{2}\right)}{\varrho^{3}-\frac{1}{\varrho^{3}}}+\frac{\cos\left(5\alpha\frac{\pi}{2}\right)}{\varrho^{5}-\frac{1}{\varrho^{5}}}-....\right\}.
\tag{207.}
\]

These last four formulas offer very simple expressions of the functions $\Phi\alpha$, $f\alpha$, $F\alpha$. By differentiations and integrations a multitude of others, more or less remarkable, may be deduced from them.

\subsection*{33.}

In the case $e=c$, the preceding formulas take a simpler form, on account of the relation $\omega=\varpi$, which holds in this case. Let, for greater simplicity, $e=c=1$. We have:
\[
r = h^{\frac{\omega\pi}{2\varpi}} = h^{\frac{\pi}{2}},\ \varrho = h^{\frac{\varpi\pi}{2\omega}} = h^{\frac{\pi}{2}},
\]
therefore on substituting, and making in (204.), (205.) $\alpha=\alpha\frac{\varpi}{2}$, there results:
\[
\varphi\left(\alpha\frac{\omega}{2}\right) = 2\frac{\pi}{\omega}\left\{\frac{h^{\frac{\alpha\pi}{2}}-h^{-\frac{\alpha\pi}{2}}}{h^{\frac{\pi}{2}}+h^{-\frac{\pi}{2}}}-\frac{h^{\frac{3\alpha\pi}{2}}-h^{-\frac{3\alpha\pi}{2}}}{h^{\frac{3\pi}{2}}+h^{-\frac{3\pi}{2}}}+\frac{h^{\frac{5\alpha\pi}{2}}-h^{-\frac{5\alpha\pi}{2}}}{h^{\frac{5\pi}{2}}+h^{-\frac{5\pi}{2}}}-.....\right\},
\]
\[
F\left(\alpha\frac{\omega}{2}\right) = 2\frac{\pi}{\omega}\left\{\frac{h^{\frac{\alpha\pi}{2}}+h^{-\frac{\alpha\pi}{2}}}{h^{\frac{\pi}{2}}-h^{-\frac{\pi}{2}}}-\frac{h^{\frac{3\alpha\pi}{2}}+h^{-\frac{3\alpha\pi}{2}}}{h^{\frac{3\pi}{2}}-h^{-\frac{3\pi}{2}}}+\frac{h^{\frac{5\alpha\pi}{2}}+h^{-\frac{5\alpha\pi}{2}}}{h^{\frac{5\pi}{2}}-h^{-\frac{5\pi}{2}}}-.....\right\},
\]
\[
\varphi\left(\alpha\frac{\omega}{2}\right) = \frac{4\pi}{\omega}\left\{\sin\left(\alpha\frac{\pi}{2}\right).\frac{h^{\frac{\pi}{2}}}{1+h^{\pi}}-\sin\left(3\alpha\frac{\pi}{2}\right).\frac{h^{\frac{3\pi}{2}}}{1+h^{3\pi}}+\sin\left(5\alpha\frac{\pi}{2}\right).\frac{h^{\frac{5\pi}{2}}}{1+h^{5\pi}}-....\right\},
\]
\[
f\left(\alpha\frac{\omega}{2}\right) = \frac{4\pi}{\omega}\left\{\cos\left(\alpha\frac{\pi}{2}\right).\frac{h^{\frac{\pi}{2}}}{h^{\pi}-1}-\cos\left(3\alpha\frac{\pi}{2}\right).\frac{h^{\frac{3\pi}{2}}}{h^{3\pi}-1}+\sin\left(5\alpha\frac{\pi}{2}\right).\frac{h^{\frac{5\pi}{2}}}{h^{5\pi}-1}-....\right\}.
\]
\ednote{In the last display the third term has a sine, where the pattern of the first two terms requires a cosine; an apparent misprint, reproduced as printed.}

\origpage{181}
The functions $\varphi$, $f$, $F$ are determined by the equations
\[
\begin{gathered}
\alpha.\frac{\omega}{2} = \int_{0}^{x}\frac{\partial x}{\sqrt{(1-x^{4})}};\ \frac{\omega}{2} = \int_{0}^{1}\frac{\partial x}{\sqrt{(1-x^{4})}};\\
x = \varphi\left(\alpha\frac{\omega}{2}\right) = \sqrt{(1-x^{2})} = f\left(\alpha\frac{\omega}{2}\right) = \sqrt{(1+x^{2})} = F\left(\alpha\frac{\omega}{2}\right).
\end{gathered}
\]

If in the last two formulas we make $\alpha=0$, and remark that the value of $\dfrac{\varphi\left(\alpha\frac{\omega}{2}\right)}{\sin\left(\alpha\frac{\pi}{2}\right)}$ is then $=\frac{\omega}{\pi}$, and that of $\dfrac{\sin\left(m\alpha\frac{\pi}{2}\right)}{\sin\left(\alpha\frac{\pi}{2}\right)}=m$, we shall find:
\[
\frac{\omega}{2} = \frac{\pi}{2}.\left\{\frac{h^{\frac{\pi}{2}}}{h^{\pi}-1}-\frac{h^{\frac{3\pi}{2}}}{h^{3\pi}-1}+\frac{h^{\frac{5\pi}{2}}}{h^{5\pi}-1}-\ldots\right\} = \int_{0}^{1}\frac{\partial x}{\sqrt{(1-x^{4})}},
\]
\[
\frac{\omega^{2}}{4} = \pi^{2}.\left\{\frac{h^{\frac{\pi}{2}}}{h^{\pi}+1}-3.\frac{h^{\frac{3\pi}{2}}}{h^{3\pi}+1}+5.\frac{h^{\frac{5\pi}{2}}}{h^{5\pi}+1}-\ldots\right\} = \left(\int_{0}^{1}\frac{\partial x}{\sqrt{(1-x^{4})}}\right)^{2}.
\]

(The continuation in the next number.)

\end{document}
